% Options for packages loaded elsewhere
\PassOptionsToPackage{unicode}{hyperref}
\PassOptionsToPackage{hyphens}{url}
\documentclass[
  english,
]{article}
\usepackage{xcolor}
\usepackage[margin=1in]{geometry}
\usepackage{amsmath,amssymb}
\setcounter{secnumdepth}{-\maxdimen} % remove section numbering
\usepackage{iftex}
\ifPDFTeX
  \usepackage[T1]{fontenc}
  \usepackage[utf8]{inputenc}
  \usepackage{textcomp} % provide euro and other symbols
\else % if luatex or xetex
  \usepackage{unicode-math} % this also loads fontspec
  \defaultfontfeatures{Scale=MatchLowercase}
  \defaultfontfeatures[\rmfamily]{Ligatures=TeX,Scale=1}
\fi
\usepackage{lmodern}
\ifPDFTeX\else
  % xetex/luatex font selection
  \setmainfont[]{STIX Two Text}
  \setmathfont[]{STIX Two Math}
\fi
% Use upquote if available, for straight quotes in verbatim environments
\IfFileExists{upquote.sty}{\usepackage{upquote}}{}
\IfFileExists{microtype.sty}{% use microtype if available
  \usepackage[]{microtype}
  \UseMicrotypeSet[protrusion]{basicmath} % disable protrusion for tt fonts
}{}
\makeatletter
\@ifundefined{KOMAClassName}{% if non-KOMA class
  \IfFileExists{parskip.sty}{%
    \usepackage{parskip}
  }{% else
    \setlength{\parindent}{0pt}
    \setlength{\parskip}{6pt plus 2pt minus 1pt}}
}{% if KOMA class
  \KOMAoptions{parskip=half}}
\makeatother
\usepackage{longtable,booktabs,array}
\newcounter{none} % for unnumbered tables
\usepackage{multirow}
\usepackage{calc} % for calculating minipage widths
% Correct order of tables after \paragraph or \subparagraph
\usepackage{etoolbox}
\makeatletter
\patchcmd\longtable{\par}{\if@noskipsec\mbox{}\fi\par}{}{}
\makeatother
% Allow footnotes in longtable head/foot
\IfFileExists{footnotehyper.sty}{\usepackage{footnotehyper}}{\usepackage{footnote}}
\makesavenoteenv{longtable}
\ifLuaTeX
\usepackage[bidi=basic,shorthands=off]{babel}
\else
\usepackage[bidi=default,shorthands=off]{babel}
\fi
\ifPDFTeX
\else
\babelfont{rm}[]{STIX Two Text}
\fi
\ifLuaTeX
  \usepackage{selnolig} % disable illegal ligatures
\fi
\setlength{\emergencystretch}{3em} % prevent overfull lines
\providecommand{\tightlist}{%
  \setlength{\itemsep}{0pt}\setlength{\parskip}{0pt}}
\usepackage{amsmath}
\usepackage{amssymb}
\usepackage{newunicodechar}
\usepackage{pdflscape}

\providecommand{\citeproctext}[1]{#1}
\providecommand{\citeproc}[2]{#2}
\providecommand{\tightlist}{%
  \setlength{\itemsep}{0pt}\setlength{\parskip}{0pt}}

\newunicodechar{✓}{\ensuremath{\checkmark}}

\setlength{\emergencystretch}{3em}
\sloppy
\usepackage{bookmark}
\IfFileExists{xurl.sty}{\usepackage{xurl}}{} % add URL line breaks if available
\urlstyle{same}
% fallback for those not using the hyperref driver hyperxmp:
\makeatletter
\@ifundefined{xmpquote}{\newcommand{\xmpquote}[1]{#1}}{}
\makeatother
\hypersetup{
  pdflang={en},
  hidelinks,
  pdfcreator={LaTeX via pandoc}}

\author{}
\date{}

\begin{document}

\begin{otherlanguage}{english}

\protect\phantomsection\label{abstracts}
\protect\phantomsection\label{ab005}
\begin{otherlanguage}{english}

\subsection{Highlights}\label{highlights}

\protect\phantomsection\label{as005}
\protect\phantomsection\label{sp0005}
\begin{itemize}
\item
  \protect\phantomsection\label{p0005}
  We introduce a bike repositioning problem with broken bikes and on-site repairers.
\item
  \protect\phantomsection\label{p0010}
  We propose a time-indexed mixed integer linear program for the problem.
\item
  \protect\phantomsection\label{p0015}
  We develop an effective hybrid algorithm to solve the problem.
\item
  \protect\phantomsection\label{p0020}
  The algorithm relies a new concept of station budgets to determine loading and repair
  instructions.
\item
  \protect\phantomsection\label{p0025}
  Hybrid genetic search is used to determine truck and repair routes.
\end{itemize}

\end{otherlanguage}

\protect\phantomsection\label{ab010}
\subsection{Abstract}\label{abstract}

\protect\phantomsection\label{as010}
\protect\phantomsection\label{sp0010}
On-site repairs by repairers can handle broken bikes in a bike-sharing system and simultaneously
satisfy usable bike demand without requiring vehicle repositioning for off-site repairs. However,
this option has not been jointly considered with vehicle-based repositioning for both usable and
broken bikes in the literature to improve repositioning performance. This paper presents a
mixed-integer linear programming model for a static bike repositioning problem that combines
vehicle-based bike delivery/collection with labor-based on-site repairs, aiming to minimize the
total cost of user dissatisfaction and carbon emissions within a time budget. A hybrid algorithm,
consisting of hybrid genetic search with adaptive diversity control and a station budget constrained
heuristic, is proposed to solve the problem. This method introduces a station budget concept to
limit the time spent at each station by considering the benefit-cost ratio function of each station
along the route. The proposed algorithm is evaluated on networks of various sizes and configurations
of the number of trucks and repairers. The results demonstrate that the algorithm can obtain optimal
solutions on small instances and can outperform a commercial solver on larger networks in achieving
superior solutions within a fraction of the computational time. Moreover, we investigated the
cost-effectiveness of introducing repairers, revealing that long repair times and a low percentage
of broken bikes diminish the effectiveness of introducing them.

\protect\phantomsection\label{kg005}
\subsection{Keywords}\label{keywords}

\protect\phantomsection\label{k0005}
Bike repositioning problem

;

\protect\phantomsection\label{k0010}
Broken bike repositioning

;

\protect\phantomsection\label{k0015}
On-site repair routing

;

\protect\phantomsection\label{k0020}
Time-indexed formulation

;

\protect\phantomsection\label{k0025}
Station budget constrained heuristic

\protect\phantomsection\label{body}
\protect\phantomsection\label{s0005}
\subsection{1. Introduction}\label{st025}

\protect\phantomsection\label{p0035}
As an indispensable part of urban public transportation, Bike Sharing Systems (BSSs) provide a
sustainable mode of transportation for citizens that can effectively reduce traffic congestion and
air pollution. However, the imbalance between demand and supply seriously impedes the smooth
operation of a BSS. For instance, after morning peak hours in a station-based BSS, some stations
have a shortage of shared bicycles or even no shared bicycles available, while some have a surplus
inventory of shared bicycles that leaves nearly no vacant docks at those stations. Consequently,
users are often left with the problem of no shared bikes available for renting or no vacant docks
available for returning the bikes. To ensure that each station has enough shared bikes or unoccupied
docks to meet the customers' rental or return demand, the bikes must be redistributed regularly. The
problem of designing effective schemes for shared bicycle redistribution is known as the Bike
Repositioning Problem (BRP).

\protect\phantomsection\label{p0040}
Currently, there are two forms of bicycle repositioning: static and dynamic. For the static
repositioning, the operation is carried out throughout the night, assuming that the demand for
shared bicycles during that period is negligible. The purpose of this operation is to rearrange the
bicycles to prepare for the next day. In contrast, the dynamic repositioning runs throughout the day
and adjusts to the time-varying demand, and the planned scheme is altered frequently (Forma et al.,
2015). These two forms of repositioning have gained the attention of researchers, and the
corresponding static and dynamic BRPs have been well-studied. However, in addition to redistributing
shared bikes between stations to meet user demand, the problem arising from broken bikes in BSSs is
also non-negligible. For example, in Barcelona, the operator reported that about 12~\% of shared
bikes were severely damaged, and about 55~\% were lightly damaged (Castro Fernández, 2011). In New
York, 2~\% of the total number of bikes in the BSS need major repairs every day (Kaspi et al.,
2016). From users' perspective, if a broken bike is found when renting a bike, time for seeking
another usable bike is required. In a worse case, when all the bikes located at the station are
broken, the users are required to go to another station. Meanwhile, a user may also encounter all
docks being occupied when returning a bike, whereas some parked bikes might be damaged, taking up
the space for returning usable bikes. As Kaspi et al. (2016) pointed out, even if the percentage of
damaged bikes in a BSS is small, it can still seriously affect user satisfaction with the system.
The decrease in user satisfaction and confidence in the system may lead to the abandonment of usage.
Therefore, it is necessary to consider handling broken bicycles when solving BRPs.

\protect\phantomsection\label{p0045}
There have been a few studies on BRPs with broken bikes, some of which focused on broken bikes only
and formulated models for the optimal routes for broken bike collection (Zhang et al., 2018, Lu et
al., 2019, 2022). The remainder considered both broken and usable bikes and designed the routes of
repositioning vehicles and loading instructions for both usable and broken bikes at each station
(Alvarez-Valdes et al., 2016, Chang et al., 2018, Wang and Szeto, 2018, Wang and Szeto, 2021, Du et
al., 2020, Zhang et al., 2020). All these studies on BRPs with broken bikes assumed that broken
bikes were only repaired at depots (representing bike shops or repair centers) and required to be
collected by repositioning trucks first. Nevertheless, in practical operations, broken bicycle
repairs are not restricted to shops or repair centers. They are also conducted in the field by
dispatching repairers to the respective locations. According to Citi Bike's monthly operating report
for June 2024 (Citi Bike, 2024), 21,637 unique bikes were checked or repaired by repairers in depots
or the field during January. On-site repairs by repairers can handle the broken bikes and satisfy
the usable bike demand simultaneously without vehicle repositioning. Therefore, it is essential to
determine how to send repairers to the field (i.e., the routes of repairers for on-site repairs) and
how many broken bikes should be handled by each repairer at each visited station in addition to the
optimal strategy combining usable bike repositioning and broken bike collection. To the best of our
knowledge, no studies have made such integration at the time of this writing. Moreover, as discussed
above, user dissatisfaction is greatly influenced by the presence of broken bikes. However, most
existing studies focus on minimizing the absolute deviations of the inventories from a desired level
or maximizing the satisfied demand for bikes (Xu et al., 2019, Chang et al., 2021, Wang and Szeto,
2021, Cai et al., 2022), which does not necessarily lead to the minimized user dissatisfaction.
Therefore, it is vital to take user dissatisfaction into account when designing the optimization
objectives.

\protect\phantomsection\label{p0050}
As stated above, regular usable bike rebalancing in a BSS is required. However, frequent usable bike
rebalancing can significantly increase the fuel consumption of rebalancing vehicles (Luo et al.,
2020), which contradicts the original aim of running a BSS that offers an eco-friendly mode of
transportation. For example, in London, due to the intensive repositioning work, the fossil-fueled
vehicle usage for rebalancing has outweighed the environmental benefit of launching the BSS as the
replacement for car trips (Fishman et al., 2014). Shui and Szeto (2018) pointed out that
repositioning plans with long routes or heavy vehicle loads may produce more air pollutants. The
problem is more severe if the collection of broken bikes is included, given that the collected
broken bikes have to be carried all the way and cannot be unloaded until the vehicles arrive at
depots. Hence, it is meaningful to take the environmental factors into account when designing the
optimal rebalancing scheme.

\protect\phantomsection\label{p0055}
In this paper, we propose a static BRP with broken bikes and on-site repairers. The problem is
finding an effective scheme for repositioning trucks and repairers. This scheme encompasses the
routes and the loading/unloading instructions of broken and usable bikes for the trucks, the routes
for the repairers, and the number of broken bikes to be repaired by each repairer at each station.
The objective is to minimize the sum of the penalty costs of the total user dissatisfaction at all
stations and the cost of the total CO\textsubscript{2} emissions of the trucks. We consider the cost
of total CO\textsubscript{2} emissions as the environmental objective, as CO\textsubscript{2} is one
of the most severe threats to the environment through the greenhouse effect. Using a time-index
formulation approach, a Mixed Integer Linear Programming (MILP) model is proposed for the problem.
To solve the problem efficiently, we propose a hybrid algorithm, abbreviated as HGSADC-SBC,
encapsulating Hybrid Genetic Search with Adaptive Dynamic Control (HGSADC) and a newly proposed
Station Budget Constrained (SBC) heuristic. The algorithmic performance is demonstrated under the
scenarios with different network sizes and numbers of repairers and trucks. Moreover, the
cost-effectiveness of introducing a repairer under different repairing times and broken bike
inventory levels was studied on a 15-station instance.

\protect\phantomsection\label{p0060}
The contributions of this paper lie in the following:

\begin{itemize}
\item
  \protect\phantomsection\label{p0065}
  to introduce a new and practical BRP with broken bikes and on-site repairers;
\item
  \protect\phantomsection\label{p0070}
  to propose a time-indexed MILP model for the problem and solve it;
\item
  \protect\phantomsection\label{p0075}
  to present new findings and operational insights:
\item
  \protect\phantomsection\label{p0080}
  We find that long repair times and a low proportion of broken bikes lead to the diminishment of
  the cost-effectiveness of introducing on-site repairers.
\item
  \protect\phantomsection\label{p0085}
  We demonstrate that the cost-effectiveness of introducing on-site repairers changes dynamically
  with the proportion of broken bikes, indicating the need for adaptable repairer allocation.
\end{itemize}

\protect\phantomsection\label{p0090}
The remainder of the paper is organized as follows. Section 2 reviews the literature on BRPs and
points out the research gaps. Section 3 presents a mathematical model for the studied problem.
Section 4 introduces the proposed solution algorithm. Section 5 describes the numerical studies and
discusses the results. Finally, Section 6 concludes this paper.

\protect\phantomsection\label{s0010}
\subsection{2. Literature review}\label{st030}

\protect\phantomsection\label{p0095}
In this section, we first briefly introduce existing research on BRPs. Then, we discuss how current
research deals with broken bikes in BRPs. Finally, we summarize the above works and point out the
research gaps.

\protect\phantomsection\label{s0015}
\subsubsection{2.1. Categories, objectives, and algorithms in existing studies on BRPs}\label{st035}

\protect\phantomsection\label{p0100}
Existing approaches to dealing with bike demand--supply imbalance can be divided into two
categories: user- and operator-based repositioning. User-based methods intend to motivate users to
participate in the bike repositioning process by encouraging them to pick up or return bikes at
designated stations in exchange for monetary incentives (Singla et al., 2015, Pan et al., 2019,
Cheng et al., 2021, Wang and Wang, 2021). While this mechanism can help align supply and demand
during system usage, leading to efficient resource allocation, its ability to increase social and
economic surplus may be limited. For instance, if the pricing scheme is set too high, the system
might achieve perfect balance but fail to generate societal benefits due to reduced participation.
Moreover, the effectiveness of this method depends on the willingness of users to participate, which
can become highly uncontrollable. As a result, the operator-based rebalancing approaches are adopted
more extensively in current BSSs (Singla et al., 2015). In operator-based bike repositioning,
vehicles depart from the depot with a load of bikes, pick up and drop off a certain number of bikes
at surplus and deficit stations, respectively, and then return to the depot where they started (Lu
et al., 2020). Specifically, as mentioned in the introduction, operator-based bike repositioning can
be further divided into static repositioning (Raviv et al., 2013, Forma et al., 2015, Liu et al.,
2018, Huang et al., 2020, Wang and Szeto, 2021) and dynamic repositioning (Caggiani and Ottomanelli,
2013, O'Mahony and Shmoys, 2015, Caggiani et al., 2018, Brinkmann et al., 2019, Chiariotti et al.,
2020, Hu et al., 2021, Chen et al., 2024, Guo et al., 2024). These two forms of operator-based
repositioning have been extensively modeled in the BRP literature.

\protect\phantomsection\label{p0105}
The objectives of existing BRP research were mainly developed according to~the critical concerns of
the bike-sharing operators, which could be generally categorized into operator- and user-oriented
ones. From the operator-oriented perspective, the objectives were set to minimize the operational
cost of repositioning, such as travel costs (Dell'Amico et al., 2016) and time costs (Lee et al.,
2020). From the user-oriented perspective, the objectives related to user dissatisfaction were
optimized, including minimizing the total unmet demand (Shui and Szeto, 2018, Huang et al., 2020,
Jia et al., 2020a), the sum of the absolute deviation of the usable bike inventory after
repositioning from the target level at each station (Rainer-Harbach et al., 2015), and the sum of
the penalty cost incurred at each station (Ho and Szeto, 2014). Apart from the above prevalent
perspectives, some researchers also took environmental influence into account, e.g., total fuel and
CO\textsubscript{2} emission costs (Wang and Szeto, 2018), to minimize the impact of fossil-fueled
vehicle repositioning. It is worth mentioning that a certain number of studies included more than
one of the objectives from above in the problem formulations, which were often handled by the
weighted sum approach (Caggiani and Ottomanelli, 2012, Regue and Recker, 2014, Dell'Amico et al.,
2018, Shui and Szeto, 2018, You, 2019, Chen and Szeto, 2023, Chen et al., 2024, Guo et al., 2024).

\protect\phantomsection\label{p0110}
To solve the aforementioned models, different algorithms were designed. Generally, algorithms
solving BRPs can be classified into exact and inexact methods. For exact algorithms, Contardo et al.
(2012) solved a dynamic repositioning problem for BSSs using column generation combined with Benders
decomposition. Caggiani and Ottomanelli (2012) solved the problem of minimizing the sum of the
relocation costs and lost user costs with the branch-and-bound algorithm. However, as the scale of
the problem grows, exact methods are intractable to provide optimal solutions within an acceptable
time. Hence, heuristics, as one class of inexact methods, have become prevalent in solving such
problems. For instance, You et al. (2019) adopted a two-phase heuristic, in which the first phase
decided on the number of bikes to be picked up and dropped off for all stations, and the second
phase decided on the bike redistribution using the repositioning trucks. Apart from heuristics,
\emph{meta}-heuristics have also been widely applied for BRPs, such as iterated tabu search (Ho and
Szeto, 2014), chemical reaction optimization (Szeto et al., 2016), hybrid large neighborhood search
(Ho and Szeto, 2017), artificial bee colony algorithm (Shui and Szeto, 2018, Jia et al., 2020a, Wang
and Szeto, 2021), and genetic algorithm (Lee et al., 2020). Meanwhile, hybrid approaches combining
exact and inexact methods have also been put forward to make a tradeoff between accuracy and
efficiency, such as hybrid ant-colony optimization with Constraint Programming (CP) (Di Gaspero et
al., 2013a), as well as hybrid Large Neighborhood Search (LNS) with CP (Di Gaspero et al., 2013b).
In more recent years, the reinforcement learning framework has also been utilized to solve BRPs,
particularly in dynamic cases, which can better handle the system's dynamicity and facilitate prompt
decision-making (Chen et al., 2024, Guo et al., 2024, Shi et al., 2024).

\protect\phantomsection\label{s0020}
\subsubsection{2.2. Current research on the BRP with broken bikes}\label{st040}

\protect\phantomsection\label{p0115}
As reviewed above, BRPs have been well-studied with different model formulations and solution
algorithms. Nevertheless, studies on BRPs considering broken bikes were relatively scarce. The first
study was conducted by Kaspi et al. (2016), who proposed a probabilistic model for unusable bicycle
detection. Later, a series of similar studies on BRPs with broken bikes came out with different
problem characteristics.

\protect\phantomsection\label{p0120}
In terms of the research scope, some studies focused on the collection of broken bikes only (Zhang
et al., 2018, Lu et al., 2019, 2022), and some considered both usable and broken bikes
(Alvarez-Valdes et al., 2016, Chang et al., 2018, Chang et al., 2023, Wang and Szeto, 2018, Wang and
Szeto, 2021, Du et al., 2020, Zhang et al., 2020, Cai et al., 2022, Jin et al., 2022). Among the
above studies, most of the literature dealt with docked BSSs (Alvarez-Valdes et al., 2016, Wang and
Szeto, 2018, Zhang et al., 2018, Zhang et al., 2020, Du et al., 2020, Cai et al., 2022, Jin et al.,
2022), with a few studies considering the dockless BSSs (Chang et al., 2018, Chang et al., 2023, Lu
et al., 2019, Usama et al., 2019, 2020). The major difficulty in coping with broken bikes in
dockless BSSs is the scatteration of the broken bikes due to the flexibility of parking in dockless
BSSs (Usama et al., 2020).

\protect\phantomsection\label{p0125}
Regarding the objectives of the existing models, a certain number of existing works set the
minimization of the travel cost for usable bike repositioning and broken bike collection as the
objective (Alvarez-Valdes et al., 2016, Chang et al., 2018, Zhang et al., 2018, Zhang et al., 2020,
Lu et al., 2019). One thing in common among the aforementioned studies is that they required a
perfect balance for usable bikes and a complete collection of broken bikes at each station. However,
given limited vehicle capacity, the perfect balance may require that trucks visit the depot and some
stations multiple times to meet only a tiny fraction of demand, leading to possible waste of
resources or even solution infeasibility, especially under the limited time budget. Therefore, some
works allowed incomplete satisfaction by including minimizing the absolute deviations of the usable
bikes from the target levels into the objective. Specifically, Usama et al. (2019) formulated user
dissatisfaction as the sum of the absolute deviations of inventory levels after rebalancing from the
corresponding target inventory levels at all the stations and minimized it. Wang and Szeto (2021)
introduced the sum of absolute deviations from the target inventory level of usable bikes at each
station and the penalty for the unhandled broken bikes at each station in their objective.
Furthermore, some works also took the environmental issue into account and incorporated the
emissions of CO\textsubscript{2} into their objectives (Wang and Szeto, 2018, Wang and Szeto, 2021,
Chang et al., 2023).

\protect\phantomsection\label{s0025}
\subsubsection{2.3. Research gaps}\label{st045}

\protect\phantomsection\label{p0130}
The existing research on BRPs with broken bikes is summarized in Table 2.1. It can be observed that
the aforementioned studies mainly concentrated on obtaining the minimum absolute deviation of the
inventory level from the target level at the minimum operational cost. Broken bikes were only
regarded as another type of bike in need of collection. However, as mentioned before, some repairing
work can be conducted on-site by repairers, who can handle the broken bikes and satisfy the usable
bike demand simultaneously without vehicle repositioning. To the best of our knowledge, no study has
so far considered the combination of repairers and repositioning vehicles in an optimization model.
We noticed one study on repairer scheduling in a BSS (Freund, 2018), which focused on broken dock
repairs only.

\clearpage
\begin{landscape}
\scriptsize
\setlength{\tabcolsep}{2pt}
\renewcommand{\arraystretch}{0.78}
\setlength{\LTpre}{0pt}
\setlength{\LTpost}{0pt}
\protect\phantomsection\label{t0005}
Table 2.1. A summary of current studies on BRPs with broken bikes.

{\def\LTcaptype{none} % do not increment counter
\begin{longtable}[]{@{}
  >{\raggedright\arraybackslash}p{(\linewidth - 26\tabcolsep) * \real{0.0714}}
  >{\raggedright\arraybackslash}p{(\linewidth - 26\tabcolsep) * \real{0.0714}}
  >{\raggedright\arraybackslash}p{(\linewidth - 26\tabcolsep) * \real{0.0714}}
  >{\raggedright\arraybackslash}p{(\linewidth - 26\tabcolsep) * \real{0.0714}}
  >{\raggedright\arraybackslash}p{(\linewidth - 26\tabcolsep) * \real{0.0714}}
  >{\raggedright\arraybackslash}p{(\linewidth - 26\tabcolsep) * \real{0.0714}}
  >{\raggedright\arraybackslash}p{(\linewidth - 26\tabcolsep) * \real{0.0714}}
  >{\raggedright\arraybackslash}p{(\linewidth - 26\tabcolsep) * \real{0.0714}}
  >{\raggedright\arraybackslash}p{(\linewidth - 26\tabcolsep) * \real{0.0714}}
  >{\raggedright\arraybackslash}p{(\linewidth - 26\tabcolsep) * \real{0.0714}}
  >{\raggedright\arraybackslash}p{(\linewidth - 26\tabcolsep) * \real{0.0714}}
  >{\raggedright\arraybackslash}p{(\linewidth - 26\tabcolsep) * \real{0.0714}}
  >{\raggedright\arraybackslash}p{(\linewidth - 26\tabcolsep) * \real{0.0714}}
  >{\raggedright\arraybackslash}p{(\linewidth - 26\tabcolsep) * \real{0.0714}}@{}}
\toprule\noalign{}
\multirow{2}{=}{\begin{minipage}[b]{\linewidth}\raggedright
Reference
\end{minipage}} & \multirow{2}{=}{\begin{minipage}[b]{\linewidth}\raggedright
Scenario
\end{minipage}} &
\multicolumn{5}{>{\raggedright\arraybackslash}p{(\linewidth - 26\tabcolsep) * \real{0.3571} + 8\tabcolsep}}{%
\begin{minipage}[b]{\linewidth}\raggedright
Objectives
\end{minipage}} & \begin{minipage}[b]{\linewidth}\raggedright
Empty Cell
\end{minipage} &
\multicolumn{3}{>{\raggedright\arraybackslash}p{(\linewidth - 26\tabcolsep) * \real{0.2143} + 4\tabcolsep}}{%
\begin{minipage}[b]{\linewidth}\raggedright
Assumptions
\end{minipage}} & \begin{minipage}[b]{\linewidth}\raggedright
Empty Cell
\end{minipage} &
\multicolumn{2}{>{\raggedright\arraybackslash}p{(\linewidth - 26\tabcolsep) * \real{0.1429} + 2\tabcolsep}@{}}{%
\begin{minipage}[b]{\linewidth}\raggedright
Repositioning quantities
\end{minipage}} \\
& & \begin{minipage}[b]{\linewidth}\raggedright
\textsuperscript{1}AD/\\
\textsuperscript{2}SD\strut
\end{minipage} & \begin{minipage}[b]{\linewidth}\raggedright
\textsuperscript{3}EUDF
\end{minipage} & \begin{minipage}[b]{\linewidth}\raggedright
\textsuperscript{4}TC/TL/TT
\end{minipage} & \begin{minipage}[b]{\linewidth}\raggedright
\textsuperscript{5}MC
\end{minipage} & \begin{minipage}[b]{\linewidth}\raggedright
\textsuperscript{6}CE
\end{minipage} & \begin{minipage}[b]{\linewidth}\raggedright
Empty Cell
\end{minipage} & \begin{minipage}[b]{\linewidth}\raggedright
\textsuperscript{7}MD
\end{minipage} & \begin{minipage}[b]{\linewidth}\raggedright
\textsuperscript{8}MV
\end{minipage} & \begin{minipage}[b]{\linewidth}\raggedright
\textsuperscript{9}OR
\end{minipage} & \begin{minipage}[b]{\linewidth}\raggedright
Empty Cell
\end{minipage} & \begin{minipage}[b]{\linewidth}\raggedright
\textsuperscript{10}UB
\end{minipage} & \begin{minipage}[b]{\linewidth}\raggedright
\textsuperscript{11}BB
\end{minipage} \\
\midrule\noalign{}
\endhead
\bottomrule\noalign{}
\endlastfoot
Kaspi (2015) & Docked, static & & ✓ & ✓ & & & & & & & & \textsuperscript{12}DV & DV \\
Alvarez-Valdes et al. (2016) & Docked, static & & & ✓ & & & & ✓ & ✓ & & & \textsuperscript{13}PB &
PB \\
Chang et al. (2018) & Dockless, static & & & ✓ & ✓ & & & & ✓ & & & \textsuperscript{14}NC & PB \\
Wang and Szeto (2018) & Docked, static & & & & & ✓ & & ✓ & ✓ & & & DV & DV \\
Zhang et al. (2018) & Docked, static & & & ✓ & & & & & ✓ & & & NC & PB \\
Lu et al. (2019) & Dockless, static & & & ✓ & & & & & ✓ & & & NC & PB \\
Xu et al. (2019) & Docked, static & ✓ & & ✓ & & & & & & & & DV & DV \\
Du et al. (2020) & Docked, static & & & ✓ & & & & ✓ & ✓ & & & PB & PB \\
Usama et al. (2020) & Dockless, static & & & ✓ & ✓ & & & & & & & PB & DV \\
Zhang et al. (2020) & Docked, static & & & ✓ & & & & & ✓ & & & PB & PB \\
Chang et al. (2021) & Dockless, dynamic & ✓ & & ✓ & ✓ & & & ✓ & & & & DV & DV \\
Wang and Szeto (2021) & Docked, static & ✓ & & & & ✓ & & & & & & DV & DV \\
Cai et al. (2022) & Docked, dynamic & ✓ & & & & & & & ✓ & & & DV & DV \\
Lu et al. (2022) & Dockless, static & & & ✓ & & & & & & & & NC & PB \\
Jin et al. (2022) & Docked, static/dynamic & & & ✓ & ✓ & & & & ✓ & & & DV & PB \\
Chang et al. (2023) & Dockless, dynamic & ✓ & & & ✓ & ✓ & & ✓ & ✓ & & & PB & PB \\
This study & Docked, static & & ✓ & & & ✓ & & & ✓ & ✓ & & DV & DV \\
\end{longtable}
}
\end{landscape}
\clearpage

\begin{description}
\item[1]
\protect\phantomsection\label{np005}
AD: Absolute Deviation of the number of usable bikes from the target inventory level.
\item[2]
\protect\phantomsection\label{np010}
SD: Satisfied Demand.
\item[3]
\protect\phantomsection\label{np015}
EUDF: Extended User Dissatisfaction Function
\item[4]
\protect\phantomsection\label{np020}
TC/TL/TT: Travel Cost/Travel Length/Travel Time.
\item[5]
\protect\phantomsection\label{np025}
MC: Manual-handling Cost.
\item[6]
\protect\phantomsection\label{np030}
CE: Carbon Emissions.
\item[7]
\protect\phantomsection\label{np035}
MD: Multiple Depots.
\item[8]
\protect\phantomsection\label{np040}
MV: Multiple Vehicles.
\item[9]
\protect\phantomsection\label{np045}
OR: On-site Repairers.
\item[10]
\protect\phantomsection\label{np050}
UB: Usable Bikes.
\item[11]
\protect\phantomsection\label{np055}
BB: Broken Bikes.
\item[12]
\protect\phantomsection\label{np060}
DV: Decision Variable.
\item[13]
\protect\phantomsection\label{np065}
PB: Perfect Balance.
\item[14]
\protect\phantomsection\label{np070}
NC: Not Considered.
\end{description}

\protect\phantomsection\label{p0135}
Regarding the model objective, most researchers focus on minimizing unmet demand for usable bikes
and the number of unhandled broken bikes when considering user-oriented objectives. However, this
approach assumes that all unmet demand and broken bikes at every station and at any time have an
equal impact on user dissatisfaction, which is not always the case in real-world scenarios. For some
more frequently used stations during peak hours, user dissatisfaction should be more sensitive to
unmet demand and unhandled broken bikes than others, and thus, higher priority should be given to
these stations in repositioning. Raviv and Kolka (2013) introduced a User Dissatisfaction Function
(UDF) to measure user dissatisfaction at each station as the expected weighted number of users who
are unable to rent or return a bicycle during a specific period, given an initial inventory at each
station. This measurement is station-based, and it captures both the shortage of bike rent/return
demand and the timing of these events. Several subsequent BRP studies also incorporated the function
into their optimization models (Raviv et al., 2013, Ho and Szeto, 2014, Ho and Szeto, 2017, Forma et
al., 2015, Zhang et al., 2017, Chen and Szeto, 2023). To address the limitation of assuming all the
bikes are usable, Kaspi (2016) extended the UDF by considering the presence of broken bikes,
creating an Extended User Dissatisfaction Function (EUDF). However, to the best of our knowledge,
apart from the study by Kaspi (2015), which proposed a model to minimize the weighted sum of user
dissatisfaction and routing costs, no other studies on BRPs with broken bikes have utilized the
EUDF. Moreover, no study has simultaneously considered both user dissatisfaction and the cost of
total carbon emissions in their models for BRPs with broken bikes, not to mention the corresponding
solution algorithm.

\protect\phantomsection\label{p0140}
Concerning the repositioning quantities, a few studies set the number of broken bikes to be handled
at each station as decision variables and minimizing the unhandled broken bikes in the objective
(Wang and Szeto, 2018, Wang and Szeto, 2021, Usama et al., 2019, Chang et al., 2023). Nevertheless,
none of these studies considered the effect of broken bikes on occupying the docks, which may lead
to a potential failure in returning shared bicycles.

\protect\phantomsection\label{p0145}
This study pioneers in filling these gaps by involving repairer routing and the number of broken
bikes assigned for repairs at every visited station in the optimization model, aiming to minimize
the sum of the penalty costs of the total user dissatisfaction at all stations in the network, and
the cost of the total carbon emissions during the repositioning process. To illustrate the
practicality and feasibility of the proposed model, we developed an effective hybrid algorithm
called HGSADC-SBC, where HGSADC developed~by Vidal et al., 2012, Vidal et al., 2013, Vidal et al.,
2014) serves as the foundation of our method and has been successfully applied to solve a similar
bike repositioning problem involving multiple bike types by Li et al. (2016). The SBC heuristic,
newly proposed in this study, is based on a novel station budget concept that assigns time at each
station by considering the Benefit-Cost Ratio Function (BCRF) for each station in the route. This
allows for a more efficient loading, unloading, and repair quantity assignment across stations in
the routes. Note that the model proposed by Li et al. (2016) focused on a single vehicle and did not
consider repairers, so our method also extends it to handle multiple trucks and repairers.

\protect\phantomsection\label{s0030}
\subsection{3. Model formulation}\label{st050}

\protect\phantomsection\label{p0150}
In this paper, we consider a BSS with one depot and a finite number of stations. Every station has a
fixed capacity, an initial inventory of usable bikes, and an initial inventory of broken bikes. The
numbers of usable bikes and empty docks at each station are assumed to be known as the current
technology of BSSs provides the system operator with this information in real-time. Moreover, the
number of broken bikes at each station is assumed to be known in this study, as we assume that
\emph{users report the existence of broken bikes once they know in the daytime,} and we only
consider a nighttime bike repositioning problem.

\protect\phantomsection\label{p0155}
A fleet of homogeneous repositioning trucks with the same capacity depart from the depot, then
travel among selected stations in the system to deliver usable bikes to visited stations or pick up
the usable bikes from visited stations, collect broken bikes at visited stations, and finally return
to the depot to drop off usable and broken bikes within a given time budget. Throughout this
process, the trucks can return to the depot to drop off usable and broken bikes when needed, but the
number of bikes carried on each truck cannot exceed the truck capacity. Meanwhile, a set of
repairers also depart from the depot, head to stations with broken bikes to conduct repairs directly
on-site, and return to the depot within the same time budget. Loading, unloading, and repairing each
bike take time, with a constant duration for each operation type. For each repositioning truck,
every station and the depot can be visited multiple times. For each repairer, each station can be
visited at most once. A station can be visited by at most one repairer. We assume the depot holds an
infinite inventory of usable bikes and can store an infinite number of usable and broken bikes.
Moreover, we assume that the repairers only repair broken bikes at stations but do not repair broken
bikes at the depot. Broken bikes at the depot are repaired by the repairing staff there but these
bikes have no urgency to be repaired within the repositioning period as we assume we have a
sufficient number of usable bikes in the system. Furthermore, we assume that the repairing staff at
the depot and the onsite repairers are hired at fixed daily salary rates. Therefore, the total
repair staffing cost is constant daily, which is independent of the number of broken bikes repaired
onsite or sent back to the depot within the repositioning period. In addition, we assume that the
repositioning process involves no rental or returning behaviors from the users, and the final
inventories after the repositioning set the starting inventories for the following day.
Additionally, we assume the repairers travel on electric tricycles that generate no emissions during
the repositioning process, and the travel speed of the electric tricycle is proportional to that of
a truck. The objective of the BRP is to minimize the cost consisting of the following components:

\begin{itemize}
\item
  \protect\phantomsection\label{p0160}
  The sum of user dissatisfaction at all stations on the following day and;
\item
  \protect\phantomsection\label{p0165}
  The CO\textsubscript{2} emissions of the repositioning trucks during the rebalancing process.
\end{itemize}

\protect\phantomsection\label{p0170}
Note that the daily repair staffing cost is not included in our optimization objective as it does
not affect the optimal repositioning decisions.

\protect\phantomsection\label{p0175}
Fig. 3.1 illustrates an example of the problem. The BSS consists of one depot (represented by a red
star and labeled with 0) and ten stations (labeled from 1 to 10). The tables in Fig. 3.1
\textbf{(a)}-\textbf{(d)} show the number of usable and broken bikes loaded onto a truck and
unloaded from a truck at each node, and the number of bikes repaired by repairers at each node. Two
trucks (Truck 1 in Fig. 3.1 \textbf{(a)} and Truck 2 in Fig. 3.1 \textbf{(b)}) with the identical
capacity of 25 bikes and two repairers (Repairer 1 in Fig. 3.1 \textbf{(c)} and Repairer 2 in Fig.
3.1 \textbf{(d)}) depart from the depot to visit these stations. Specifically, Truck 1 follows the
route 0--4--6--9--0--2--0, with one intermediate return to the depot to drop off 18 broken bikes and
pick up three usable bikes. Truck 2 follows the route 0--5--4--0--7--5--9--0, with one intermediate
return to the depot to drop off eight usable bikes and 13 broken bikes. Repairer 1 follows the route
0--6--8--0, performing on-site repairs before returning to the depot, while Repairer 2 follows the
route 0--1--3--0. Node 5 is visited by Truck 2 twice because the total number of broken and surplus
usable bikes to be loaded at the station (26) is more than the truck capacity (25). To reduce
vehicle emissions, it is wise to drop off all broken bikes when Truck 2 visits the depot during the
repositioning process. The travel times between nodes are labeled on the arcs in the figures, and
the travel times between nodes are not assumed to be symmetric (e.g., the back-and-forth travel
times for Truck 1 between nodes 0 to 2 are different). Note that Station 10 is not visited by both
trucks and repairers.

\begin{figure}
\centering
Fig. 3.1. Problem illustration
\caption{}\label{f0005}
\end{figure}

\protect\phantomsection\label{p0180}
Fig. 3.2 shows the repositioning and repair process of Truck 1 and Repairer 1 from the example in
Fig. 3.1 from a temporal perspective. The repositioning time budget is 7200~s as marked on the time
axis. The figure shows the routes taken, inventory variations at the stations, and temporal aspects
of the operations (e.g., loading and repairing times, arrival times). Loading or unloading one bike
takes 60~s while repairing one broken bike takes 300~s. Note that (\emph{p, b}) \ensuremath{\rightarrow} (\emph{p', b'}) at
each station represents the variations of usable bike inventory from \emph{p} to \emph{p'} and
broken bike inventory from \emph{b} to \emph{b'} after the operations of trucks or
repairers.~\textless~\emph{A}, \emph{B}~\textgreater~on the axis represents the truck inventory of
\emph{A} usable bikes and \emph{B} broken bikes when it travels between the nodes. The
round-cornered shapes show the variation of the user dissatisfaction under the inventories before
and after the operations. Specifically, Truck 1 follows the route 0--4--6--9--0--2--0 as indicated
above the time axis. It first visits Station 4 to collect two usable bikes and one broken bike, then
travels to Station 6 to collect five usable and 17 broken bikes, reaching its full capacity of 25
bikes. Next, the truck travels to Station 9 to unload the seven usable bikes. Following that, it
returns to the depot to unload all broken bikes, loads three usable bikes, and proceeds to Station 2
to deliver the three usable bikes and collect nine broken bikes. Finally, it returns to the depot,
unloading all broken bikes and completing its operations in 7123~s, calculated by adding up the
arrival time at the depot (6583~s) and the operational time for unloading the nine broken bikes
(540~s). Meanwhile, Repairer 1 performs on-site repairs along the route 0--6--8--0, repairing four
broken bikes at Station 6 and 12 at Station 8. It concludes its operations in 7180~s. Note that
Station 6 is visited by both Truck 1 and Repairer 1, and all 21 broken bikes at the station are
handled, with four repaired by the repairer before the truck arrives at the station and 17 loaded by
the truck. The four repaired bikes, together with one usable bike originally at Station 6 and two
usable bikes at Station 4 are used to partially address the deficit of bikes at Station 9. The
remaining deficit of eight usable bikes at Station 9 is satisfied by the usable bikes collected by
Truck 2 from other stations and delivered to the station.

\begin{figure}
\centering
Fig. 3.2. The repositioning and repair process of Truck 1 and Repairer 1.
\caption{}\label{f0010}
\end{figure}

\protect\phantomsection\label{p0185}
The first component of the objective function of the problem involves the calculation of user
dissatisfaction throughout the day following the rebalancing process based on the starting
inventories resulting from that process. The method proposed by Kaspi et al. (2017) is used for this
calculation, which is briefly introduced here. Given station \emph{i} with \emph{C\textsubscript{i}}
docks and a finite period for the following day, at the beginning of that day, there are
\emph{p\textsubscript{i}} usable bikes and \emph{b\textsubscript{i}} broken bikes present at the
station. Throughout this timeframe, users intending to rent or return a bike arrive at the station
based on a stochastic process. When a user's demand for a bike or an available dock is met, the bike
is rented or returned, subsequently updating the station's inventory level. Conversely, if the
demand cannot be fulfilled, it is assumed that the user promptly leaves the station, corresponding
to a failure in renting or returning a bike. The user dissatisfaction at a station is calculated as
the weighted expected number of occurrences where users are unable to rent or return bikes during
the specified period, provided that the usable and broken bike inventories at the beginning of the
period are \emph{p\textsubscript{i}} and \emph{b\textsubscript{i}}, respectively. It is assumed that
changes in inventory levels are solely caused by users' rental or return actions, with no external
interventions such as truck repositioning or repairs. An Extended User Dissatisfaction Function
(EUDF), denoted as \(F_{i}\left( {p_{i},b_{i}} \right)\), is defined to represent the
dissatisfaction at station \emph{i} given the inventories \emph{p\textsubscript{i}} and
\emph{b\textsubscript{i}}. The domain of this bivariate function is
\(\Theta_{i} = \left\{ {\left( {p_{i},b_{i}} \right) \mid p_{i} \geq 0,b_{i} \geq 0,p_{i} + b_{i} \leq C_{i}} \right\}\).
Kaspi et al. (2017) proposed a method for determining a convex polyhedral function that has nearly
identical values to the EUDF values in its range, enabling the integration of the EUDF into the
subsequent linear optimization model. Specifically, for each
\(\left( {p_{i},b_{i}} \right) \in \Theta_{i},\) a plane
\(\alpha_{p_{i},b_{i}}^{i} \cdot p_{i} + \beta_{p_{i},b_{i}}^{i} \cdot b_{i} + \gamma_{p_{i},b_{i}}^{i}\)
that passes as close as possible to \(F_{i}\left( {p_{i},b_{i}} \right)\) can be fitted, where
\(\alpha_{p_{i},b_{i}}^{i},\beta_{p_{i},b_{i}}^{i},\) and \(\gamma_{p_{i},b_{i}}^{i}\) are the
coefficients of the plane that can be computed for every
\(\left( {p_{i},b_{i}} \right) \in \theta_{i}\). The computation details of the coefficients of the
fitted plane are given in Appendix A. Due to limited space, we omit the lengthy details on
generating the station EUDF values. Interested readers can consult Kaspi et al. (2017) for
generating them.

\protect\phantomsection\label{p0190}
The second component requires the calculation of the CO\textsubscript{2} emissions of a
repositioning truck as it moves. Wang and Szeto (2021) used the following emission model to
calculate the amount of CO\textsubscript{2} emissions \emph{e} produced by a truck when it carries
bikes over a
distance:\[\tag{1}e = CER \cdot \left( {\rho_{0} + \frac{\rho^{\ast} - \rho_{0}}{Q}\psi} \right) \cdot d\]where
\emph{CER} is the CO\textsubscript{2} emission rate, \(\rho^{\ast}\) and \(\rho_{0}\) are the fuel
consumption rates when the vehicle is empty and fully loaded, respectively, \emph{Q} is the vehicle
capacity, \(\psi\) is the number of bikes carried by the truck, and \emph{d} is the distance
traveled by the truck. Given the similarities between their and our research context, we chose to
adopt this emission model for the calculation in the second component.

\protect\phantomsection\label{p0195}
As aforementioned, each station can be visited by each truck an arbitrary number of times. Hence,
the classic arc-indexed formulation is not suitable for our problem. As proposed by Raviv et al.
(2013), the time-indexed formulation requires the whole time budget \emph{T} to be discretized into
multiple periods with equal length, which allows us to keep track of each visit to a station and the
inventory of bikes on each truck and at each station during every period, thereby enabling the
formulation (1) for multiple visits to a station and the depot by vehicles and repairers and (2) to
decide the pickup quantities of usable bikes by trucks and the repair quantities of broken bikes by
repairs based on the actual inventory status. Therefore, we formulate our problem based on the
time-indexed formulation.

\protect\phantomsection\label{s0035}
\subsubsection{3.1. Notations}\label{st055}

\protect\phantomsection\label{p0200}
\protect\phantomsection\label{t0080}
{\def\LTcaptype{none} % do not increment counter
\begin{longtable}[]{@{}ll@{}}
\toprule\noalign{}
\endhead
\bottomrule\noalign{}
\endlastfoot
\multicolumn{2}{@{}l@{}}{%
\textbf{Sets}} \\
\emph{S} & The set of stations, indexed by \emph{i~=}~1, \ldots, \textbar{}\emph{S}\textbar{} \\
\emph{S}\textsubscript{0} & The set of nodes, including the stations and the depot (denoted by 0) \\
\emph{K} & The set of rebalancing trucks \\
\emph{R} & The set of repairers \\
\multicolumn{2}{@{}l@{}}{%
\textbf{Parameters}} \\
\emph{p\textsubscript{i}}\textsuperscript{0} & The initial inventory of usable bikes at node
\(i \in S_{0}\) \\
\emph{b\textsubscript{i}}\textsuperscript{0} & The initial quantity of broken bikes at node
\(i \in S_{0}\) \\
\emph{ρ}\textsubscript{0} & The empty-load fuel consumption rate of each truck (liters/km) \\
\emph{ρ*} & The full-load fuel consumption rate of each truck (liters/km) \\
\emph{CER} & The CO\textsubscript{2} emission rate (kg/liter) \\
\emph{Q\textsubscript{k}} & The capacity of truck \emph{k} \\
\emph{C\textsubscript{i}} & The capacity of node \(i \in S_{0}\) \\
\(\tau\) & The length of a discretized period (second) \\
\(t^{\text{load}}\) & Time for loading/unloading a bike (\(t^{\text{load}} < \tau\)) \\
\(t^{\text{repair}}\) & Time for repairing a bike on-site (\(t^{\text{repair}} < \tau\)) \\
\emph{T} & Rebalancing time budget (second) \\
\(t_{\mathit{ij}}^{\text{v}}\) & Travel time for a truck from node \emph{i} to node \emph{j} \\
\(t_{\mathit{ij}}^{\text{r}}\) & Travel time for a repairer from node \emph{i} to node \emph{j} \\
\(t_{\mathit{ij}}'\) & The number of periods required by a truck to travel from node \emph{i} to
node \emph{j}, which is calculated by \(t_{\mathit{ij}}' = \left\{ \begin{array}{lc}
{\left\lceil t_{\mathit{ij}}^{\text{v}}/\tau \right\rceil,} & {\text{if}\; i \neq j} \\
{1,} & {\text{if}\; i = j}
\end{array} \right.\) \\
\(t_{\mathit{ij}}^{''}\) & The number of periods required by a repairer to travel from node \emph{i}
to node \emph{j}, which is calculated by \(t_{\mathit{ij}}^{''} = \left\{ \begin{array}{lc}
{\left\lceil t_{\mathit{ij}}^{\text{r}}/\tau \right\rceil,} & {\text{if}\; i \neq j} \\
{1,} & {\text{if}\; i = j}
\end{array} \right.\) \\
\(d_{\mathit{ij}}\) & The travel distance from node \emph{i} to node \emph{j} (km) \\
\(\alpha_{p_{i},b_{i}}^{i}\) & The first coefficient of the fitted plane to approximate the EUDF at
station \emph{i}, given the usable bike inventory level of \emph{p\textsubscript{i}} and the broken
bike inventory level of \emph{b\textsubscript{i}} \\
\(\beta_{p_{i},b_{i}}^{i}\) & The second coefficient of the fitted plane to approximate the EUDF at
station \emph{i}, given the usable bike inventory level of \emph{p\textsubscript{i}} and the broken
bike inventory level of \emph{b\textsubscript{i}} \\
\(\gamma_{p_{i},b_{i}}^{i}\) & The third coefficient of the fitted plane to approximate the EUDF at
station \emph{i}, given the usable bike inventory level of \emph{p\textsubscript{i}} and the broken
bike inventory level of \emph{b\textsubscript{i}} \\
\(T'\) & The number of periods in the rebalancing time budget such that
\(T' = \left\lfloor T/\tau \right\rfloor\) \\
\(c^{\text{p}}\) & The penalty cost for a user unable to rent/return a bike (CNY) \\
\(c^{\text{e}}\) & The unit carbon trading cost (CNY/kg) \\
\multicolumn{2}{@{}l@{}}{%
\textbf{Decision variables}} \\
\(x_{i,j,t,k}^{\text{v}}\) & If truck \emph{k} travels directly from node \emph{i} to node \emph{j}
during period \emph{t}, \(x_{i,j,t,k}^{\text{v}}\)~=~1, 0 otherwise \\
\(x_{i,j,t,m}^{\text{r}}\) & If repairer \emph{m} travels directly from node \emph{i} to node
\emph{j} during period \emph{t}, \(x_{i,j,t,m}^{\text{r}}\)~=~1, 0 otherwise \\
\(p_{i,t}\) & The number of usable bikes at node \emph{i} at the end of period \emph{t} \\
\(b_{i,t}\) & The number of broken bikes at node \emph{i} at the end of period \emph{t} \\
\(y_{i,t,k}^{\text{u}^{+}}\) & The number of usable bikes delivered to node \emph{i} by truck
\emph{k} during period \emph{t} \\
\(y_{i,t,k}^{\text{u}^{-}}\) & The number of usable bikes collected from node \emph{i} by truck
\emph{k} during period \emph{t} \\
\(y_{i,t,k}^{\text{b}^{+}}\) & The number of broken bikes delivered to node \emph{i} by truck
\emph{k} during period \emph{t} \\
\(y_{i,t,k}^{\text{b}^{-}}\) & The number of broken bikes collected from node \emph{i} by truck
\emph{k} during period \emph{t} \\
\(g_{i,t,m}\) & The number of broken bikes repaired by repairer \emph{m} at node \emph{i} during
period \emph{t} \\
\(\psi_{i,j,t,k}^{\text{u}}\) & The number of usable bikes on truck \emph{k} when it departs from
node \emph{i} to node \emph{j} during period \emph{t} \\
\(\psi_{i,j,t,k}^{\text{b}}\) & The number of broken bikes on truck \emph{k} when it departs from
node \emph{i} to node \emph{j} during period \emph{t} \\
\(e_{i,j,t,k}\) & The CO\textsubscript{2} emissions generated when truck \emph{k} travels from node
\emph{i} to node \emph{j} during period \emph{t} \\
\(\partial_{k}^{\text{v}}\) & The period in which truck \emph{k} returns to the depot and remains
stationary until the rebalancing time budget is used up \\
\(\partial_{m}^{\text{r}}\) & The period in which repairer \emph{m} returns to the depot and remains
stationary until the rebalancing time budget is used up \\
\(\varsigma_{t,k}^{\text{v}}\) & If truck \emph{k} remains stationary at the depot from period \emph{t} to
the end of the repositioning process, \(\varsigma_{t,k}^{\text{v}} = 1\), 0 otherwise \\
\(\varsigma_{t,m}^{\text{r}}\) & If repairer \emph{m} remains stationary at the depot from period \emph{t}
to the end of the repositioning process, \(\varsigma_{t,m}^{\text{r}} = 1\), 0 otherwise \\
\({\widehat{F}}_{i}\) & The auxiliary variable for the linearization of the EUDF of station
\emph{i} \\
\multicolumn{2}{@{}l@{}}{%
Function} \\
\(F_{i}\left( {p_{i},b_{i}} \right)\) & The EUDF defined over the usable bike inventory level
\emph{p\textsubscript{i}} and broken bike inventory level \emph{b\textsubscript{i}} at station
\emph{i} \\
\end{longtable}
}

\protect\phantomsection\label{s0040}
\subsubsection{3.2. Time-indexed formulation}\label{st060}

\protect\phantomsection\label{p0205}
The time-indexed formulation of the BRP separating the routes of trucks and repairers is as follows.

\protect\phantomsection\label{p0210}
\[\tag{2}\begin{array}{l}
{\text{minimize~}c^{\text{p}}\sum_{i \in S}F_{i}\left( {p_{i,T^{'}},b_{i,T^{'}}} \right) + c^{\text{e}}\sum_{i \in S_{0}}\sum_{j \in S_{0}}\sum_{t = t_{i,j}^{'} + 1}^{T^{'}}\sum_{k \in K}e_{i,j,t - t_{i,j}^{'}k} +} \\
{\theta \cdot \left( \begin{array}{l}
{\sum_{i \in S_{0}}\sum_{j \in S_{0}}\sum\limits_{t = t_{i,j}^{'} + 1}^{T^{'}}\sum_{k \in K}\left( {t_{i,j}^{\text{v}} \cdot x_{i,j,t - t_{i,j}^{'},k}^{\text{v}}} \right) + \sum_{i \in S_{0}}\sum_{t = t_{i,j}^{'} + 1}^{T^{'}}\sum_{k \in K}t^{\text{load}} \cdot \left( {y_{i,t,k}^{\text{u}^{+}} + y_{i,t,k}^{\text{u}^{-}} + y_{i,t,k}^{\text{b}^{+}} + y_{i,t,k}^{\text{b}^{-}}} \right) +} \\
{\sum_{i \in S_{0}}\sum_{j \in S_{0}}\sum\limits_{t = t_{i,j}^{''} + 1}^{T^{'}}\sum_{m \in R}\left( {t_{i,j}^{\text{r}} \cdot x_{i,j,t - t_{i,j}^{''},m}^{\text{r}}} \right) + \sum_{i \in S_{0}}\sum_{t = t_{i,j}^{''} + 1}^{T^{'}}\sum_{m \in R}t^{\text{repair}} \cdot g_{i,t,m}}
\end{array} \right)}
\end{array}\]Subject to

\[\tag{3}p_{i,1} = p_{i}^{0} - \sum\limits_{k \in K}y_{i,1,k}^{\text{u}^{-}} + \sum\limits_{k \in K}y_{i,1,k}^{\text{u}^{+}} + \sum\limits_{m \in R}g_{i,1,m},\forall i \in S_{0},\]

\[\tag{4}p_{i,t} = p_{i,t - 1} - \sum\limits_{k \in K}y_{i,t,k}^{\text{u}^{-}} + \sum\limits_{k \in K}y_{i,t,k}^{\text{u}^{+}} + \sum\limits_{m \in R}g_{i,t,m},\forall i \in S_{0},t = 2,...,T',\]

\[\tag{5}b_{i,1} = b_{i}^{0} - \sum\limits_{k \in K}y_{i,1,k}^{\text{b}^{-}} + \sum\limits_{k \in K}y_{i,1,k}^{\text{b}^{+}} - \sum\limits_{m \in R}g_{i,1,m},\forall i \in S_{0},\]

\[\tag{6}b_{i,t} = b_{i,t - 1} - \sum\limits_{k \in K}y_{i,t,k}^{\text{b}^{-}} + \sum\limits_{k \in K}y_{i,t,k}^{\text{b}^{+}} - \sum\limits_{m \in R}g_{i,t,m},\forall i \in S_{0},t = 2,\ldots,T',\]

\[\tag{7}p_{i,t} + b_{i,t} \leq C_{i},\forall i \in S,t = 1,....T',\]

\[\tag{8}\sum\limits_{j \in S_{0}}\psi_{0,j,1,k}^{\text{u}} = y_{0,1,k}^{\text{u}^{-}},\forall k \in K,\]

\[\tag{9}\sum\limits_{j \in S_{0}}\psi_{i,j,t,k}^{\text{u}} = \sum\limits_{j \in S_{0},t - t_{j,i}' > 0}\psi_{j,i,t - t_{j,i}',k}^{\text{u}} - y_{i,t,k}^{\text{u}^{+}} + y_{i,t,k}^{\text{u}^{-}},\forall i \in S_{0},t = 2,...,T',k \in K,\]

\[\tag{10}\psi_{i,j,T',k}^{\text{u}} = 0,\forall i \in S_{0},j \in S_{0},k \in K,\]

\[\tag{11}\psi_{i,j,t,k}^{\text{u}} \leq Q_{k}x_{i,j,t,k}^{\text{v}},\forall i \in S_{0},j \in S_{0},t = 1,...,T',k \in K,\]

\[\tag{12}\psi_{i,j,1,k}^{\text{b}} = 0,\forall i \in S_{0},j \in S_{0},k \in K,\]

\[\tag{13}\sum\limits_{j \in S_{0}}\psi_{i,j,t,k}^{\text{b}} = \sum\limits_{j \in S_{0},t - t_{j,i}' > 0}\psi_{j,i,t - t_{j,i}',k}^{\text{b}} - y_{i,t,k}^{\text{b}^{+}} + y_{i,t,k}^{\text{b}^{-}},\forall i \in S_{0},t = 2,...,T',k \in K,\]

\[\tag{14}\psi_{i,j,T',k}^{\text{b}} = 0,\forall i \in S_{0},j \in S_{0},k \in K,\]

\[\tag{15}\psi_{i,j,t,k}^{\text{b}} \leq Q_{k}x_{i,j,t,k}^{\text{v}},\forall i \in S_{0},j \in S_{0},t = 1,...,T',k \in K,\]

\[\tag{16}\psi_{i,j,t,k}^{\text{u}} + \psi_{i,j,t,k}^{\text{b}} \leq Q_{k},\forall i \in S_{0},j \in S_{0},t = 1,...,T',k \in K,\]

\[\tag{17}y_{i,t,k}^{\text{u}^{+}} \leq \min\left\{ {C_{i},Q_{k}} \right\} \cdot \sum\limits_{j \in S_{0}}x_{i,j,t,k}^{\text{v}},\forall i \in S_{0},t = 1,...,T',k \in K,\]

\[\tag{18}y_{i,t,k}^{\text{u}^{-}} \leq \min\left\{ {C_{i},Q_{k}} \right\} \cdot \sum\limits_{j \in S_{0}}x_{i,j,t,k}^{\text{v}},\forall i \in S_{0},t = 1,...,T',k \in K,\]

\[\tag{19}y_{0,t,k}^{\text{b}^{+}} \leq Q_{k} \cdot \sum\limits_{j \in S_{0}}x_{0,j,t,k}^{\text{v}},\forall t = 1,...,T',k \in K,\]

\[\tag{20}y_{i,t,k}^{\text{b}^{+}} = 0,\forall i \in S,t = 1,...,T',k \in K,\]

\[\tag{21}y_{0,t,k}^{\text{b}^{-}} = 0,\forall t = 1,...,T',k \in K,\]

\[\tag{22}y_{i,t,k}^{\text{b}^{-}} \leq \min\left\{ {b_{i}^{0},Q_{k}} \right\} \cdot \sum\limits_{j \in S_{0}}x_{i,j,t,k}^{\text{v}},\forall i \in S,t = 1,...,T',k \in K,\]

\[\tag{23}g_{0,t,m} = 0,\forall t = 1,...,T',m \in R,\]

\[\tag{24}g_{i,t,m} \leq b_{i}^{0} \cdot \sum\limits_{j \in S_{0}}x_{i,j,t,m}^{\text{r}},\forall i \in S,t = 1,...,T',m \in R,\]

\[\tag{25}\sum\limits_{t = 1}^{T'}{\sum\limits_{k \in K}y_{i,t,k}^{\text{b}^{-}}} + \sum\limits_{t = 1}^{T'}{\sum\limits_{m \in R}g_{i,t,m}} \leq b_{i}^{0},\forall i \in S,\]

\[\tag{26}\sum\limits_{j \in S_{0}}x_{0,j,1,k}^{\text{v}} = 1,\forall k \in K,\]

\[\tag{27}\sum\limits_{j \in S}x_{0,j,1,m}^{\text{r}} = 1,\forall m \in R,\]

\[\tag{28}\sum\limits_{i \in S_{0}}x_{i,0,T^{'} - t_{i,0}^{'},k}^{\text{v}} = 1,\forall k \in K,\]

\[\tag{29}\sum\limits_{i \in S_{0}}x_{i,0,T^{'} - t_{i,0}^{''},m}^{\text{r}} = 1,\forall m \in R,\]

\[\tag{30}\sum\limits_{j \in S_{0},t - t_{\mathit{ji}}' > 0}x_{j,i,t - t_{\mathit{ji}}',k}^{\text{v}} = \sum\limits_{u \in S_{0}}x_{i,u,t,k}^{\text{v}},\forall i \in S_{0},t = 2,...,T',k \in K,\]

\[\tag{31}\sum\limits_{j \in S_{0},t - t_{\mathit{ji}}^{''} > 0}x_{j,i,t - t_{\mathit{ji}}^{''},m}^{\text{r}} = \sum\limits_{u \in S_{0}}x_{i,u,t,m}^{\text{r}},\forall i \in S_{0},t = 2,...,T',m \in R,\]

\[\tag{32}\left. \varsigma_{t,k}^{\text{v}} \geq \left( t - \partial_{k}^{\text{v}} + 1 \right)/T',\forall t = 1,\ldots,T',k \in K, \right.\]

\[\tag{33}\left. \varsigma_{t,k}^{\text{v}} \leq \left( t - \partial_{k}^{\text{v}} + 0.5 \right)/T' + 1,\forall t = 1,\ldots,T',k \in K, \right.\]

\[\tag{34}\left. \varsigma_{t,m}^{\text{r}} \geq \left( t - \partial_{m}^{\text{r}} + 1 \right)/T',\forall t = 1,\ldots,T',m \in R, \right.\]

\[\tag{35}\left. \varsigma_{t,m}^{\text{r}} \leq \left( t - \partial_{m}^{\text{r}} + 0.5 \right)/T' + 1,\forall t = 1,\ldots,T',m \in R, \right.\]

\[\tag{36}\sum\limits_{i,j \in S_{0}}\sum\limits_{t \leq w,t > t_{ij}^{'}}t_{i,j}^{\text{v}} \cdot x_{i,j,t - t_{ij}^{'},k}^{\text{v}} + \sum\limits_{i \in S_{0}}\sum\limits_{t \leq w}t^{\text{load}}\left( {y_{i,t,k}^{\text{u}^{+}} + y_{i,t,k}^{\text{u}^{-}} + y_{i,t,k}^{\text{b}^{+}} + y_{i,t,k}^{\text{b}^{-}}} \right) \leq w \cdot \tau,\forall w = 1,...,T',k \in K,\]

\[\tag{37}\sum\limits_{i,j \in S_{0}}\sum\limits_{t \leq w,t > 0}t_{i,j}^{\text{v}} \cdot x_{i,j,t,k}^{\text{v}} + \sum\limits_{i \in S_{0}}\sum\limits_{t \leq w}t^{\text{load}}\left( {y_{i,t,k}^{\text{u}^{+}} + y_{i,t,k}^{\text{u}^{-}} + y_{i,t,k}^{\text{b}^{+}} + y_{i,t,k}^{\text{b}^{-}}} \right) \geq \left( {w - 2} \right) \cdot \tau - \varsigma_{w,k}^{\text{v}}T,\forall w = 1,...,T^{'},k \in K,\]

\[\tag{38}\sum\limits_{i,j \in S_{0}}{\sum\limits_{t \leq w,t > t_{\mathit{ij}}^{''}}t_{i,j}^{\text{r}}} \cdot x_{i,j,t - t_{\mathit{ij}}^{''},m}^{\text{r}} + \sum\limits_{i \in S}{\sum\limits_{t \leq w}{t^{\text{repair}}g_{i,t,m}}} \leq w \cdot \tau,\forall w = 1,...,T',m \in R,\]

\[\tag{39}\sum\limits_{i,j \in S_{0}}{\sum\limits_{t \leq w,t > 0}t_{i,j}^{\text{r}}} \cdot x_{i,j,t,m}^{\text{r}} + \sum\limits_{i \in S}{\sum\limits_{t \leq w}{t^{\text{repair}}g_{i,t,m}}} \geq \left( {w - 2} \right) \cdot \tau - \varsigma_{w,m}^{\text{r}}T,\forall w = 1,...,T',m \in R,\]

\[\tag{40}x_{0,0,t,k}^{\text{v}} \leq 1 - \left( \varsigma_{t + 1,k}^{\text{v}} - \varsigma_{t,k}^{\text{v}} \right),\forall t = 1,\ldots,T' - 1,k \in K,\]

\[\tag{41}x_{0,0,t,m}^{\text{r}} \leq 1 - \left( \varsigma_{t + 1,m}^{\text{r}} - \varsigma_{t,m}^{\text{r}} \right),\forall t = 1,\ldots,T' - 1,m \in R,\]

\[\tag{42}x_{i,i,t,k}^{\text{v}} \leq \left( {y_{i,t,k}^{\text{u}^{+}} + y_{i,t,k}^{\text{u}^{-}} + y_{i,t,k}^{\text{b}^{+}} + y_{i,t,k}^{\text{b}^{-}}} \right) + \left\lfloor \frac{\tau}{t^{load}} \right\rfloor \varsigma_{t,k}^{\text{v}},\forall i \in S_{0},t = 1,\ldots,T^{'},k \in K,\]

\[\tag{43}x_{i,i,t,m}^{\text{r}} \leq g_{i,t,m} + \left\lfloor \frac{\tau}{t^{\text{repair}}} \right\rfloor \varsigma_{t,m}^{\text{r}},\forall i \in S_{0},t = 1,\ldots,T^{'},m \in R,\]

\[\tag{44}\sum\limits_{i \in S_{0}}{\sum\limits_{j \in S_{0}}x_{i,j,t,k}^{\text{v}}} \leq 1,\forall t = 1,...,T',k \in K,\]

\[\tag{45}\sum\limits_{i \in S_{0}}{\sum\limits_{j \in S_{0}}x_{i,j,t,m}^{\text{r}}} \leq 1,\forall t = 1,...,T',m \in R,\]

\[\tag{46}\sum\limits_{j \in S_{0},j \neq i}{\sum\limits_{t = t_{j,i}^{''} + 1}^{T'}{\sum\limits_{m \in R}x_{j,i,t - t_{j,i}^{''},m}^{\text{r}}}} \leq 1,\forall i \in S_{0},\]

\[\tag{47}e_{i,j,t,k} \geq CER \cdot \left\lbrack {\rho_{0} \cdot x_{i,j,t,k}^{\text{v}} + \frac{\rho^{\ast} - \rho_{0}}{Q_{k}} \cdot \left( {\psi_{i,j,t,k}^{\text{u}} + \psi_{i,j,t,k}^{\text{b}}} \right)} \right\rbrack \cdot d_{\mathit{ij}},\mspace{100mu}\mspace{100mu}\forall i,j \in S_{0},t = 1,...,T',k \in K,\]

\[\tag{48}x_{i,j,t,k}^{\text{v}} \in \left\{ 0,1 \right\},\forall i,j \in S_{0},t = 1,...,T',k \in K,\]

\[\tag{49}x_{i,j,t,m}^{\text{r}} \in \left\{ 0,1 \right\},\forall i,j \in S_{0},t = 1,...,T',m \in R,\]

\[\tag{50}\varsigma_{t,k}^{\text{v}} \in \left\{ 0,1 \right\},\forall t = 1,...,T',k \in K,\]

\[\tag{51}\varsigma_{t,m}^{\text{r}} \in \left\{ 0,1 \right\},\forall t = 1,...,T',m \in R,\]

\[\tag{52}p_{i,t} \geq 0,\forall i \in S_{0},t = 1,...,T',\]

\[\tag{53}b_{i,t} \geq 0,\forall i \in S_{0},t = 1,...,T',\]

\[\tag{54}y_{i,t,k}^{\text{u}^{+}} \geq 0,\text{integer},\forall i \in S_{0},t = 1,...,T',k \in K,\]

\[\tag{55}y_{i,t,k}^{\text{u}^{-}} \geq 0,\text{integer},\forall i \in S_{0},t = 1,...,T',k \in K,\]

\[\tag{56}y_{i,t,k}^{\text{b}^{+}} \geq 0,\text{integer},\forall i \in S_{0},t = 1,...,T',k \in K,\]

\[\tag{57}y_{i,t,k}^{\text{b}^{-}} \geq 0,\text{integer},\forall i \in S_{0},t = 1,...,T',k \in K,\]

\[\tag{58}g_{i,t,m} \geq 0,\text{integer},\forall i \in S_{0},t = 1,...,T',m \in R,\]

\[\tag{59}\partial_{k}^{\text{v}} \geq 0,\text{integer},\forall k \in K,\]

\[\tag{60}\partial_{m}^{\text{r}} \geq 0,\text{integer},\forall m \in R,\]

\[\tag{61}\partial_{k}^{\text{v}} \leq T',\text{integer},\forall k \in K,\]

\[\tag{62}\partial_{m}^{\text{r}} \leq T',\text{integer},\forall m \in R,\]

\[\tag{63}\psi_{i,j,t,k}^{\text{u}} \geq 0,\forall i,\mspace{100mu} j \in S_{0},t = 1,...,T',k \in K,\]

\[\tag{64}\psi_{i,j,t,k}^{\text{b}} \geq 0,\forall i,\mspace{100mu} j \in S_{0},t = 1,...,T',k \in K,\quad\mspace{600mu}\text{and}\]

\[\tag{65}e_{i,j,t,k} \geq 0,\forall i,\mspace{100mu} j \in S_{0},t = 1,...,T',k \in K.\]

\protect\phantomsection\label{p0220}
The objective (2) is to minimize the sum of the penalty cost of the total user dissatisfaction at
all stations and the cost of the total CO\textsubscript{2} emissions of all the trucks. Note that we
add a third term to penalize the sum of the total travel time of trucks and repairers, bike
loading/unloading times for trucks, and repairing times for the repairers. The inclusion of the
third term helps discourage non-intuitive or redundant operations that, while feasible, may not be
necessary given the available resources and time. For example, as shown in Fig. 3.3 \textbf{(a)},
the first case illustrates a truck loading eight usable bikes at a station in period
\emph{t}\textsubscript{0} and then unloading two bikes at the same station in the next period
\emph{t}\textsubscript{0}~+~1. This solution is feasible and could be optimal when the third term in
the objective was not included, but such an operation is redundant as it can be simplified by
loading only six bikes at once. The second case in Fig. 3.3 \textbf{(b)} shows a repairer repairing
one broken bike at one station, traveling to another station to repair one broken bike, and then
returning to the original station to repair one additional bike. Again, while feasible, this
solution is non-intuitive because the unnecessary return visit to the same station can be avoided.
Overall, the third term in the objective penalizes such operations, ensuring a more straightforward
and efficient rebalancing process. To avoid introducing unnecessary bias when the experiments
involve the scenario without repairers, we set the coefficient \emph{θ} for this penalty term to be
1e-8. This ensures that the third term does not affect solution optimality. The largest possible
value for this term is \(\left( \middle| K \middle| + \middle| R \middle| \right)T\), and under our
experimental setting, the maximum value of the term is 43,200 (\textbar{}\emph{K}\textbar{} = 2,
\textbar{}\emph{R}\textbar{} = 2, and \emph{T} = 10800), which is small enough to avoid any
noticeable impact on results that are rounded to two decimal places.

\begin{figure}
\centering
Fig. 3.3. Examples of using the third component in the objective to avoid redundant operations.
\caption{}\label{f0015}
\end{figure}

\protect\phantomsection\label{p0225}
Constraints (3), (4), (5), (6) keep track of the numbers of usable and broken bikes at each node at
the end of each period. Constraints (7) stipulate that the total number of usable and broken bikes
at each station is no more than its capacity. Constraints (8) specify the initial load of usable
bikes on each truck during the first period. Constraints (9) update the load of usable bikes on each
truck at each node during each period. Constraints (10) require that each truck holds no usable
bikes at the end of the last period. Constraints (11) ensure that the quantity of usable bikes
carried by each truck in any period does not exceed its capacity. Similar constraints for broken
bikes are provided in constraints (12), (13), (14), (15). Constraints (16) ensure that the total
number of usable and broken bikes on each truck during each period is no more than its capacity.

\protect\phantomsection\label{p0230}
Constraints (17), (18) make sure that the total quantity of usable bikes delivered to or collected
from a visited station in a period is less than or equal to the minimum of the station capacity and
the truck capacity. Constraints (19) limit that the number of broken bikes that can be deposited at
the depot in a period is less than or equal to the truck capacity. Constraints (20) ensure no broken
bikes are delivered to any station. Constraints (21) enforce that no broken bikes are collected from
the depot. Constraints (22) make sure that the total quantity of broken bikes collected from a
visited station in a period is less than or equal to the minimum of the initial broken bike
inventory at the station and the truck capacity. Constraints (23) guarantee that no broken bike is
repaired by the on-site repairers at the depot. Constraints (24) limit that the number of broken
bikes repaired at any station is less than or equal to the initial broken bike inventory at the
station. Constraints (25) ensure that the total number of broken bikes collected by trucks and
repaired by repairers at each station does not exceed the initial broken bike inventory.

\protect\phantomsection\label{p0235}
Constraints (26), (27), (28), (29) require that both trucks and repairers depart from the depot
initially and return to the depot eventually. Constraints (30), (31) are the
flow~conservation~conditions for trucks and repairers, respectively. Constraints (32), (33), (34),
(35) relate \(\varsigma_{t,k}^{\text{v}}\)(the variable used to define whether a truck or a repairer stays
still at the depot from period \emph{t} to the end of the rebalancing process) to
\(\partial_{k}^{\text{v}}\) (the period of the truck's or repairer's final return to the depot).

\protect\phantomsection\label{p0240}
Constraints (36), (37), (38), (39) restrict the upper and lower bounds of the operational time, and
are modified from the study of Liu et al. (2018). Specifically, constraints (36) imply that up to
each period, the sum of route travel time and the service time at each node cannot exceed the time
from the start of repositioning to the end of that period, and constraints (37) imply that up to
each period, the sum of route travel time and the service time at each node cannot be lower than the
time from the start of repositioning to the end of two periods before that period. Similar
constraints are also imposed on the repairers in constraints (38), (39). Note that constraints (37),
(39) are only activated before the period that the truck and the repairer return to the depot for
the last time, respectively.

\protect\phantomsection\label{p0245}
Constraints (40), (41) ensure the truck or the repairer stays at the depot after the final return.
Constraints (42), (43) limit that before the final return to the depot, if the truck or repairer
stays at the same node for another period, it must conduct at least one loading, unloading, or bike
repair here. Constraints (44) enforce that a truck can travel between at most one pair of nodes
within a period. A similar set of constraints for repairers is given in constraints (45).
Constraints (46) ensure that each station can be visited at most once by at most one repairer in the
whole repositioning horizon. Constraints (47) define the carbon emissions of each truck produced in
each period when traveling between any two nodes. Constraints (48), (49), (50), (51), (52), (53),
(54), (55), (56), (57), (58), (59), (60), (61), (62), (63), (64), (65) are the domain constraints.

\protect\phantomsection\label{p0250}
Following Kaspi (2015), the EUDF is linearized by introducing a new variable \({\widehat{F}}_{i}\)
for each station \(i \in S\), and the following linear constraints are added to the
formulation.\[\tag{66}{\hat{F}}_{i} \geq \alpha_{p_{i},b_{i}}^{i} \cdot p_{i,T^{'}} + \beta_{p_{i},b_{i}}^{i} \cdot b_{i,T^{'}} + \gamma_{p_{i},b_{i}}^{i},\forall\left( {p_{i},b_{i}} \right) \in \left\{ {\left( {p_{i},b_{i}} \right) \in {\mathbb{Z}}^{2} \mid p_{i} \geq 0,b_{i} \geq 0,p_{i} + b_{i} \leq C_{i}} \right\},\forall i \in S.\]

\protect\phantomsection\label{p0255}
The resultant objective is

\[\tag{67}\begin{array}{l}
{\text{minimize}\mspace{6mu} c^{\text{p}}\sum\limits_{i \in S}{\hat{F}}_{i} + c^{\text{e}}\sum\limits_{i \in S_{0}}\sum\limits_{j \in S_{0}}\sum\limits_{t = t_{i,j}^{'} + 1}^{T^{'}}\sum\limits_{k \in K}e_{i,j,t - t_{i,j}^{'},k} +} \\
{\theta \cdot \left( \begin{array}{l}
{\sum_{i \in S_{0}}\sum_{j \in S_{0}}\sum\limits_{t = t_{i,j}^{'} + 1}^{T^{'}}\sum_{k \in K}\left( {t_{i,j}^{\text{v}} \cdot x_{i,j,t - t_{i,j}^{'},k}^{\text{v}}} \right) + \sum_{i \in S_{0}}\sum_{t = t_{i,j}^{'} + 1}^{T^{'}}\sum_{k \in K}t^{\text{load}} \cdot \left( {y_{i,t,k}^{\text{u}^{+}} + y_{i,t,k}^{\text{u}^{-}} + y_{i,t,k}^{\text{b}^{+}} + y_{i,t,k}^{\text{b}^{-}}} \right)} \\
{\sum_{i \in S_{0}}\sum_{j \in S_{0}}\sum\limits_{t = t_{i,j}^{''} + 1}^{T^{'}}\sum_{m \in R}\left( {t_{i,j}^{\text{r}} \cdot x_{i,j,t - t_{i,j}^{''},m}^{\text{r}}} \right) + \sum_{i \in S_{0}}\sum_{t = t_{i,j}^{''} + 1}^{T^{'}}\sum_{m \in R}t^{\text{repair}} \cdot g_{i,t,m}}
\end{array} \right)}
\end{array}.\]

\protect\phantomsection\label{s0045}
\subsection{4. The solution algorithm}\label{st065}

\protect\phantomsection\label{p0260}
This study developed a hybrid algorithm (HGSADC-SBC) consisting of Hybrid Genetic Search with
Adaptive Diversity Control (HGSADC) and a Station Budget Constrained (SBC) heuristic to solve the
proposed model. The HGSADC is the backbone of the algorithm, and the SBC heuristic is encapsulated
to determine the quantities of loading, unloading, and repairing of bikes for a given route.

\protect\phantomsection\label{s0050}
\subsubsection{4.1. Solution representation}\label{st070}

\protect\phantomsection\label{p0265}
Each solution is represented by
\(\phi = \left( {S_{\text{R}},S_{\text{K}},\Omega_{\text{R}},\Omega_{\text{K}}} \right)\), where
\(S_{\text{R}} = \left( {r_{1},r_{2},...,r_{|R|}} \right)\) denotes the list of repairer routes
\(r_{i}\mspace{100mu}\mspace{100mu}\left( {i \in R} \right)\), and
\(S_{\text{K}} = \left( {r_{1}',r_{2}',...,r_{|K|}'} \right)\) denotes the list of truck routes
\(r_{k}'\mspace{100mu}\left( {k \in K} \right)\). The repairing quantities are represented by
\(\Omega_{\text{R}} = \left( {G_{1},G_{2},\ldots,G_{|R|}} \right)\), with \emph{G\textsubscript{i}}
\(\left( {i \in R} \right)\) specifying the repairing quantities associated with the repairer route
\(r_{i}\mspace{100mu}\mspace{100mu}\left( {i \in R} \right)\). Similarly, the (un)loading quantities
are represented by \(\Omega_{\text{K}} = \left( {Y_{1},Y_{2},\ldots,Y_{|K|}} \right)\), where
\emph{Y\textsubscript{k}} \(\left( {k \in K} \right)\) specifies the (un)loading quantities
associated with the truck route \(r_{k}'\mspace{100mu}\left( {k \in K} \right)\).

\protect\phantomsection\label{p0270}
The route of each repairer is represented as
\(r_{i} = \left( {s_{i_{0}},s_{i_{1}},s_{i_{2}},...,s_{i_{|r_{i}{| - 2}}},s_{i_{|r_{i}{| - 1}}}} \right),i \in R\),
and the route of each truck is represented as
\(r_{k}' = \left( {s_{k_{0}}',s_{k_{1}}',s_{k_{2}}',...,s_{k_{{|r_{k}'|} - 2}}',s_{k_{{|r_{k}'|} - 1}}'} \right),k \in K,\)
where \(s_{i_{j}} \in S_{0}\) and \(s_{k_{l}}' \in S_{0}\)
\(\left( {0 \leq j \leq \left| r_{i} \right| - 1,0 \leq l \leq \left| r_{k}' \right| - 1} \right)\),
and the subscripts \emph{j} and \emph{l} are used to define the order of nodes being visited. Every
route comprises a permutation of visited nodes, with the first and last being the depot, labeled
with 0, i.e.,
\(s_{i_{0}} = s_{i_{|r_{i}{| - 1}}} = 0\mspace{100mu}\mspace{100mu}\mspace{100mu}\left( {i \in R} \right)\)
and
\(s_{k_{0}}' = s_{k_{|r_{k}'{| - 1}}}' = 0\mspace{100mu}\mspace{100mu}\left( {k \in K} \right)\).
For any of the two routes in \emph{S}\textsubscript{R}, no common stations are allowed between the
initial and ending depot.

\protect\phantomsection\label{p0275}
For repairer \emph{i} \(\left( {i \in R} \right)\), the repairing quantities associated with route
\(r_{i}\) are represented as
\(G_{i} = \left( {g_{s_{i_{0}},i},g_{s_{i_{1}},i},\ldots,g_{s_{i_{|r_{i}{| - 1}}},i}} \right)\),
where \(g_{s_{i_{j}},i}\) denotes the number of broken bikes repaired at node \(s_{i_{j}}\) in
\(r_{i}\) \(\left( {0 \leq j \leq \left| r_{i} \right| - 1} \right)\). Similarly, for truck \emph{k}
\(\left( {k \in K} \right)\), the (un)loading quantities associated with route \(r_{k}'\) are
denoted by
\(Y_{k} = \left( {\left( {y_{0,k}^{\text{u}^{+}},y_{0,k}^{\text{u}^{-}},y_{0,k}^{\text{b}^{+}},y_{0,k}^{\text{b}^{-}}} \right),\left( {y_{1,k}^{\text{u}^{+}},y_{1,k}^{\text{u}^{-}},y_{1,k}^{\text{b}^{+}},y_{1,k}^{\text{b}^{-}}} \right),\ldots,\left( {y_{{|r_{k}'|} - 1,k}^{\text{u}^{+}},y_{{|r_{k}'|} - 1,k}^{\text{u}^{-}},y_{{|r_{k}'|} - 1,k}^{\text{b}^{+}},y_{{|r_{k}'|} - 1,k}^{\text{b}^{-}}} \right)} \right)\),
where \(y_{l,k}^{\text{u}^{+}}\) and \(y_{l,k}^{\text{u}^{-}}\) denote the number of usable bikes
unloaded and loaded at node \(s_{k_{l}}\) respectively, and \(y_{l,k}^{\text{b}^{+}}\) and
\(y_{l,k}^{\text{b}^{-}}\) represent the number of broken bikes unloaded and loaded at \(s_{k_{l}}\)
\(\left( {0 \leq l \leq \left| r_{k}' \right| - 1} \right)\) respectively.

\protect\phantomsection\label{p0280}
The solution that corresponds to the one shown in Fig. 3.1 in a network with ten stations and one
depot with two repairers and two trucks can be represented
as\[\phi = \left( {S_{\text{R}},S_{\text{K}},\Omega_{\text{R}},\Omega_{\text{K}}} \right) = \left( {\left( {r_{1},r_{2}} \right),\left( {r_{1}',r_{2}'} \right),\left( {G_{1},G_{2}} \right),\left( {Y_{1},Y_{2}} \right)} \right)\]where:

\begin{itemize}
\item
  \protect\phantomsection\label{p0285}
  \emph{r}\textsubscript{1} = (0, 6, 8, 0) and \emph{G}\textsubscript{1} = (0, 4, 12, 0);
\item
  \protect\phantomsection\label{p0290}
  \emph{r}\textsubscript{2} = (0, 1, 3, 0) and \emph{G}\textsubscript{2} = (0, 3, 2, 0);
\item
  \protect\phantomsection\label{p0295}
  \emph{r′}\textsubscript{1} = (0, 4, 6, 9, 0, 2, 0) and \emph{Y}\textsubscript{1} = ((0, 0, 0, 0),
  (0, 2, 0, 1), (0, 5, 0, 17), (7, 0, 0, 0), (0, 3, 18, 0), (3, 0, 0, 9), (0, 0, 9, 0);
\item
  \protect\phantomsection\label{p0300}
  \emph{r′}\textsubscript{2} = (0, 5, 4, 0, 7, 5, 9, 0) and \emph{Y}\textsubscript{2} = ((0, 0, 0,
  0), (0, 12, 0, 10), (0, 0, 0, 3), (8, 0, 13, 0), (0, 3, 0, 2), (0, 1, 0, 3), (8, 0, 0, 0), (0, 0,
  5, 0)).
\end{itemize}

\protect\phantomsection\label{p0305}
Note that the zeros at the beginning and end of each route are included solely for representation
purposes. In the implementation, these zeros are excluded during operations at each step, with only
the sequence of the visited nodes in between being retained to enhance computational efficiency.

\protect\phantomsection\label{s0055}
\subsubsection{4.2. Overall algorithmic procedure}\label{st075}

\protect\phantomsection\label{p0310}
The algorithmic procedure for HGSADC-SBC is given in \textbf{Algorithm 1}. The algorithm starts with
the initialization of a population \emph{Pop} consisting of a feasible subpopulation
\emph{Pop}\textsubscript{f} and an infeasible subpopulation \emph{Pop}\textsubscript{i} (Line 1).
Then, the population of individuals iteratively evolves through successive operations (Line 3 to
Line 29). The best fitness value and its corresponding solution up to each iteration are kept with
\emph{current\_best\_value} and \emph{current\_best\_sol}, respectively. The algorithm terminates
and returns the best known solution and fitness value when the number of iterations without
improvement (denoted as \emph{iter\_no\_imp}) reaches the limit \emph{It}\textsubscript{NI}, or the
running time (denoted as \emph{cpu\_time}) reaches the limit \emph{T}\textsubscript{max}. A
diversification process is activated every \emph{It}\textsubscript{div} iterations without
improvement to maintain a high level of diversity among individuals (Line 26 to Line 28).

\protect\phantomsection\label{t0100}
{\def\LTcaptype{none} % do not increment counter
\begin{longtable}[]{@{}lll@{}}
\toprule\noalign{}
\multicolumn{3}{@{}l@{}}{%
\textbf{Algorithm 1: HGSADC-SBC}} \\
\midrule\noalign{}
\endhead
\bottomrule\noalign{}
\endlastfoot
\textbf{1} & & Initialization: Generate a population \emph{Pop} consisting of a feasible
subpopulation \emph{Pop}\textsubscript{f} and an infeasible subpopulation
\emph{Pop}\textsubscript{i} \\
\textbf{2} & & Set \emph{iter\_no\_imp~=}~0 and \emph{current\_best\_value} = +\ensuremath{\infty} \\
\textbf{3} & & \textbf{while} \emph{iter\_no\_imp}~\textless~\emph{It}\textsubscript{NI} and
\emph{cpu\_time~\textless~T}\textsubscript{max} \textbf{do} \\
\textbf{4} & & Randomly select two individuals \emph{Ind}\textsubscript{1} and
\emph{Ind}\textsubscript{2} from \emph{Pop} \\
\textbf{5} & & Generate the routes of offspring \emph{C}\textsubscript{1} and
\emph{C}\textsubscript{2} using ordered crossover on the routes of \emph{Ind}\textsubscript{1} and
\emph{Ind}\textsubscript{2}, then apply the SBC heuristic on the routes to determine the (un)loading
and repairing quantities for \emph{C}\textsubscript{1} and \emph{C}\textsubscript{2} \\
\textbf{6} & & Educate on both \emph{C}\textsubscript{1} and \emph{C}\textsubscript{2} by applying
local search procedures to generate new routes and using the SBC heuristic to determine the
(un)loading and repairing quantities \\
\textbf{7} & & \textbf{if} \emph{C\textsubscript{i}} (\emph{i}~=~1, 2) is feasible \textbf{then} \\
\textbf{8} & & ~Add \emph{C\textsubscript{i}} to the feasible subpopulation
\emph{Pop}\textsubscript{f} \\
\textbf{9} & & \textbf{else} \\
\textbf{10} & & ~Add \emph{C\textsubscript{i}} to the infeasible subpopulation
\emph{Pop}\textsubscript{i} and repair \emph{C\textsubscript{i}} with the probability
\emph{P}\textsubscript{repair} using local search procedures and the SBC heuristic \\
\textbf{11} & & \textbf{end if} \\
\textbf{12} & & \textbf{if} the individuals become feasible after repairs \textbf{then} \\
\textbf{13} & & ~The repaired individuals are added to the feasible subpopulation
\emph{Pop}\textsubscript{f} \\
\textbf{14} & & \textbf{else} \\
\textbf{15} & & ~The repaired individuals are discarded, and the original infeasible solution is
kept in the infeasible subpopulation \emph{Pop}\textsubscript{i} \\
\textbf{16} & & \textbf{end if} \\
\textbf{17} & & \textbf{if} the size of \emph{Pop}\textsubscript{f} or \emph{Pop}\textsubscript{i}
reaches the maximum subpopulation size \textbf{then} \\
\textbf{18} & & ~Select survivors from the corresponding subpopulation \\
\textbf{19} & & \textbf{end if} \\
\textbf{20} & & Set \emph{bs}~=~the solution with the best fitness value from
\emph{Pop}\textsubscript{f} and set \emph{bv}~=~the fitness value of \emph{bs} \\
\textbf{21} & & \textbf{if} \emph{bv} \(\geq\) \emph{current\_best\_value} \textbf{then} \\
\textbf{22} & & ~Set \emph{iter\_no\_imp}~=~\emph{iter\_no\_imp}~+~1 \\
\textbf{23} & & \textbf{Else} \\
\textbf{24} & & ~Set \emph{current\_best\_value}~=~\emph{bv}, \emph{current\_best\_sol}~=~\emph{bs},
and \emph{iter\_no\_imp}~=~0 \\
\textbf{25} & & \textbf{end if} \\
\textbf{26} & & \textbf{if} \emph{iter\_no\_imp} \(\geq It_{\text{div}}\) \textbf{then} \\
\textbf{27} & & ~Diversify the population consisting of \emph{Pop}\textsubscript{f} and
\emph{Pop}\textsubscript{i} and adjust the penalty parameter \\
\textbf{28} & & \textbf{end if} \\
\textbf{29} & & \textbf{end while} \\
\textbf{30} & & \textbf{return} \emph{current\_best\_sol}, \emph{current\_best\_value} \\
\end{longtable}
}

The algorithm structure follows the one given by Vidal et al., 2012, Vidal et al., 2013, Vidal et
al., 2014). The route representation, crossover operator, and education operator are based on the
study by Li et al. (2016) with some modifications, mostly on extending their methods to address the
scenario of multiple routes for multiple trucks and repairers from the scenario of a single route
for only one truck. Moreover, in this algorithm, the crossover and local search operations in
education and the repair process are applied to the routes of the solutions. Once new routes are
generated through these operations, the SBC heuristic is subsequently applied to determine the
loading, unloading, and repair quantities, completing each solution with both routes and associated
quantities. The SBC heuristic is developed based on a novel station budget concept, and is detailed
in Section 4.4.

\protect\phantomsection\label{s0060}
\subsubsection{4.3. Population initialization}\label{st080}

\protect\phantomsection\label{p0315}
To initialize a population, we generate solutions one by one until a specified population size is
reached. The generation of each solution requires initializing routes for both repairers and trucks,
followed by applying the SBC heuristic to determine the associated (un)loading and repairing
quantities. Each solution then undergoes a feasibility check (as detailed in Section 4.6) and is
added to the feasible or infeasible subpopulation based on its feasibility.

\protect\phantomsection\label{p0320}
To initialize a route, the depot is set as the first node. Then, new nodes are iteratively added to
the route based on the last node of the current route, denoted as \emph{n}\textsubscript{last}.
Specifically, we first find the set of nodes that allow the truck/repairer to reach there from
\emph{n}\textsubscript{last}, conduct at least one loading/unloading/repairing there, and return to
the depot within the time budget. The set is defined as the \emph{reachable list} of
\emph{n}\textsubscript{last}. Note that \emph{n}\textsubscript{last} itself is not included in the
set. We randomly choose one node from the \emph{reachable list} and append it to the end of the
route. The procedure repeats until the reachable list of \emph{n}\textsubscript{last} is empty, and
we finally append the depot at the end. Other routes are generated in the same manner. Note that the
routes are generated one by one, in the order of repairers first and trucks second to make the best
use of the repairers to repair broken bikes before truck arrivals. It should also be noted that any
station that has already been visited by a repairer is not considered when generating routes for
other repairers.

\protect\phantomsection\label{s0065}
\subsubsection{4.4. Station budget constrained (SBC) heuristic for quantity
determination}\label{st085}

\protect\phantomsection\label{p0325}
For each route of a truck or repairer, we have to determine the quantities of bikes loaded,
unloaded, and repaired at designated stations to form complete solutions and calculate their fitness
values. To determine these quantities, an SBC heuristic is proposed. Some notations and
terminologies are defined first to better describe the SBC heuristic:

\begin{itemize}
\item
  \protect\phantomsection\label{p0330}
  \(q_{s}\) represents the target usable inventory level at station \emph{s}
  \(\left( {s \in S} \right)\), which is determined as the usable bike inventory at station \emph{s}
  that corresponds to the lowest EUDF value. It together with the initial inventory of usable bikes
  \(p_{s}^{0}\) can be used to define the station type. A station is categorized into a deficit
  station, a surplus station, and a balanced station if \(p_{s}^{0} < q_{s}\),
  \(p_{s}^{0} > q_{s}\), and \(p_{s}^{0} = q_{s}\) hold, respectively.
\item
  \protect\phantomsection\label{p0335}
  \(p_{s}\) and \(b_{s}\) are used to track the inventory of usable and broken bikes at station
  \emph{s} (\(s \in S\)), respectively. At the beginning of the SBC heuristic, they are assigned to
  the initial inventories before the repositioning starts. During the procedure of the heuristic,
  they are updated when bikes are loaded or unloaded from the trucks or when repairers conduct
  repairs.
\item
  \protect\phantomsection\label{p0340}
  \(\psi_{l,k}^{\text{u}}\) and \(\psi_{l,k}^{\text{b}}\) are used to track the inventory of usable
  and broken bikes on truck \emph{k} (\(k \in K\)) when it leaves the \emph{l}-th visited node
  \(s_{k_{l}}'\left( {l = 0,...,\left| r_{k}' \right| - 1} \right)\).
\item
  \protect\phantomsection\label{p0345}
  \(\mathit{maxTim}e_{s,i}^{\text{r}}\) represents the station repairing budget assigned to station
  \(s \in S\) for repairer \(i \in R\).
\item
  \protect\phantomsection\label{p0350}
  \(\mathit{maxTim}e_{s,k}^{\text{t}}\) represents the station (un)loading budget assigned to
  station \(s \in S\) for truck \(k \in K\).
\item
  \protect\phantomsection\label{p0355}
  \(\mathit{maxTimeRemai}n_{s,k}^{\text{t}}\) represents the remaining station (un)loading budget
  assigned to station \(s \in S\) for truck \(k \in K\).
\item
  \protect\phantomsection\label{p0360}
  \(\mathit{unsa}t_{s,k}^{\text{load,u}}\) represents the number of usable bikes not loaded onto
  truck \emph{k} due to the limited station loading budget at station \(s \in S\).
\item
  \protect\phantomsection\label{p0365}
  \(\mathit{unsa}t_{s,k}^{\text{unload,u}}\) represents the number of usable bikes not unloaded to
  station \(s \in S\) from truck \emph{k} due to the limited station unloading budget.
\item
  \protect\phantomsection\label{p0370}
  \(\mathit{unsa}t_{s,k}^{\text{load},\text{b}}\) represents the number of broken bikes not loaded
  onto truck \emph{k} due to the limited station loading budget at station \(s \in S\).
\item
  \protect\phantomsection\label{p0375}
  \(\mathit{maxTimeIni}t_{s,k}^{\text{t}}\) represents the initial station time budget assigned to
  station \(s \in S\) for truck \emph{k} before determining truck loading quantities.
\item
  \protect\phantomsection\label{p0380}
  \(\mathit{maxTimeNe}w_{s,k}^{\text{t}}\) represents the updated station time budget assigned to
  station \(s \in S\) for truck \emph{k} after adjustment.
\end{itemize}

\protect\phantomsection\label{p0385}
We point out that the variables presented above are global variables that are shared across all
subsequent algorithms, which means any changes to these variables are effective in all algorithms.
For the sake of brevity, we directly use and assign values to these variables within the algorithms
without redefining them or including them as arguments in each algorithmic procedure. In addition,
in the following algorithm description, we use a notation where a variable is enclosed in
parentheses and subscripted by its index domains to represent a list of variables indexed by the
subscripts. For example, \(\left( {\mathit{maxTim}e_{s,i}^{\text{r}}} \right)_{s \in S}\) represents
a one-dimensional list that stores the station budget at each station in route \(r_{i}\) for
repairer \emph{i.}

\protect\phantomsection\label{p0390}
The SBC heuristic procedure is formed by the following steps:

\begin{itemize}
\item
  \protect\phantomsection\label{p0395}
  \textbf{Step 1: Initialization}. Set initial values for station inventories.
\item
  \protect\phantomsection\label{p0400}
  \textbf{Step 2: Assign repairing quantities for repairers.} For each repairer, conduct the
  following three subprocedures:
\item
  \protect\phantomsection\label{p0405}
  Assign a station repair budget to each visited station in the repairer's route.
\item
  \protect\phantomsection\label{p0410}
  Determine the number of bike repairs to be conducted at each station, constrained by the
  quantities allowed to be handled under the assigned station budget, the broken bike inventory at
  the station, and the deviation of the usable bike inventory to the station target.
\item
  \protect\phantomsection\label{p0415}
  Update the station\textquotesingle s usable and broken bike inventories to reflect the completed
  repairs.
\item
  \protect\phantomsection\label{p0420}
  \textbf{Step 3: First-round truck loading and unloading assignment.} For each truck, conduct the
  following three subprocedures:
\item
  \protect\phantomsection\label{p0425}
  Allocate a station (un)loading budget for each station in the truck's route.
\item
  \protect\phantomsection\label{p0430}
  Assign the initial quantities for loading and unloading usable and broken bikes at each station,
  based on station status (surplus or deficit) and station budget.
\item
  \protect\phantomsection\label{p0435}
  Update truck and station inventories accordingly.
\item
  \protect\phantomsection\label{p0440}
  \textbf{Step 4: Adjust station budgets based on unused time.} Adjust the station (un)loading
  budgets to redistribute any unused time to stations where demand remains unmet.
\item
  \protect\phantomsection\label{p0445}
  \textbf{Step 5: Second-round truck loading and unloading assignment.} Revert station inventories
  to their values before the first round of loading/unloading assignment, and reassign the
  quantities for the second round using the adjusted station budgets obtained from \textbf{Step 4}.
\item
  \protect\phantomsection\label{p0450}
  \textbf{Step 6: Final amendments to routes and quantities.} Make final adjustments to the truck
  and repairer routes, along with their associated loading, unloading, and repairing quantities, to
  optimize the use of any time budget left.
\end{itemize}

\protect\phantomsection\label{p0455}
For the detailed procedure of the SBC heuristic, readers can refer to \textbf{Algorithm 2 in}
\textbf{Appendix B}.

\protect\phantomsection\label{p0460}
In the following subsections, we delve into the details of each component of the algorithm,
including station budget assignment in Steps 2.1 and 3.1 (Section 4.4.1), repairer repairing
quantity assignment in Step 2.2 and truck (un)loading quantity assignment in Step 3.2 and (Section
4.4.2), station (un)loading budget and truck (un)loading quantity adjustments in Steps 4 and 5
respectively (Section 4.4.3), and the final amendments to the routes, repairing, and (un)loading
quantities in Step 6 (Section 4.4.4).

\protect\phantomsection\label{s0070}
\paragraph{4.4.1. Station budget allocation}\label{st090}

\protect\phantomsection\label{p0465}
The SBC heuristic allocates truck (un)loading time and repairer repairing time at stations by
leveraging the EUDF to define a Benefit-Cost Ratio Function (BCRF). The EUDF captures user
dissatisfaction as a function of station inventories, offering a convex framework for assessing the
impact of inventory adjustments. The BCRF builds upon the EUDF by quantifying the effectiveness of
resource allocation at each station to reduce user dissatisfaction per unit of time spent.
Specifically, the BCRF for a station is computed as the change in EUDF value---reflecting the
benefit of truck (un)loading or repairs---divided by the operational time required to achieve that
change.

\protect\phantomsection\label{p0470}
Given the usable bike inventory \emph{p}, the broken bike inventory \emph{b}, and the target usable
bike inventory \emph{q} at station \emph{s}, the BCRF value for a truck at that station (denoted as
\(\mathit{BCR}F_{s}^{\text{t}}\left( {p,b} \right)\)) is defined
as\[\tag{68}\mathit{BCR}F_{s}^{\text{t}}\left( {p,b} \right) = \left\{ \begin{array}{lc}
{0,} & {p = q,b = 0;} \\
{\frac{F_{s}\left( {p,b} \right) - F_{s}\left( {q,0} \right)}{t^{\text{load}} \cdot \left( \left| q - p \middle| + b \right. \right)},} & {\text{Otherwise}\text{.}}
\end{array} \right.\]

\protect\phantomsection\label{p0475}
For balanced stations without broken bikes, the BCRF value for a truck at that station is defined as
zero, as indicated in the first condition of the piecewise function. In the second condition, the
denominator is the time required to handle the usable and broken bikes at the station to achieve the
target inventory and is considered to be a cost, assuming the value of time is fixed. The numerator
is the penalty reduction and is considered to be a benefit. \(F_{s}\left( {q,0} \right)\) is the
minimum EUDF value so the numerator is always non-negative.

\protect\phantomsection\label{p0480}
Fig. 4.1 shows an example illustrating the calculation of the values of the BCRF for trucks based on
the EUDF. The three columns in this illustration represent Stations 1, 2, and 3, respectively.

\begin{figure}
\centering
Fig. 4.1. An illustration of calculating the BCRF values of three stations based on the EUDF values.
\caption{}\label{f0020}
\end{figure}

\protect\phantomsection\label{p0485}
In the first row, the bar charts display the current inventory of usable bikes, the target inventory
for usable bikes, and the current inventory of broken bikes at each station. Arrows indicate the
deviations of the current inventories of usable bikes from the corresponding target inventory
levels. The second row contains heat maps showing the EUDF values for various combinations of usable
and broken bike inventories at each station. Darker colors correspond to higher dissatisfaction, as
captured by the EUDF. The red dots represent the current inventory levels, and the green dots
indicate the target inventory levels. The arrows between these points show how the current inventory
can be adjusted to match the target, specifying how many usable and broken bikes need to be loaded
and unloaded.

\protect\phantomsection\label{p0490}
For Station 3, the change in color between the current and target inventory levels is subtle,
suggesting that spending time at this station has little effect on reducing the EUDF value. In
contrast, Station 2 shows a significant color shift, indicating a large potential for reducing
dissatisfaction by adjusting the inventory level.

\protect\phantomsection\label{p0495}
The observations in the last paragraph are further quantified in the BCRF values on the last row,
which are calculated by dividing the reduction in EUDF value by the time required to adjust the
inventory from the current to the target level. Station 2′s BCRF value is nearly ten times greater
than that of Station 3, emphasizing that more time should be spent on Station 2 to maximize the
reduction in user dissatisfaction per unit of time.

\protect\phantomsection\label{p0500}
Similarly, for each repairer, we also have \(\mathit{BCR}F_{s}^{\text{r}}\left( {p,b} \right)\)
defined as\[\tag{69}\mathit{BCR}F_{s}^{\text{r}}\left( {p,b} \right) = \left\{ \begin{array}{lc}
{0,} & {p \geq q,b = 0;} \\
{\frac{F_{s}\left( {p,b} \right) - F_{s}\left( {p + \min\left\{ {b,\max\left\{ {q - p,0} \right\}} \right\},0} \right)}{t^{\text{repair}} \cdot \min\left\{ {b,\max\left\{ {q - p,0} \right\}} \right\}},} & {\text{Otherwise}\text{.}}
\end{array} \right.\]

\protect\phantomsection\label{p0505}
For balanced and surplus stations without broken bikes, the BCRF value for a truck at that station
is defined as zero, as indicated in the first condition of the piecewise function. In the second
condition, the denominator \(\min\left\{ {b,\max\left\{ {q - p,0} \right\}} \right\}\) is the
maximum number of broken bikes repaired at the concerned station, which cannot be larger than the
number of usable bikes required at that station because the repaired bikes become usable. That is
why the repaired bikes are also included in the second term of the numerator.

\protect\phantomsection\label{p0510}
The station repair budget is allocated based on the calculated \emph{BCRF}\textsuperscript{r} for
each station, which is determined by prioritizing those with higher BCRF values that reflect a
greater reduction in user dissatisfaction per unit of time. Specifically, the available repair time
is distributed proportionally across stations according to the \emph{BCRF}\textsuperscript{r} of the
station, ensuring a balanced allocation of resources while preventing excessive time spent at any
single station. To make the best use of the time budget, any overly allocated time is redistributed
to unsatisfied stations with unmet repair needs following the descending order of
\emph{BCRF}\textsuperscript{r}. The detailed algorithm for this is presented in \textbf{Algorithm 3}
in Appendix C. The station (un)loading budget allocation algorithm for trucks is similar to the one
for repairers, as presented in \textbf{Algorithm 4} in Appendix C, except that it considers the
truck's loading and unloading operations instead of repair operations. To illustrate the main idea
with a simple example, consider a route consisting of the above three stations 0--1-2--3-0, where 0
is the depot. After accounting for the route's travel time, the remaining time for bike loading and
unloading is distributed proportionally across the stations based on their BCRF values. The results
of this distribution are depicted in Fig. 4.2, which are calculated as follows:\[\begin{array}{l}
{\text{Station 1:}\frac{0.0103}{0.0103 + 0.0344 + 0.0032} = 21.5\%,} \\
{\text{Station 2:}\frac{0.0344}{0.0103 + 0.0344 + 0.0032} = 71.8\%,\;\text{and}} \\
{\text{Station 3:}\frac{0.0032}{0.0103 + 0.0344 + 0.0032} = 6.7\%.}
\end{array}\]

\begin{figure}
\centering
Fig. 4.2. Proportion of operational time allocated to each station based on their BCRF values.
\caption{}\label{f0025}
\end{figure}

\protect\phantomsection\label{s0075}
\paragraph{4.4.2. Loading, unloading, and repairing quantity assignment}\label{st095}

\protect\phantomsection\label{p0515}
With the station (un)loading budgets determined, we proceed to assign the quantities of usable bikes
to be loaded or unloaded by trucks, broken bikes to be loaded by trucks, and broken bikes to be
repaired by repairers. To improve the clarity of the main text, the detailed algorithms for these
assignments (from \textbf{Algorithm 5} to \textbf{Algorithm 9}) are provided in Appendix D.
Different from the traditional greedy heuristic, which assigns quantities based solely on the
truck's inventory and the station's demand, our heuristic incorporates the station budget as an
additional factor. The budget limits the maximum time a truck or repairer can spend at a station for
operations. By considering this factor, we prevent excessive time consumption at early stations that
may not necessarily justify such a large time investment, and ensure a more balanced and efficient
allocation of resources across all stations.

\protect\phantomsection\label{s0080}
\paragraph{4.4.3. Station (un)loading budget and truck (un)loading quantity
adjustments}\label{st100}

\protect\phantomsection\label{p0520}
In \textbf{Algorithm 2} in Appendix B, after the assignment of loading and unloading quantities, it
is observed that some stations may still have a loading budget unused while others encounter a
shortfall. To make use of the extra time, an adjustment to the station (un)loading budget
allocations for trucks is necessary. The adjustment aims to redistribute the unused station
(un)loading budgets from stations where the allocated time budget exceeds the actual loading and
unloading time requirements to those stations where the budget is insufficient. The algorithm is
given in Appendix E.

\protect\phantomsection\label{s0085}
\paragraph{4.4.4. Amendments to the routes and assigned quantities for trucks and
repairers}\label{st105}

\protect\phantomsection\label{p0525}
When assigning repair, loading, and unloading quantities, a conservative rounding-down strategy is
applied (e.g., Line 3 in Algorithm 6 in Appendix D). This can lead to the availability of unused
time budgets, which can be utilized for more bike repairs and loading. Therefore, this section
provides amendments to the routes and assigned quantities for trucks and repairers.

\protect\phantomsection\label{p0530}
The amendment procedure begins by removing stations that are not assigned with any loading,
unloading, and repairing quantities, then calculating the leftover time for each truck and repairer
based on the outputs from \textbf{Algorithm 2}. It then focuses on the stations which are still not
balanced. For each repairer at every station, the amendment procedure explores the possible broken
bike repairing quantity that ranges from zero to the maximum allowed by the extra time. It selects
the repairing quantity that results in the lowest \(\mathit{BCR}F_{s}^{\text{r}}\). Meanwhile, each
truck checks the surplus stations and redistributes as many surplus bikes as possible to deficit
stations. The choice of the surplus and deficit stations prioritizes those with the highest
\(\mathit{BCR}F_{s}^{\text{t}}\). For each surplus station checked by the truck, the candidate
destinations are set as the stations between the current station and the next depot along the
truck's route, with surplus bikes unloaded to the next depot if unloading at any candidate station
would violate station or truck capacity constraints. After checking all surplus stations, if extra
time remains, supplementary loadings of broken bikes are conducted similarly, with the distinction
that broken bikes can only be unloaded at the depot.

\protect\phantomsection\label{s0090}
\subsubsection{4.5. Feasibility check and solution evaluation}\label{st110}

\protect\phantomsection\label{p0535}
After determining the loading, unloading, and repairing quantities, we can evaluate solution
\(\phi\) by\[\tag{70}Z(\phi) = f(\phi) + \eta\varepsilon(\phi)\]where \(f(\phi)\) is the sum of the
penalty costs of the total user dissatisfaction at all stations in the network and the total cost of
carbon emissions during the rebalancing process. \(\varepsilon(\phi)\) is related to the penalties
for solution infeasibility and \(\eta\) is the penalty coefficient of \(\varepsilon(\phi)\). If
\(\varepsilon(\phi)\) is 0, then solution \(\phi\) is feasible.

\protect\phantomsection\label{p0540}
The infeasibility we consider here originates from the fact that the heuristic proposed in Section
4.4 determines the quantities route by route separately, without considering the interdependencies
between different routes and their effects on station inventories. This strategy is used to make the
decision-making process more straightforward and easier to implement. It assumes that the arrival
times at stations follow the visiting order of stations in the route, with all visits by trucks with
lower indices occurring before any visits by trucks with higher indices. This assumption, however,
can lead to invalid quantities being assigned.

\protect\phantomsection\label{p0545}
Consider a problem where two trucks perform the repositioning task in a network with one depot. For
simplicity, we assume there are no repairers in this example. The route for Truck 1 is
(0--1--3--5--0), and the route for Truck 2 is (0--2--5--0). We notice that Station 5 is visited by
both trucks. It is unclear which truck would arrive first at Station 5 when the arrival times at
stations, which depend on quantity assignment and travel times, are not calculated. The method
proposed in Section 4.4 presumed that the first truck arrived at Station 5 ahead of the second
truck, updating the inventories and assigning the repositioning quantities accordingly, which can
lead to an invalid assignment. For example, as shown in Fig. 4.3 \textbf{(a)}, suppose the capacity
of the station is 20, with a usable bike inventory of 3, a broken bike inventory of 10, and a target
usable bike inventory of 13. If Truck 1 arrives at Station 5 from Station 3 with a truckload of 0
usable bikes and 0 broken bikes, according to the strategy in Section 4.4, it collects 10 broken
bikes (provided the remaining time is enough) at Station 5, and thus the broken bike inventory of
Station 5 is updated to 0. Then, for Truck 2, if it arrives at Station 5 with 12 usable bikes, it
delivers 10 usable bikes to achieve the target usable bike inventory. However, if Truck 1 actually
arrives at Station 5 later than Truck 2, the decision to deliver 10 usable bikes by Truck 2 results
in a total inventory of usable and broken bikes of 23 (3~+~10~+~10), which exceeds the station
capacity (see Fig. 4.3 \textbf{(b)}). The scenario could become more complicated if the repairers
also conduct repairs at the station. Therefore, we need to validate the feasibility of the solution
by examining the inventories of the multiple visited stations after the operation of each visited
truck or repairer.

\begin{figure}
\centering
Fig. 4.3. An example of solution infeasibility in the quantity assignment.
\caption{}\label{f0030}
\end{figure}

\protect\phantomsection\label{p0550}
After the operations at a node, the usable or broken bike inventory may be below 0, or the total
inventory of usable and broken bikes may be over the capacity of the station. We set
\(\varepsilon(\phi)\) to the total number of bikes that exceeds the capacity or is below 0, which is
calculated following the procedure below:

\begin{itemize}
\item
  \protect\phantomsection\label{p0555}
  \textbf{Step 1:} Calculate the arrival time \emph{ar\textsubscript{s}} for every station
  \(s \in S\) visited by each repairer and truck according to solution \(\phi\) and the quantities
  determined in Section 4.4, and record the results as a list of tuples
  \(\left( {n,s,y_{s}^{\text{b}^{-}},y_{s}^{\text{u}^{+}},y_{s}^{\text{u}^{-}},g_{s},ar_{s}} \right)\),
  where \emph{n} represents the index of the repairer or truck operating at station \emph{s}, and
  the number of loaded broken bikes, loaded usable bikes, unloaded usable bikes, and repaired broken
  bikes are represented as \(y_{s}^{\text{b}^{-}},y_{s}^{\text{u}^{-}},y_{s}^{\text{u}^{+}}\), and
  \(g_{s}\), respectively.
\item
  \protect\phantomsection\label{p0560}
  \textbf{Step 2:} Sort the list of tuples according to the arrival time \(ar_{s}\).
\item
  \protect\phantomsection\label{p0565}
  \textbf{Step 3:} Select the nodes in the list one by one in the ascending order of arrival time
  and perform the loading/unloading/repairing to update the inventory of usable and broken bikes.
  Update \(\varepsilon(\phi)\) if, after the operations at some node, the usable bike inventory
  \(p_{s}\) or the broken bike inventory \(b_{s}\) is below 0, or the total inventory of usable and
  broken bikes exceeds the capacity of the station \(C_{s}\) by
\end{itemize}

\[\left. \varepsilon(\phi)\leftarrow\varepsilon(\phi) - \left( {\min\left\{ {p_{s},0} \right\} + \min\left\{ {b_{s},0} \right\} + \min\left\{ {C_{s} - p_{s} - b_{s},0} \right\}} \right). \right.\]

\protect\phantomsection\label{p0570}
Next, we define the biased fitness value of a solution to consider the diversity in each
subpopulation, which is calculated
with\[\tag{71}BF(\phi) = rank_{\text{fit}}(\phi) + \left( {1 - \frac{N_{\text{elite}}}{N_{\text{indiv}}}} \right)rank_{\text{sim}}(\phi)\]where
\emph{N\textsubscript{elite}}~is the number of elite individuals that survive to the next
generation, and~\emph{N\textsubscript{indiv}}~is the actual number of individuals in the
subpopulation. \(\mathit{ran}k_{\text{fit}}(\phi)\) is the rank of solution \(\phi\) in its
subpopulation regarding \(Z(\phi)\) sorted in the ascending order, and
\(\mathit{ran}k_{\text{sim}}(\phi)\) is the rank of solution \(\phi\) regarding the similarity with
other solutions in the subpopulation.

\protect\phantomsection\label{p0575}
In our problem, for each solution \(\phi\), we gather the arc set of all truck routes in \(\phi\) to
the set \(TA_{\phi}\), and the arc set of all repairer routes in \(\phi\) to a set \(RA_{\phi}\).
Let \(\chi\) be the subpopulation that \(\phi\) belongs to. Depending on the feasibility of
\(\phi\), \(\chi = Pop_{\text{f}}\) if \(\phi\) is feasible, and \(\chi = Pop_{\text{i}}\)
otherwise. Then, the similarity of \(\phi\) with respect to the subpopulation \(\chi\) it belongs to
is defined
as\[\tag{72}\left| {TA_{\phi} \cap \left( {\bigcup\limits_{\phi' \in \chi \smallsetminus {\{\phi\}}}{TA_{\phi'}}} \right)} \right| + \left| {RA_{\phi} \cap \left( {\bigcup\limits_{\phi' \in \chi \smallsetminus {\{\phi\}}}{RA_{\phi'}}} \right)} \right|.\]which
is the sum of the number of common arcs between the truck routes in the solution and the union of
truck routes from every other solution in the corresponding subpopulation and the number of common
arcs between repairer routes in the solution and the union of repairer routes from every other
solution in the corresponding subpopulation.

\protect\phantomsection\label{s0095}
\subsubsection{4.6. Selection, crossover, education, and repairing}\label{st115}

\protect\phantomsection\label{p0580}
In each iteration, two individuals are selected using a binary tournament, which randomly picks two
individuals from the population and reserves the one with a better-biased fitness value. The
crossover and local search strategies for routes in education and repair are basically the same as
those proposed by Li et al. (2016). Therefore, we omit the details here. The only difference is that
in our solution, there can be multiple truck routes and multiple repairer routes. Hence, for each
solution, the crossover, education, and repair are carried out route by route. For cases with more
than one truck or repairer, we only consider the crossover of the routes with the same index. For
example, the route of the first repairer in solution \(\phi_{1}\) only makes a crossover with the
route of the first repairer in solution \(\phi_{2}\). The repair operation follows the method by
Vidal (2022) except that the SBC heuristic is used after each new route is generated.

\protect\phantomsection\label{s0100}
\subsubsection{4.7. Population management and search guidance}\label{st120}

\protect\phantomsection\label{p0585}
The population management mechanism works with selection, crossover, and education operators to help
propagate the characteristics of excellent solutions, increase population diversity, and provide
guidance for a more efficient search. At the beginning of the algorithm, \(4\mu\) individuals are
initialized (Li et al., 2016), where \(\mu\) is the minimum subpopulation size. Each of them
directly goes through an education process using local search procedures and the SBC heuristic, and
then falls into a feasible or infeasible subpopulation based on the feasibility of the solution. The
repair operation is activated for infeasible individuals with a probability of
\emph{P}\textsubscript{repair}. The repair operation runs local search procedures and the SBC
heuristic under a penalty factor of 10 times the original one, aiming to convert an infeasible
solution into a feasible one (Vidal, 2022). If the repair is successful, the resulting individual is
added to the feasible subpopulation. Otherwise, the infeasible one before repair is preserved in the
infeasible subpopulation. The algorithm keeps generating new individuals under the process of
selection, crossover, education, and repair until the time limit or the maximum number of iterations
without improvement is reached.

\protect\phantomsection\label{p0590}
During the above process, individuals generated through crossover, education, or repair are added to
the corresponding subpopulation based on their feasibility. If the size of either subpopulation
reaches \(\mu + \lambda\), where \(\lambda\) is the number of offspring in a generation, survivor
selection is conducted to remove the \(\lambda\) solutions with the highest biased fitness values to
ensure that the size equals \(\mu\).

\protect\phantomsection\label{p0595}
As in the study of Li et al. (2016), the penalty coefficient \(\eta\) in Eq. (70) is modified every
100 iterations to adjust the proportion of feasible solutions. Let \(\xi^{\text{REF}}\) be the
target proportion of feasible individuals. If the proportion of feasible individuals with respect to
\(\varepsilon(\phi)\) is less than \(\xi^{\text{REF}} - 5\%\) or greater than
\(\xi^{\text{REF}} + 5\%\), then we change \(\eta\) to \(1.2\eta\) or \(0.85\eta\), respectively (Li
et al., 2016).

\protect\phantomsection\label{p0600}
The diversification operation is triggered if the best solution has not been improved for
\emph{It}\textsubscript{div} iterations. In the diversification, the best \(\frac{\mu}{3}\)
solutions of each subpopulation are preserved (Li et al., 2016), and \(4\mu\) new individuals are
generated, followed by the survivor selection to adjust the size of each subpopulation to \(\mu\).

\protect\phantomsection\label{s0105}
\subsection{5. Numerical studies}\label{st125}

\protect\phantomsection\label{p0605}
The proposed model was implemented in Python 3.11 using Gurobi Python API (Gurobi Optimization, LLC,
2023). The hybrid algorithm HGSDAC-SBC was implemented with C++. All numerical experiments were
conducted on a computer equipped with an Intel Core i9-12900 processor (12 cores, 24 threads) and
64~GB of physical memory.

\protect\phantomsection\label{p0610}
In Section 5.1, we introduce the parameter tuning and setting for HGSDAC-SBC. In Section 5.2, we
demonstrate the effectiveness of the proposed algorithm by presenting results and running times for
small instances involving a single truck and repairer. In Section 5.3, we evaluate the algorithm's
performance on larger networks with multiple trucks and repairers. Section 5.4 focuses on the impact
of varying repair times on the cost-effectiveness of including a repairer through experiments on
different proportions of broken bikes.

\protect\phantomsection\label{p0615}
The parameters regarding the network are from the dataset used by Ho and Szeto (2017), which is
available at https://www.dropbox.com/s/jbsfefi2wqsssrr/hlns-set3.zip?dl=0. The dataset contains the
travel times between nodes, the capacity of each station, and the initial inventory of usable bikes
at each station. The travel times used in the experiments are also from the dataset above, in which
the asymmetric travel times (in seconds) between the stations were obtained from the Open Source
Routing Machine project (https://project-osrm.org/) on the travel times of trucks between stations.
However, since the project does not provide travel times for electric tricycles (the vehicle assumed
to be used by repairers in this study), we assume the speed of the electric tricycles is
proportional to that of the trucks. The proportionality constant is determined based on values from
the existing literature that provide relevant speed information for both vehicles (see Table 5.1).
As a result, the travel times for repairers between nodes can be calculated accordingly based on the
constant.

\protect\phantomsection\label{t0010}
Table 5.1. Parameter setting of the numerical experiments.

{\def\LTcaptype{none} % do not increment counter
\begin{longtable}[]{@{}llll@{}}
\toprule\noalign{}
\multicolumn{2}{@{}l}{%
Parameter} & Value & References \\
\midrule\noalign{}
\endhead
\bottomrule\noalign{}
\endlastfoot
\emph{ρ}\textsubscript{0} & The empty-load fuel consumption rate of each truck (liters/km) & 0.252 &
Franceschetti et al., 2013, Shui and Szeto, 2018 \\
\emph{ρ*} & The full-load fuel consumption rate of each truck (liters/km) & 0.258 & Franceschetti et
al., 2013, Shui and Szeto, 2018 \\
\emph{CER} & The CO\textsubscript{2} emission rate (kg/liter) & 2.61 & Wang and Szeto (2018) \\
\(Q_{k}\) & The capacity of truck \emph{k} (bikes) & 25 & Ho and Szeto (2017) \\
\(t^{\text{load}}\) & Time for loading/unloading a bike (seconds) & 60 & Liu et al. (2018) \\
\(t^{\text{repair}}\) & Time for repairing a bike on-site (seconds) & 300 & Jin et al. (2022) \\
\(\frac{v^{\text{t}}}{v^{\text{m}}}\) & The ratio of speed between a truck and an electric tricycle
& 1.68 & Wu, 2021, Fang et al., 2019 \\
\(c^{\text{p}}\) & The penalty cost for a user unable to rent/return a bike (CNY) & 2 & Chen et al.
(2020) \\
\(c^{\text{e}}\) & The unit carbon trading cost (CNY/kg) & 0.06 & Y. Jia et al. (2020) \\
\(\theta\) & The penalty coefficient used to eliminate unnecessary loading or unloading operations &
1e-8 & − \\
\end{longtable}
}

\protect\phantomsection\label{p0620}
The EUDF value of each station was calculated based on the historical data of Citi Bike, which
consists of the number of bike rentals and returns at each station, recorded every 15~min from 6:00
a.m. to 12:00p.m., over the course of a single day.

\protect\phantomsection\label{p0625}
All instances used in this section were generated by randomly selecting a certain number of stations
from the 518-station dataset. Specifically, in parameter tuning in Section 5.1, a 60-station network
was generated. In the experiment in Section 5.2, five different instances were generated for each
network size (6, 10, and 15 stations). In the experiment in Section 5.3, instances with network
sizes ranging from 60 to 500 stations were created, with one instance generated for each size. In
the experiment in Section 5.4, one of the 15-station networks generated in Section 5.2 was used. As
the dataset provides no information on the initial inventory of broken bikes, we set the initial
quantity of broken bikes \(b_{i}^{0}\) at station \emph{i} as follows: in the experiments in 5.1
Parameter tuning and setting for the algorithm, 5.2 Comparison of the results between HGSDAC-SBC and
Gurobi, 5.3 Comparison of results from HGSDAC-SBC and Gurobi on larger networks, \(b_{i}^{0}\) was
assigned a pre-generated random value within the range
\(\left\lbrack {0,C_{i} - p_{i}^{0}} \right\rbrack\), while in the experiment in Section 5.4, it was
set proportional to the residual capacity \(C_{i} - p_{i}^{0}\). Note that the time-indexed
formulation for the MILP converts all travel times into integer intervals by dividing the actual
travel time by \emph{τ} and rounding up to the nearest integer, which can lead to an inflated time
representation caused by such an approximation. Consequently, the lower bounds obtained from Gurobi
in this section may not be the accurate bounds for the actual (continuous-time) problem.

\protect\phantomsection\label{p0630}
The setting for other parameters of the experiments is listed in Table 5.1. Note that except for the
experiment in Section 5.3, in which we vary the number of trucks and repairers on larger networks,
we fixed the number of trucks and repairers to one. For the repositioning time budget, we uniformly
set it to two hours in 5.1 Parameter tuning and setting for the algorithm, 5.2 Comparison of the
results between HGSDAC-SBC and Gurobi, 5.4 Cost-effectiveness of introducing a repairer under
varying repair times and broken bike inventory levels. In Section 5.3, for larger instances,
experiments were conducted under both two-hour and three-hour time budgets. In the time
discretization, we used 600~s as the length of each period \emph{τ} for most experiments, except in
Section 5.4, where repair times vary between 300 and 1200~s. To accommodate the largest repair time
in these settings, in Section 5.4, we set the length of each period \emph{τ} to 1200~s, ensuring
that repair times do not exceed the period as defined in the parameter table in Section 3. The codes
and the generated datasets used in the experiments are available at
https://github.com/rqhu1995/brpwr.

\protect\phantomsection\label{s0110}
\subsubsection{5.1. Parameter tuning and setting for the algorithm}\label{st130}

\protect\phantomsection\label{p0635}
The parameter setting mostly follows that of Li et al. (2016), who referred to the research
conducted by Vidal et al., 2012, Vidal et al., 2013, Vidal et al., 2014). The only parameter that
differs from theirs and requires tuning is \emph{η}, the one related to the penalty of infeasible
solutions. \emph{η} was set to take values from the set \{10, 50, 100, 150, 200, 300, 400, 500\}. We
ran the algorithm, setting \(\eta\) to be each value from the candidate set. The result of the
average objective value in 20 runs under different values of \(\eta\) is given in Fig. 5.1., and 100
was chosen because it yields the best average result within a shorter CPU time. The detailed
parameter setting is summarized in Table 5.2.

\begin{figure}
\centering
Fig. 5.1. Parameter tuning results for the penalty coefficient \emph{η}
(\textbar{}\emph{S}\textbar{} = 60, \emph{T} = 2~h, \emph{\textbar K\textbar{} =
\textbar R}\textbar{} = 1)\emph{.}
\caption{}\label{f0035}
\end{figure}

\protect\phantomsection\label{t0015}
Table 5.2. The parameter setting for HGSDAC-SBC.

{\def\LTcaptype{none} % do not increment counter
\begin{longtable}[]{@{}ll@{}}
\toprule\noalign{}
Parameter & Value \\
\midrule\noalign{}
\endhead
\bottomrule\noalign{}
\endlastfoot
The minimum subpopulation size \(\mu\) & 25 \\
The number of offspring in a generation \(\lambda\) & 40 \\
The target proportion of feasible individuals \emph{ξ}\textsuperscript{REF} & 0.2 \\
The maximum number of iterations without improvement \emph{It}\textsubscript{NI} & 5000 \\
The number of successive iterations that triggers the diversification operation
\emph{It}\textsubscript{div} & 200 (0.4\emph{It}\textsubscript{NI}) \\
The time limit for executing the algorithm \emph{T}\textsubscript{max} & 7200~s \\
Probability for repairing an infeasible solution \emph{P}\textsubscript{repair} & 0.5 \\
The number of elite individuals that survive to the next generation \emph{N}\textsubscript{elite} &
10 \\
Penalty coefficient for solution infeasibility \(\eta\) & 100 \\
\end{longtable}
}

\protect\phantomsection\label{s0115}
\subsubsection{\texorpdfstring{5.2. Comparison of the results between HGSDAC-SBC and
Gurobi\textsuperscript{a}}{5.2. Comparison of the results between HGSDAC-SBC and Gurobia}}\label{st135}

\protect\phantomsection\label{p0640}
This section aims to compare the results obtained from HGSDAC-SBC with the optimal solutions
obtained by Gurobi for small network instances. Across all instances, the optimal solutions could be
found using both the solver and the heuristic for all the tested instances, and when applied with
different random seeds 20 times to a single instance, the heuristic consistently found the optimal
solution. Therefore, we only report the CPU times using Gurobi and the average CPU times over 20
runs with HGSDAC-SBC in Table 5.3.

\protect\phantomsection\label{t0020}
Table 5.3. A comparison between the CPU times (s) of Gurobi and HGSDAC-SBC on small networks
(Repositioning time budget \emph{T} = 2~h, \emph{\textbar K\textbar{} = \textbar R}\textbar{} = 1,
\emph{τ~=}~600~s).

{\def\LTcaptype{none} % do not increment counter
\begin{longtable}[]{@{}llll@{}}
\toprule\noalign{}
\textbar{}\emph{S}\textbar{} & instance & Gurobi & HGSDAC-SBC \\
\midrule\noalign{}
\endhead
\bottomrule\noalign{}
\endlastfoot
\multirow{5}{*}{6} & 1 & 13.41 & 6.29 \\
& 2 & 13.85 & 5.75 \\
& 3 & 66.97 & 6.22 \\
& 4 & 2.29 & 3.63 \\
& 5 & 22.29 & 4.84 \\
\multirow{5}{*}{10} & 1 & 928.50 & 3.87 \\
& 2 & 1409.89 & 5.27 \\
& 3 & 13768.72 & 3.02 \\
& 4 & 2404.03 & 5.16 \\
& 5 & 4172.75 & 4.17 \\
\multirow{5}{*}{15} & 1 & 2602.24 & 7.80 \\
& 2 & 4209.67 & 4.89 \\
& 3 & 6283.58 & 4.72 \\
& 4 & 3374.38 & 5.32 \\
& 5 & 2371.40 & 5.08 \\
\end{longtable}
}

\protect\phantomsection\label{p0645}
In most instances, HGSDAC-SBC outperformed Gurobi in terms of CPU times, producing optimal solutions
in significantly less time. Noticeably, in one of the 10-station instances (instance 3), HGSDAC-SBC
completed the computations in as little as 3.02~s on average, whereas Gurobi required over 13,000~s
(nearly four hours) for the same instance. However, in a smaller 6-station instance (instance 4),
Gurobi required slightly less CPU time than HGSDAC-SBC. This is likely because a significant portion
of the computing time of the heuristic was spent on the time required to check the convergence.
Despite this, the overall computational benefits of HGSDAC-SBC are evident as the network size
increases, making it a more efficient choice for large-scale problems. It is also worth noting that
the time-index formulation can make the model computationally hard, even for instances comprising
just a small number of stations, further showing the necessity of using heuristics to solve the
problem efficiently.

\protect\phantomsection\label{s0120}
\subsubsection{5.3. Comparison of results from HGSDAC-SBC and Gurobi on larger
networks}\label{st140}

\protect\phantomsection\label{p0650}
In this section, we proceed to compare the results obtained from HGSDAC-SBC with the best feasible
solutions found by Gurobi for larger network instances with different numbers of trucks and
repairers. Gurobi was given a two-hour time limit for computation in all instances in this section.
In some cases, Gurobi terminated earlier due to running out of memory. Hence, no CPU times are
reported for those instances (marked as N/A). The results are reported in Table 5.4\textbf{.}

\protect\phantomsection\label{t0025}
Table 5.4. A comparison between the results of Gurobi and HGSDAC-SBC on larger networks
(Repositioning time budget \emph{T} = 2~h, \emph{τ~=}~600~s).

{\def\LTcaptype{none} % do not increment counter
\begin{longtable}[]{@{}lllllllllllll@{}}
\toprule\noalign{}
\multirow{2}{*}{\textbar{}\emph{S}\textbar{}} & \multirow{2}{*}{\textbar{}\emph{K}\textbar{}} &
\multirow{2}{*}{\textbar{}\emph{R}\textbar{}} & \multicolumn{3}{l}{%
Gurobi} & \multirow{2}{*}{Empty Cell} & \multicolumn{4}{l}{%
HGSDAC-SBC} & \multirow{2}{*}{Empty Cell} & \multirow{2}{*}{\textsuperscript{5}Gap (\%)} \\
& & & UB & LB & CPU (s) & & \textsuperscript{1}Obj. Avg. & \textsuperscript{2}Obj. Best. &
\textsuperscript{3}Obj. Std. & \textsuperscript{4}CPU (s) \\
\midrule\noalign{}
\endhead
\bottomrule\noalign{}
\endlastfoot
\multirow{4}{*}{60} & 1 & 1 & 2604.63 & 2517.12 & 7200.00 & & 2581.29 & 2579.82 & 0.86 & 6.37 & &
0.90 \\
& 1 & 2 & 2533.44 & 2420.32 & 7200.00 & & 2494.24 & 2486.32 & 4.09 & 9.20 & & 1.55 \\
& 2 & 1 & 3581.80 & 2284.28 & 7200.00 & & 2400.65 & 2379.21 & 10.18 & 9.87 & & 32.98 \\
& 2 & 2 & 3517.71 & 2202.55 & 7200.00 & & 2324.50 & 2302.98 & 11.76 & 10.90 & & 33.92 \\
\multirow{4}{*}{90} & 1 & 1 & 4510.58 & 3893.04 & 7200.00 & & 4038.28 & 4014.71 & 13.59 & 8.81 & &
10.47 \\
& 1 & 2 & 4159.43 & 3782.32 & 7200.00 & & 3941.78 & 3915.27 & 14.03 & 13.40 & & 5.23 \\
& 2 & 1 & 4510.58 & 3656.81 & 7200.00 & & 3862.44 & 3829.09 & 19.77 & 9.56 & & 14.37 \\
& 2 & 2 & 4488.64 & 3557.92 & 7200.00 & & 3768.58 & 3733.90 & 21.59 & 13.28 & & 16.04 \\
\multirow{4}{*}{120} & 1 & 1 & 5113.72 & 4247.68 & 7200.00 & & 4303.05 & 4293.75 & 3.58 & 9.13 & &
15.85 \\
& 1 & 2 & 5103.87 & 4169.64 & 7200.00 & & 4237.52 & 4226.39 & 4.17 & 14.96 & & 16.97 \\
& 2 & 1 & 5113.72 & 4024.31 & 7200.00 & & 4167.75 & 4129.41 & 16.19 & 11.45 & & 18.50 \\
& 2 & 2 & 5103.87 & 3951.26 & 7200.00 & & 4097.19 & 4071.74 & 11.87 & 15.98 & & 19.72 \\
\multirow{4}{*}{200} & 1 & 1 & 8657.30 & 2919.70 & 7200.00 & & 7821.12 & 7793.53 & 19.35 & 12.36 & &
9.66 \\
& 1 & 2 & 8657.30 & 7563.69 & 7200.00 & & 7719.58 & 7694.50 & 12.08 & 23.96 & & 10.83 \\
& 2 & 1 & 8657.30 & 7378.11 & 7200.00 & & 7641.82 & 7573.84 & 28.23 & 14.19 & & 11.73 \\
& 2 & 2 & 8657.30 & 7286.86 & 7200.00 & & 7561.16 & 7511.34 & 22.96 & 23.90 & & 12.66 \\
\multirow{4}{*}{300} & 1 & 1 & 11717.13 & 10898.80 & 7200.00 & & 11169.22 & 11134.10 & 15.16 & 16.55
& & 4.68 \\
& 1 & 2 & 11716.99 & 10757.43 & 7200.00 & & 11057.76 & 10990.90 & 19.38 & 25.67 & & 5.63 \\
& 2 & 1 & 11717.13 & 10561.75 & 7200.00 & & 10989.39 & 10953.90 & 19.84 & 19.59 & & 6.21 \\
& 2 & 2 & 11716.99 & 10442.63 & 7200.00 & & 10888.38 & 10848.40 & 20.63 & 33.70 & & 7.07 \\
\multirow{4}{*}{400} & 1 & 1 & 15068.74 & 6026.88 & 7200.00 & & 14280.67 & 14267.10 & 5.89 & 25.27 &
& 5.23 \\
& 1 & 2 & 15060.58 & 6026.88 & 7200.00 & & 14178.63 & 14149.00 & 12.33 & 34.05 & & 5.86 \\
& 2 & 1 & 15068.74 & 6026.59 & 7200.00 & & 14169.14 & 14104.30 & 21.84 & 25.92 & & 5.97 \\
& 2 & 2 & 15060.58 & 6026.59 & 7200.00 & & 14059.18 & 14033.60 & 15.29 & 43.77 & & 6.65 \\
\multirow{4}{*}{500} & 1 & 1 & 19793.68 & 7716.08 & 7200.00 & & 18961.58 & 18927.40 & 18.77 & 27.36
& & 4.20 \\
& 1 & 2 & 19785.07 & 7716.08 & 7200.00 & & 18793.14 & 18734.00 & 22.68 & 41.60 & & 5.01 \\
& 2 & 1 & 19793.68 & 7715.69 & N/A & & 18752.12 & 18618.10 & 50.56 & 33.65 & & 5.26 \\
& 2 & 2 & 19785.07 & N/A & N/A & & 18612.62 & 18505.40 & 37.57 & 51.13 & & 5.93 \\
\end{longtable}
}

\begin{description}
\item[1]
\protect\phantomsection\label{np075}
the average objective value for each instance in 20 runs.
\item[2]
\protect\phantomsection\label{np080}
the best objective value for each instance in 20 runs.
\item[3]
\protect\phantomsection\label{np085}
the standard deviation of the objective value for each instance in 20 runs.
\item[4]
\protect\phantomsection\label{np090}
average computational time for each instance in 20 runs in seconds.
\item[5]
\protect\phantomsection\label{np095}
(UB -- Obj. Avg) / UB.
\end{description}

\protect\phantomsection\label{p0655}
As shown in both Table 5.4 and Table 5.5, Gurobi was unable to find optimal solutions within the
two-hour CPU time limit in any of the instances. For all instances where memory did not run out,
Gurobi used the entire two-hour time limit to provide the best feasible solution (UB). The
difficulty in solving these larger instances again highlights the computational complexity of the
problem.

\protect\phantomsection\label{t0030}
Table 5.5. A comparison between the results of Gurobi and HGSDAC-SBC on larger networks
(Repositioning time budget \emph{T~=}~3~h, \emph{τ~=}~600~s).

{\def\LTcaptype{none} % do not increment counter
\begin{longtable}[]{@{}lllllllllllll@{}}
\toprule\noalign{}
\multirow{2}{*}{\textbar{}\emph{S}\textbar{}} & \multirow{2}{*}{\textbar{}\emph{K}\textbar{}} &
\multirow{2}{*}{\textbar{}\emph{R}\textbar{}} & \multicolumn{3}{l}{%
Gurobi} & \multirow{2}{*}{Empty Cell} & \multicolumn{4}{l}{%
HGSDAC-SBC} & \multirow{2}{*}{Empty Cell} & \multirow{2}{*}{\textsuperscript{5}Gap (\%)} \\
& & & UB & LB & CPU (s) & & \textsuperscript{1}Obj. Avg. & \textsuperscript{2}Obj. Best. &
\textsuperscript{3}Obj. Std. & \textsuperscript{4}CPU (s) \\
\midrule\noalign{}
\endhead
\bottomrule\noalign{}
\endlastfoot
\multirow{4}{*}{60} & 1 & 1 & 3553.62 & 2337.38 & 7200.00 & & 2433.29 & 2412.09 & 3.14 & 10.50 & &
31.53 \\
& 1 & 2 & 3507.20 & 2221.33 & 7200.00 & & 2364.56 & 2352.33 & 5.67 & 12.33 & & 32.58 \\
& 2 & 1 & 3581.80 & 2035.17 & 7200.00 & & 2264.57 & 2248.72 & 8.76 & 13.65 & & 36.78 \\
& 2 & 2 & 3507.20 & 1930.24 & 7200.00 & & 2214.40 & 2196.06 & 6.44 & 12.99 & & 36.86 \\
\multirow{4}{*}{90} & 1 & 1 & 4503.66 & 3705.60 & 7200.00 & & 3903.04 & 3865.03 & 16.13 & 10.06 & &
13.34 \\
& 1 & 2 & 4481.06 & 3563.25 & 7200.00 & & 3778.54 & 3735.60 & 20.69 & 14.04 & & 15.68 \\
& 2 & 1 & 4503.66 & 3398.88 & 7200.00 & & 3705.04 & 3652.58 & 20.54 & 15.00 & & 17.73 \\
& 2 & 2 & 4481.06 & 3264.80 & 7200.00 & & 3606.92 & 3563.30 & 21.66 & 19.82 & & 19.51 \\
\multirow{4}{*}{120} & 1 & 1 & 5113.72 & 4085.94 & 7200.00 & & 4203.49 & 4180.03 & 9.81 & 8.84 & &
17.80 \\
& 1 & 2 & 5103.87 & 3982.62 & 7200.00 & & 4106.56 & 4087.78 & 8.83 & 18.28 & & 19.54 \\
& 2 & 1 & 5113.72 & 3784.05 & 7200.00 & & 4023.59 & 4007.56 & 10.96 & 13.59 & & 21.32 \\
& 2 & 2 & 5103.87 & 3685.32 & 7200.00 & & 3952.46 & 3909.70 & 17.55 & 17.09 & & 22.56 \\
\multirow{4}{*}{200} & 1 & 1 & 8657.30 & 7457.23 & 7200.00 & & 7695.07 & 7636.22 & 22.75 & 14.47 & &
11.11 \\
& 1 & 2 & 8657.30 & 7322.35 & 7200.00 & & 7575.25 & 7539.85 & 14.29 & 25.61 & & 12.50 \\
& 2 & 1 & 8657.30 & 2916.37 & 7200.00 & & 7471.22 & 7406.93 & 32.94 & 19.45 & & 13.70 \\
& 2 & 2 & 8657.30 & 2916.37 & 7200.00 & & 7393.52 & 7333.03 & 30.51 & 27.84 & & 14.60 \\
\multirow{4}{*}{300} & 1 & 1 & 11715.76 & 4568.18 & 7200.00 & & 11024.28 & 10973.30 & 21.87 & 18.06
& & 5.90 \\
& 1 & 2 & 11715.62 & 10462.88 & 7200.00 & & 10901.05 & 10858.60 & 16.88 & 38.59 & & 6.95 \\
& 2 & 1 & 11715.76 & 4568.16 & 7200.00 & & 10830.35 & 10756.00 & 23.79 & 27.03 & & 7.56 \\
& 2 & 2 & 11715.62 & 4568.16 & 7200.00 & & 10692.76 & 10638.50 & 25.19 & 44.10 & & 8.73 \\
\multirow{4}{*}{400} & 1 & 1 & 15068.74 & 6026.59 & 7200.00 & & 14137.72 & 14090.00 & 20.16 & 21.88
& & 6.18 \\
& 1 & 2 & 15060.58 & 6026.59 & 7200.00 & & 14009.58 & 13974.40 & 15.25 & 45.86 & & 6.98 \\
& 2 & 1 & 15068.74 & 6026.58 & N/A & & 13968.64 & 13926.70 & 21.06 & 32.79 & & 7.30 \\
& 2 & 2 & 15060.58 & N/A & N/A & & 13851.78 & 13822.50 & 12.96 & 55.96 & & 8.03 \\
\multirow{4}{*}{500} & 1 & 1 & 19793.68 & 7715.71 & 7200.00 & & 18818.29 & 18716.10 & 37.12 & 21.88
& & 4.93 \\
& 1 & 2 & 19785.07 & N/A & N/A & & 18596.33 & 18539.40 & 26.59 & 56.24 & & 6.01 \\
& 2 & 1 & 19793.68 & N/A & N/A & & 18572.48 & 18493.30 & 38.87 & 35.34 & & 6.17 \\
& 2 & 2 & 19785.07 & N/A & N/A & & 18407.06 & 18250.20 & 44.37 & 62.35 & & 6.96 \\
\end{longtable}
}

\begin{description}
\item[1]
\protect\phantomsection\label{np100}
the average objective value for each instance in 20 runs.
\item[2]
\protect\phantomsection\label{np105}
the best objective value for each instance in 20 runs.
\item[3]
\protect\phantomsection\label{np110}
the standard deviation of the objective value for each instance in 20 runs.
\item[4]
\protect\phantomsection\label{np115}
average computational time for each instance in 20 runs in seconds.
\item[5]
\protect\phantomsection\label{np120}
(UB -- Obj. Avg) / UB.
\end{description}

\protect\phantomsection\label{p0660}
On the other hand, HGSDAC-SBC consistently outperformed Gurobi by finding better feasible solutions
(as reflected by a positive gap) in significantly less time. As the repositioning time budget
increased from two hours to three hours in the model, the performance gap between HGSDAC-SBC and
Gurobi widened, especially for instances comprising 60 stations. For example, in the 60-station
instance with one truck and two repairers, the gap increased from 1.55~\% under the two-hour budget
to 32.58~\% under the three-hour budget. This indicates that the increased complexity of the model
with longer time budgets makes it more difficult for Gurobi to find better solutions, while
HGSDAC-SBC is still efficient in obtaining better solutions. Additionally, we notice that the
computation for some very large instances (\textbar{}\emph{S}\textbar{} = 400 and 500) used up the
64~GB memory of the computer before the two-hour time limit was reached. This suggests that as
network size increases and resources become a limitation for Gurobi, HGSDAC-SBC shows an additional
advantage by requiring less memory and less CPU time, making it more suitable for real-world BRP
with large networks.

\protect\phantomsection\label{s0125}
\subsubsection{5.4. Cost-effectiveness of introducing a repairer under varying repair times and
broken bike inventory levels}\label{st145}

\protect\phantomsection\label{p0665}
In this section, we analyze the cost-effectiveness of involving a repairer under different repair
times and broken bike quantities. The experiment was conducted on a fixed 15-station network, with
the number of broken bikes at each station being proportional to the station's residual capacity
with a proportional coefficient denoted by \(\ell\). The initial broken bike quantity at station
\emph{i} was set to
\(b_{i}^{0} = \left\lfloor {\ell \cdot \left( {C_{i} - p_{i}^{0}} \right)} \right\rfloor\). \(\ell\)
ranges from 10~\% to 100~\% in an increment of 10~\%. For each value of \(\ell\), the repair time
per bike varies from 300~s to 1200~s, with a step size of 100~s. The length of period \emph{τ} is
set to 1200~s in this section to guarantee \emph{t}\textsuperscript{repair}~\emph{\textless~τ}.

\protect\phantomsection\label{p0670}
To capture the cost-effectiveness of a repairer, an indicator \(\nu\) is introduced, which is
defined as the ratio of the reduction cost in user dissatisfaction (when a repairer is introduced)
to the wage paid to the repairer, which is calculated
as\[\tag{73}\nu = \frac{c^{\text{p}}\sum_{i \in S}\left( {{\widehat{F}}_{i} - {\widehat{F}}_{i}'} \right)}{\left. \varpi \cdot T/3600 \right.}\]where
\({\widehat{F}}_{i}\) and \({\widehat{F}}_{i}'\) are the user dissatisfaction after the
repositioning at station \emph{i} with and without a repairer, respectively, \(c^{\text{p}}\) is the
unit cost of the dissatisfaction, \(\varpi\) is the hourly wage of a repairer, and \emph{T} is the
time budget for the repositioning in seconds. The division by 3600 converts the time from seconds to
hours, ensuring consistency with the wage unit in CNY/h. The numerator is the reduction in the
dissatisfaction cost, while the denominator is the cost paid to the repairer. Note that the
dissatisfaction after the repositioning with no repairers \({\widehat{F}}_{i}'\) was calculated by
setting \textbar{}\emph{R}\textbar{} = 0 in the optimization model. The hourly wage was determined
based on a piece of job-hunting information online,\textsuperscript{b} and was set to 33.33 CNY/h.

\protect\phantomsection\label{p0675}
For each setting of \(\ell\), experiments under varying repair times were conducted, and the results
of the variations in \(\nu\) over repair times are plotted in Fig. 5.2. The results show a clear
trend across all instances: As repair time increases, the cost-effectiveness ratio decreases
consistently, meaning that on-site repairs become less cost-effective. This implies that a longer
repair time leads to a smaller reduction in user dissatisfaction, as fewer repairs can be completed
by the repairer within the available time. In the extreme case, when the repair time exceeds a
certain threshold, introducing on-site repairs is cost-ineffective.

\begin{figure}
\centering
Fig. 5.2. Variations in the cost-effectiveness ratio over repair times under different broken bike
inventories.
\caption{}\label{f0040}
\end{figure}

\protect\phantomsection\label{p0680}
One key observation is the existence of a critical point at which the cost-effectiveness ratio
\emph{ν} equals 1.0 at some level of broken bikes. For lower levels of broken bikes (10\%\ensuremath{\leq} ℓ \ensuremath{\leq}
20\%), this ratio remains below 1.0 regardless of the repair times, indicating that the cost of
employing a repairer outweighs the benefits. In such cases, the system can still be effectively
managed by trucks alone for repositioning, as the low number of broken bikes does not justify the
additional cost of involving a repairer. However, for systems with higher levels of broken bikes
(ℓ\ensuremath{\geq}30\%), a turning point emerges where the cost-effectiveness ratio shifts from above to below 1.0.
This signifies that as the proportion of broken bikes increases, the benefit of reducing user
dissatisfaction gradually begins to outweigh the cost of employing a repairer at lower ranges of
repair times. Larger proportions of broken bikes shift the incorporation of a repairer from being
cost-ineffective to cost-effective at lower ranges of repair times, reflecting the growing necessity
for on-site repairs as the average damage level within the system rises. The changes in the
cost-effectiveness ratio over the broken bike inventory levels suggest the necessity for a dynamic
assignment of repairers based on the actual level of damage within the system, which includes
varying the number of repairers in response to the dynamic change in the broken bike levels in the
system to ensure the cost-effectiveness. Specifically, part-time repairer scheduling can serve as a
flexible strategy to further enhance resource allocation, allowing systems to adjust the input for
repairs dynamically.

\protect\phantomsection\label{p0685}
Another important observation is the shift in the repair time at which the cost-effectiveness ratio
\(\nu\) equals 1.0. As the level of broken bikes increases, this critical point moves toward longer
repair times. For example, when the broken bike inventory reaches 50~\%, the point at which
\(\nu\)~=~1.0 occurs at approximately 492~s, whereas for the 90~\% level, it occurs at around 580~s.
This suggests that with a higher number of broken bikes, longer repair times can still yield a
favorable cost-effectiveness ratio before the repairer's cost outweighs the benefit. However, once
repair times exceed around 645~s, employing a repairer becomes less cost-effective across all
instances, regardless of the level of broken bike inventory. Although our model assumes the same
repair time for all bikes, this can be viewed as an average repair time, representing the overall
damage level of the system. Hence, such a relatively long repair time suggests that the system is
experiencing significant damage. Therefore, it may be more effective to shift resources from
reactive repairs to preventive maintenance. This strategy can be adapted based on the system's
changing damage degrees, helping to reduce the number of broken bikes and the average repair time
over time and preventing the system from reaching a state where repairs become less cost-effective.

\protect\phantomsection\label{s0130}
\subsection{6. Conclusion}\label{st150}

\protect\phantomsection\label{p0690}
We investigate a new and practical static BRP with broken bikes and on-site repairers. The problem
is to find the best routes and loading/unloading instructions of both usable and broken bikes for
the repositioning trucks while determining the best routes for on-site repairers and the number of
broken bikes repaired by each repairer at each station. The goal is to reduce the sum of the penalty
costs of user dissatisfaction and the cost of total CO\textsubscript{2} emissions. The problem is
formulated as a time-index Mixed Integer Linear Programming (MILP) model. A hybrid algorithm
HGSDAC-SBC is proposed to solve the problem efficiently for large instances. The algorithm relies on
a novel concept of station budget, which utilizes the Benefit-Cost Ratio Function (BCRF) to
determine (un)loading and repair quantities at stations, ensuring that time is spent effectively
across the network.

\protect\phantomsection\label{p0695}
Numerical experiments on small networks of up to 15 stations demonstrated that HGSDAC-SBC
consistently achieves optimal solutions, matching Gurobi's results while requiring significantly
less computational time in most cases. For larger networks, ranging from 60 to 500 stations,
HGSDAC-SBC outperformed Gurobi in finding a feasible solution obtained under a two-hour time limit,
demonstrating the algorithm's computational efficiency. Additionally, we analyzed the
cost-effectiveness of introducing a repairer under varying repair times and broken bike inventory
levels. The findings show that long repair times reduce cost-effectiveness. In the extreme case,
when the repair time exceeds a certain level, introducing on-site repairs is cost-ineffective.
Moreover, when the number of broken bikes in the system is small, introducing on-site repairs can
become detrimental. The results also underscore the need for a dynamic repairer allocation strategy
based on the system's actual damage level and suggest that preventive maintenance strategies can
reduce the average number of broken bikes in the system and the average repair time over time.

\protect\phantomsection\label{p0700}
It is worth noting that we conducted additional experiments comparing a full station set with a
reduced station set, where stations with no broken bicycles and inventory levels that (nearly)
minimize the EUDF were removed. The results showed no statistically significant difference in terms
of objective values or CPU times between the two approaches. This indicates that the
heuristic\textquotesingle s natural allocation of negligible station budgets to such stations
achieves a similar effect to explicitly removing them from the problem, supporting the efficiency
and robustness of the current approach. Since the manuscript is already lengthy, we omit the
detailed experimental results in the manuscript.

\protect\phantomsection\label{p0705}
While this study addresses key aspects of the problem, two limitations must be acknowledged to
provide a comprehensive perspective on the findings. The time-indexed formulation, although
necessary for the MILP model to handle multiple visits to stations by both trucks and repairers,
introduces a discrepancy between the discrete and continuous-time problems by approximating travel
times through integer conversion, limiting the accuracy of the model in reflecting the
continuous-time problem. Additionally, the SBC heuristic assumes that bicycles are (un)loaded or
repaired directly at their respective stations without intermediate stops or second-time transfers.
This simplification streamlines implementation but restricts the heuristic from leveraging potential
efficiency gains through temporary inventories, which could yield better solutions. Future research
will develop methodologies to address these limitations.

\protect\phantomsection\label{p0710}
\textbf{Declarations of interest}: This research was jointly supported by grants from the National
Natural Science Foundation of China (No. 71771194), the Research Grants Council of the Hong Kong
Special Administrative Region of China (17206322), and a grant from the University Research
Committee of the University of Hong Kong (2202100693).

\protect\phantomsection\label{s0140}
\subsection{CRediT authorship contribution statement}\label{st160}

\protect\phantomsection\label{p0720}
\textbf{Runqiu Hu:} Writing -- original draft, Visualization, Methodology, Investigation, Formal
analysis, Data curation, Conceptualization. \textbf{W.Y. Szeto:} Writing -- review \& editing,
Supervision, Project administration, Methodology, Funding acquisition. \textbf{Sin C. Ho:} Writing
-- review \& editing, Formal analysis, Data curation.

\protect\phantomsection\label{coi005}
\subsection{Declaration of competing interest}\label{st165}

\protect\phantomsection\label{p0725}
The authors declare that they have no known competing financial interests or personal relationships
that could have appeared to influence the work reported in this paper.

\protect\phantomsection\label{ak005}
\subsection{Acknowledgments}\label{st170}

\protect\phantomsection\label{p0730}
This research was jointly supported by grants from the National Natural Science Foundation of China
(No. 71771194), the Research Grants Council of the Hong Kong Special Administrative Region of China
(17206322), and a grant from the University Research Committee of the University of Hong Kong
(2202100693). The authors are grateful to the constructive comments from the reviewers.

\protect\phantomsection\label{s0145}
\subsection{Appendix A. Linearization of the EUDF}\label{st175}

\protect\phantomsection\label{p0735}
The Extended User Dissatisfaction Function (EUDF) is a bivariate function that represents a weighted
sum of the expected shortages of usable bikes and the expected shortages of lockers at a
bike-sharing station in a day. It is calculated based on the quantities of both usable and unusable
bikes at the station. The calculation of the EUDF values follows a Markov process, which captures
the dynamics of bicycle inventory levels over time. For detailed steps for how the EUDF values are
calculated, readers can refer to the approach outlined in the paper of Kaspi et al. (2017).

\protect\phantomsection\label{p0740}
Let \(F_{i}\left( {p_{i}',b_{i}'} \right)\) be the EUDF, where \(p_{i}'\) and \(b_{i}'\) are the
usable and broken inventory levels at the station, respectively. The domain of EUDF at each station
\emph{i} is
\(\Theta_{i} = \left\{ {\left( {p_{i}',b_{i}'} \right) \mid p_{i}' \geq 0,b_{i}' \geq 0,p_{i}' + b_{i}' \leq C_{i}} \right\}\),
where \emph{C\textsubscript{i}} is the station capacity. For each point
\(\left( {p_{i}',b_{i}'} \right) \in \Theta_{i},\) a plane
\(\alpha_{p_{i},b_{i}}^{i} \cdot p_{i}' + \beta_{p_{i},b_{i}}^{i} \cdot b_{i}' + \gamma_{p_{i},b_{i}}^{i}\)
that passes as close as possible to \(F_{i}\left( {p_{i}',b_{i}'} \right)\) can be fitted, where
\(\alpha_{p_{i},b_{i}}^{i},\beta_{p_{i},b_{i}}^{i},\) and \(\gamma_{p_{i},b_{i}}^{i}\) are the
coefficients of the plane that can be computed for every
\(\left( {p_{i},b_{i}} \right) \in \Theta_{i}\). These coefficients are obtained by solving a linear
programming (LP) model which ensures that the polyhedral
function\({\overset{\sim}{F}}_{i}\left( {p_{i}^{'},b_{i}^{'}} \right)\)~=~\(\max\limits_{{({p_{i},b_{i}})} \in \theta}\left( {\alpha_{p_{i},b_{i}}^{i} \cdot p_{i}' + \beta_{p_{i},b_{i}}^{i} \cdot b_{i}' + \gamma_{p_{i},b_{i}}^{i}} \right)\)remains
below or equal to the EUDF values at all points except for
\(\left( {p_{i},b_{i}} \right) \in \theta_{i}\), with the objective of minimizing the gap between
the fitted function and the actual EUDF at each point. The LP model is given as follows:

\protect\phantomsection\label{p0745}
\textbf{Decision variables:}

\protect\phantomsection\label{p0750}
\(\alpha_{p_{i},b_{i}}^{i},\beta_{p_{i},b_{i}}^{i},\)\(\gamma_{p_{i},b_{i}}^{i}\) Coefficients of
the fitted plane

\protect\phantomsection\label{p0755}
\(s\) Gap at the point \(\left( {p_{i},b_{i}} \right)\) between the fitted plane and the EUDF value

\protect\phantomsection\label{p0760}
\textbf{The LP formulation:}

\protect\phantomsection\label{p0765}
Minimize \(s\) (A.74)

\protect\phantomsection\label{p0770}
s.t.\[\tag{A.75}\alpha_{p_{i},b_{i}}^{i} \cdot p_{i}' + \beta_{p_{i},b_{i}}^{i} \cdot b_{i}' + \gamma_{p_{i},b_{i}}^{i} \leq F\left( {p_{i}',b_{i}'} \right)\quad\forall\left( {p_{i}',b_{i}'} \right)\left. \in \theta\backslash \right.\left\{ \left( {p_{i},b_{i}} \right) \right\}\]

\[\tag{A.76}\alpha_{p_{i},b_{i}}^{i} \cdot p_{i} + \beta_{p_{i},b_{i}}^{i} \cdot b_{i} + \gamma_{p_{i},b_{i}}^{i} = F\left( {p_{i},b_{i}} \right) - s\]

\[\tag{A.77}\alpha_{p_{i},b_{i}}^{i},\beta_{p_{i},b_{i}}^{i},\gamma_{p_{i},b_{i}}^{i}\quad{free},\]

\[\tag{A.78}s \geq 0\]

\protect\phantomsection\label{p0775}
The objective (A.74) is to minimize the gap between the fitted plane and the EUDF value. Constraints
(A.75) ensure that the fitted plane passes under the EUDF at all integer points except for
\(\left( {p_{i},b_{i}} \right)\). Constraint (A.76) defines the gap at the point
\(\left( {p_{i},b_{i}} \right)\) between the fitted plane and the EUDF value. Constraints (A.77) and
(A.78) are the domain constraints for the decision variables.

\protect\phantomsection\label{p0780}
After the polyhedral function \({\overset{\sim}{F}}_{i}\left( {p_{i},b_{i}} \right)\) is formed at
each point \(\left( {p_{i},b_{i}} \right) \in \Theta_{i},\) we replace the EUDF
\(F_{i}\left( {p_{i,T'},b_{i,T'}} \right)\) in the objective function with an auxiliary variable
\({\widehat{F}}_{i}\) and add the following constraints to the
model:\[\tag{A.79}{\hat{F}}_{i} \geq \alpha_{p_{i},b_{i}}^{i} \cdot p_{i,T^{'}} + \beta_{p_{i},b_{i}}^{i} \cdot b_{i,T^{'}} + \gamma_{p_{i},b_{i}}^{i},\forall\left( {p_{i},b_{i}} \right) \in \left\{ {\left( {p_{i},b_{i}} \right) \in {\mathbb{Z}}^{2} \mid p_{i} \geq 0,b_{i} \geq 0,p_{i} + b_{i} \leq C_{i}} \right\},\forall i \in S.\]

\protect\phantomsection\label{s0150}
\subsection{Appendix B. SBC heuristic for the quantity assignment}\label{st180}

\protect\phantomsection\label{p0785}
\protect\phantomsection\label{t0035}
{\def\LTcaptype{none} % do not increment counter
\begin{longtable}[]{@{}llllll@{}}
\toprule\noalign{}
\multicolumn{6}{@{}l@{}}{%
\textbf{Algorithm 2:} \(SBCAssignment\left( {S_{\text{R}},S_{\text{K}}} \right)\)} \\
\midrule\noalign{}
\endhead
\bottomrule\noalign{}
\endlastfoot
\textbf{1} & \multicolumn{5}{l@{}}{%
\(p_{s} = p_{s}^{0},b_{s} = b_{s}^{0},\forall s \in S\)// initialization} \\
& \multicolumn{5}{l@{}}{%
// lines 2 to 9: repairing quantity assignment} \\
\textbf{2} & \multicolumn{5}{l@{}}{%
\textbf{for} \(i = 1\;\text{to}\;|R|\) \textbf{do}} \\
\textbf{3} & & \multicolumn{4}{l@{}}{%
// assign the repairing budget for repairer \emph{i} at each visited station (\textbf{Algorithm
3})} \\
& & \multicolumn{4}{l@{}}{%
\(\left( {\mathit{maxTim}e_{s,i}^{\text{r}}} \right)_{s \in S}\)=
\(\mathit{allocateRepairBudget}\left( r_{i} \right)\)} \\
& & \multicolumn{4}{l@{}}{%
// determine the repairing quantity for repairer \emph{i} at each visited station in route \(r_{i}\)
(\textbf{Algorithm 5})} \\
\textbf{4} & & \multicolumn{4}{l@{}}{%
\(\left( g_{s,i} \right)_{s \in S}\)=\(assignRepair\left( {r_{i},\left( {\mathit{maxTim}e_{s,i}^{\text{r}}} \right)_{s \in S}} \right)\)} \\
\textbf{8} & & \multicolumn{4}{l@{}}{%
/* lines 5 to 8: station inventory update (create a duplicate of the usable and broken bike
inventory after repairs in \(p_{s}'\) and \(b_{s}'\) , respectively, in preparation for restoring
the bike inventories to the state after repairs before the second-round truck (un)loading assignment
(Line 37) */} \\
\textbf{5} & & \multicolumn{4}{l@{}}{%
\textbf{for} \(s = 1\;\text{to}\;|S|\) \textbf{do}} \\
\textbf{6} & & & \multicolumn{3}{l@{}}{%
\(p_{s} = p_{s} + g_{s,i}\),\(p_{s}' = p_{s}\)} \\
\textbf{7} & & & \multicolumn{3}{l@{}}{%
\(b_{s} = b_{s} - g_{s,i}\),\(b_{s}' = b_{s}\)} \\
\textbf{8} & & \multicolumn{4}{l@{}}{%
\textbf{end for}} \\
\textbf{9} & \multicolumn{5}{l@{}}{%
\textbf{end for}} \\
& \multicolumn{5}{l@{}}{%
// lines 10 to 36: truck (un)loading quantity assignment (the first round)} \\
\textbf{10} & \multicolumn{5}{l@{}}{%
\textbf{for} \(k = 1\;\text{to}\;|K|\) \textbf{do}} \\
\textbf{11} & & \multicolumn{4}{l@{}}{%
\(\mathit{unsa}t_{s,k}^{\text{load,u}} = unsat_{s,k}^{\text{unload,u}} = unsat_{s,k}^{\text{load,b}} = 0,\forall s \in S\)//
initialization} \\
& & \multicolumn{4}{l@{}}{%
/* assign the station (un)loading budget for truck \emph{k} at each station (\textbf{Algorithm 4})
and record it in \(\left( {\mathit{maxTimeIni}t_{s,k}^{\text{t}}} \right)_{s \in S}\) */} \\
\textbf{12} & & \multicolumn{4}{l@{}}{%
\(\left( {\mathit{maxTimeIni}t_{s,k}^{\text{t}}} \right)_{s \in S}\)=\(\mathit{allocateLoadingBudget}\left( r_{k}' \right)\)} \\
\textbf{13} & & \multicolumn{4}{l@{}}{%
\(\left( {\mathit{maxTimeRemai}n_{s,k}^{\text{t}}} \right)_{s \in S} = \left( {\mathit{maxTimeIni}t_{s,k}^{\text{t}}} \right)_{s \in S}\)
// initialization} \\
& & \multicolumn{4}{l@{}}{%
/* at the first departure from the depot, we load no usable bikes by default, and the adjustment on
depot loading quantities is made in \textbf{Algorithm 8} when we assign unloading quantities at
subsequent deficit stations */} \\
\textbf{14} & & \multicolumn{4}{l@{}}{%
\(y_{0,k}^{\text{u}^{+}} = y_{0,k}^{\text{u}^{-}} = y_{0,k}^{\text{b}^{+}} = y_{0,k}^{\text{b}^{-}} = 0\),\(\psi_{0,k}^{\text{u}} = \psi_{0,k}^{\text{b}} = 0\)} \\
& & \multicolumn{4}{l@{}}{%
// lines 15 to 33: assign loading and unloading quantities for each node in the route} \\
\textbf{15} & & \multicolumn{4}{l@{}}{%
\textbf{for} \(l = 1\;\text{to}\;\left| r_{k}' \right| - 1\) \textbf{do}} \\
\textbf{16} & & & \multicolumn{3}{l@{}}{%
\(\psi_{l,k}^{\text{u}} = \psi_{l,k}^{\text{b}} = 0\)} \\
\textbf{17} & & & \multicolumn{3}{l@{}}{%
\(s' = s_{k_{l}}'\)} \\
\textbf{18} & & & \multicolumn{3}{l@{}}{%
\(y_{l,k}^{\text{u}^{+}} = y_{l,k}^{\text{u}^{-}} = y_{l,k}^{\text{b}^{+}} = y_{l,k}^{\text{b}^{-}} = 0\)//
all (un)loading quantities are initialized with 0} \\
\textbf{19} & & & \multicolumn{3}{l@{}}{%
\textbf{if} \(s' = 0\) \textbf{then}} \\
& & & & \multicolumn{2}{l@{}}{%
/* at intermediate visits to the depot, we load no usable bikes by default, and the adjustment on
the depot (un)loading quantities is made in \textbf{Algorithm 8} when we assign unloading quantities
at subsequent deficit stations */} \\
\textbf{20} & & & & \multicolumn{2}{l@{}}{%
\(y_{l,k}^{\text{u}^{+}} = \psi_{l - 1,k}^{\text{u}},y_{l,k}^{\text{b}^{+}} = \psi_{l - 1,k}^{\text{b}}\)} \\
\textbf{21} & & & \multicolumn{3}{l@{}}{%
\textbf{else}} \\
\textbf{22} & & & & \multicolumn{2}{l@{}}{%
\textbf{if} \(p_{s'} > q_{s'}\) \textbf{then} // for surplus stations, load usable bikes
(\textbf{Algorithm 6}).} \\
\textbf{23} & & & & & \(\begin{matrix}
{\left( {\mathit{maxTimeRemai}n_{s,k}^{\text{t}}} \right)_{s \in S},\left( {\mathit{unsa}t_{s,k}^{\text{load,u}}} \right)_{s \in S},y_{l,k}^{\text{u}^{-}} =} \\
{assignLoadU\left( {r_{k}',l,\left( {\mathit{maxTimeRemai}n_{s,k}^{\text{t}}} \right)_{s \in S},\left( {\mathit{unsa}t_{s,k}^{\text{load,u}}} \right)_{s \in S}} \right)}
\end{matrix}\) \\
\textbf{24} & & & & \textbf{else} & \\
\textbf{25} & & & & & \textbf{if} \(p_{s'} < q_{s'}\) \textbf{then} // for deficit stations, unload
usable bikes (\textbf{Algorithms 7}--\textbf{8}) \\
\textbf{26} & & & & & \(\begin{matrix}
{\left( {\mathit{maxTimeRemai}n_{s,k}^{\text{t}}} \right)_{s \in S},\left( {\mathit{unsa}t_{s,k}^{\text{load,u}}} \right)_{s \in S},\left( {\mathit{unsa}t_{s,k}^{\text{unload,u}}} \right)_{s \in S},y_{l,k}^{\text{u}^{+}} =} \\
{assignUnloadU\begin{pmatrix}
{r_{k}',l,\left( {\mathit{maxTimeRemai}n_{s,k}^{\text{t}}} \right)_{s \in S},\left( {\mathit{unsa}t_{s,k}^{\text{load,u}}} \right)_{s \in S},} \\
\left( {\mathit{unsa}t_{s,k}^{\text{unload,u}}} \right)_{s \in S}
\end{pmatrix}}
\end{matrix}\) \\
\textbf{27} & & & & & \textbf{end if} \\
\textbf{28} & & & & \multicolumn{2}{l@{}}{%
\textbf{end if}} \\
& & & & \multicolumn{2}{l@{}}{%
// station and truck inventory update} \\
\textbf{29} & & & & \multicolumn{2}{l@{}}{%
\(p_{s'} = p_{s'} - y_{l,k}^{\text{u}^{-}} + y_{l,k}^{\text{u}^{+}},\psi_{l,k}^{\text{u}} = \psi_{l,k}^{\text{u}} + y_{l,k}^{\text{u}^{-}} - y_{l,k}^{\text{u}^{+}}\)} \\
& & & & \multicolumn{2}{l@{}}{%
// for both types of stations, load broken bikes (Algorithm 9)} \\
\textbf{30} & & & & \multicolumn{2}{l@{}}{%
\(\begin{matrix}
{\left( {\mathit{maxTimeRemai}n_{s,k}^{\text{t}}} \right)_{s \in S},\left( {\mathit{unsa}t_{s,k}^{\text{load},\text{b}}} \right)_{s \in S},y_{l,k}^{\text{b}^{-}} =} \\
{assignLoadB\left( {r_{k}',l,\left( {\mathit{maxTimeRemai}n_{s,k}^{\text{t}}} \right)_{s \in S},\left( {\mathit{unsa}t_{s,k}^{\text{unload},\text{b}}} \right)_{s \in S}} \right)}
\end{matrix}\)} \\
\textbf{31} & & & & \multicolumn{2}{l@{}}{%
\(b_{s'} = b_{s'} - y_{l,k}^{\text{b}^{-}},\psi_{l,k}^{\text{b}} = \psi_{l,k}^{\text{b}} + y_{l,k}^{\text{b}^{-}}\)//
station and truck inventory update} \\
\textbf{32} & & & \multicolumn{3}{l@{}}{%
\textbf{end if}} \\
\textbf{33} & & \multicolumn{4}{l@{}}{%
\textbf{end for}} \\
& & \multicolumn{4}{l@{}}{%
// \(\mathit{BCR}F_{s}^{\text{t}}\) is updated here in preparation for the reassignment of the
station (un)loading budget} \\
\textbf{34} & & \multicolumn{4}{l@{}}{%
Update \(\mathit{BCR}F_{s}^{\text{t}}\) for each station based on the current station
inventories} \\
& & \multicolumn{4}{l@{}}{%
/* adjust the station un(loading) budget of truck \emph{k} by distributing the unused time to
unbalanced stations with bikes unable to be (un)loaded due to insufficient station (un)loading
budgets (\textbf{Algorithm 10}) */} \\
\textbf{35} & & \multicolumn{4}{l@{}}{%
\(\begin{matrix}
{\left( {\mathit{maxTimeNe}w_{s,k}^{\text{t}}} \right)_{s \in S} =} \\
{adjustLoadingBudget\begin{pmatrix}
{\left( {\mathit{maxTimeIni}t_{s,k}^{\text{t}}} \right)_{s \in S},\left( {\mathit{maxTimeRemai}n_{s,k}^{\text{t}}} \right)_{s \in S},} \\
{\left( {\mathit{unsa}t_{s,k}^{\text{load},\text{u}}} \right)_{s \in S},\left( {\mathit{unsa}t_{s,k}^{\text{unload},\text{u}}} \right)_{s \in S},\left( {\mathit{unsa}t_{s,k}^{\text{load},\text{b}}} \right)_{s \in S}}
\end{pmatrix}}
\end{matrix}\)} \\
\textbf{36} & \multicolumn{5}{l@{}}{%
\textbf{end for}} \\
& \multicolumn{5}{l@{}}{%
/* restore the station inventories to the inventories before the first-round truck (un)loading
quantity assignment */} \\
\textbf{37} & \multicolumn{5}{l@{}}{%
\(p_{s} = p_{s}',b_{s} = b_{s}'\)} \\
\textbf{38} & \multicolumn{5}{l@{}}{%
\textbf{for} \(k = 1\;\text{to}\;|K|\) \textbf{do}} \\
& & \multicolumn{4}{l@{}}{%
/* assign the quantities for bike (un)loading for the second round using the updated station
(un)loading budget */} \\
\textbf{39} & & \multicolumn{4}{l@{}}{%
Set
\(\left( {\mathit{maxTimeRemai}n_{s,k}^{\text{t}}} \right)_{s \in S} = \left( {\mathit{maxTimeNe}w_{s,k}^{\text{t}}} \right)_{s \in S}\)
and repeat the procedure from Line 15 to Line 33} \\
\textbf{40} & \multicolumn{5}{l@{}}{%
\textbf{end for}} \\
\textbf{41} & \multicolumn{5}{l@{}}{%
Conduct amendments to both trucks and repairers regarding the routes and their associated repairing
or (un)loading quantities using the algorithm described in Section 4.4.4.} \\
\textbf{42} & \multicolumn{5}{l@{}}{%
\textbf{return} \emph{G}\textsubscript{R}, \emph{Y}\textsubscript{K}} \\
\end{longtable}
}

The algorithm begins by assigning quantities of repairs for each repairer and updating the station
inventories (Line 1 to Line 9). Following this, the algorithm iteratively traverses the route of
each truck, determining the quantities of loading/unloading based on the status of each station
(Line 10 to Line 40). A key feature of this SBC heuristic is setting a maximum time for bike
loading/unloading for every truck and a maximum time for bike repairing for every repairer, referred
to as the \emph{station budget}. Specifically, we refer to the budget at each station as the
\emph{station} (\emph{un})\emph{loading budget} for trucks and the \emph{station repairing budget}
for repairers. The introduction of the station budget prevents the truck/repairer from excessively
spending time at the initial stations, which can lead to insufficient time for the latter stations.
For each truck, the first round of (un)loading quantity assignments is conducted (Line 15 to Line
33), followed by an adjustment on the station (un)loading budgets based on unsatisfied demand and
unused time budgets in the first round (Lines 34 and 35). Afterward, the algorithm restores the
station inventories to the values before the truck (un)loading quantity assignment (Line 37) and
conducts a second-round assignment to adjust the truck (un)loading quantities (Line 38 to Line 40).
Next, amendments to both trucks and repairers regarding the routes and their associated repairing or
(un)loading quantities are made to utilize the unused time budget resulting from the rounding-down
strategies applied during quantity assignments (Line 41). Finally, the algorithm returns the
assigned repairing and (un)loading quantities for usable and broken bikes (Line 42).

\protect\phantomsection\label{s0155}
\subsection{Appendix C. Assignment algorithms for station repairing budget and station (un)loading
budget}\label{st185}

\protect\phantomsection\label{p0790}
\protect\phantomsection\label{t0040}
{\def\LTcaptype{none} % do not increment counter
\begin{longtable}[]{@{}llll@{}}
\toprule\noalign{}
\multicolumn{4}{@{}l@{}}{%
\textbf{Algorithm 3:} \(\mathit{allocateRepairBudget}\left( r_{i} \right)\)} \\
\midrule\noalign{}
\endhead
\bottomrule\noalign{}
\endlastfoot
\textbf{1} & \multicolumn{3}{l@{}}{%
Initialization:\(\mathit{maxTim}e_{s,i}^{\text{r}} = 0,\forall s \in S\)} \\
\textbf{2} & \multicolumn{3}{l@{}}{%
\(T^{\text{op}} = T - \sum\limits_{j = 0}^{|r_{i}{| - 1}}t_{s_{i_{j}},s_{i_{j + 1}}}^{\text{r}}\)//
total time budget left for repairing bikes} \\
\textbf{3} & \multicolumn{3}{l@{}}{%
\textbf{while} \(T^{\text{op}} < t^{\text{repair}}\) \textbf{do}} \\
& & \multicolumn{2}{l@{}}{%
// remove the station with the lowest benefit-cost ratio function value from the route} \\
\textbf{4} & & \multicolumn{2}{l@{}}{%
Remove station \(s' = \arg\min_{s \in S}BCRF_{s}^{\text{r}}\) from \(r_{i}\)} \\
\textbf{5} & & \multicolumn{2}{l@{}}{%
\(T^{\text{op}} = T - \sum\limits_{j = 0}^{|r_{i}{| - 1}}t_{s_{i_{j}},s_{i_{j + 1}}}^{\text{r}}\)} \\
\textbf{6} & \multicolumn{3}{l@{}}{%
\textbf{end while}} \\
& \multicolumn{3}{l@{}}{%
/* the set \(\mathit{unsat}\) stores the stations where broken bikes cannot be all repaired due to
insufficient station repairing budget */} \\
\textbf{7} & \multicolumn{3}{l@{}}{%
\(\mathit{unsat} = \varnothing\)} \\
& \multicolumn{3}{l@{}}{%
// record the overly assigned station repairing budget} \\
\textbf{8} & \multicolumn{3}{l@{}}{%
\(\mathit{extraTime} = 0\)} \\
& \multicolumn{3}{l@{}}{%
// lines 9 to 18: assign the maximum repairing quantity at each station
(\(1 \leq j \leq \left| r_{i} \right| - 1\))} \\
\textbf{9} & \multicolumn{3}{l@{}}{%
\textbf{for} \emph{j} \textbf{in} 1 to \(\left| r_{i} \right| - 2\) \textbf{do}} \\
& & \multicolumn{2}{l@{}}{%
// distribute the total station repairing budget using the ratio of the \emph{BCRF} value of each
station} \\
\textbf{10} & & \multicolumn{2}{l@{}}{%
\(\mathit{maxTim}e_{s_{i_{j}},i}^{\text{r}} = \frac{T^{\text{op}} \times BCRF_{s_{i_{j}}}^{\text{r}}\left( {p_{s_{i_{j}}},b_{s_{i_{j}}}} \right)}{\sum_{j' = 1}^{|r_{i}{| - 2}}{\mathit{BCR}F_{s_{i_{j'}}}^{\text{r}}\left( {p_{s_{i_{j'}}},b_{s_{i_{j'}}}} \right)}}\)} \\
\textbf{11} & & \multicolumn{2}{l@{}}{%
\(\mathit{repai}r_{s_{i_{j}}}^{\max} = \min\left\{ {\max\left\{ {q_{s_{i_{j}}} - p_{s_{i_{j}}},0} \right\},b_{s_{i_{j}}}} \right\}\)//
maximum repairing quantity} \\
\textbf{12} & & \multicolumn{2}{l@{}}{%
\textbf{if}
\(\mathit{maxTim}e_{s_{i_{j}}}^{\text{r}} < repair_{s_{i_{j}}}^{\max} \cdot t^{\text{repair}}\)
\textbf{then}} \\
Empty Cell & & & // record stations with insufficient repairing budgets to repair all broken
bikes \\
\textbf{13} & & & \(\mathit{unsat} = unsat \cup \left\{ s_{i_{j}} \right\}\) \\
\textbf{14} & & \multicolumn{2}{l@{}}{%
\textbf{else}} \\
Empty Cell & & & // gather the excess repairing budget and adjust the repairing budget \\
\textbf{15} & & &
\(\mathit{extraTime} = extraTime + maxTime_{s_{i_{j}},i}^{\text{r}} - repair_{s_{i_{j}}}^{\max} \cdot t^{\text{repair}}\) \\
\textbf{16} & & &
\(\mathit{maxTim}e_{s_{i_{j}}}^{\text{r}} = repair_{s_{i_{j}}}^{\max} \cdot t^{\text{repair}}\) \\
\textbf{17} & & \multicolumn{2}{l@{}}{%
\textbf{end if}} \\
\textbf{18} & \multicolumn{3}{l@{}}{%
\textbf{end for}} \\
& \multicolumn{3}{l@{}}{%
// lines 19 to 25: reassign the total excess time to nodes with insufficient budgets} \\
\textbf{19} & \multicolumn{3}{l@{}}{%
\textbf{while} \(\mathit{extraTime} > t^{\text{repair}}\) \textbf{and}
\(\mathit{unsat} \neq \varnothing\) \textbf{do}} \\
\textbf{20} & & \multicolumn{2}{l@{}}{%
\(\widehat{s} =\) The node in \(\mathit{unsat}\) with the highest \(\mathit{BCR}F_{s}^{\text{r}}\)
(break ties by choosing the smallest \(\widehat{s}\))} \\
\textbf{21} & & \multicolumn{2}{l@{}}{%
\emph{added}
=\(\min\left\{ {\mathit{extraTime},repair_{\widehat{s}}^{\max} \cdot t^{\text{repair}} - maxTime_{\widehat{s},i}^{\text{r}}} \right\}\)} \\
\textbf{22} & & \multicolumn{2}{l@{}}{%
\(\mathit{maxTim}e_{\widehat{s},i}^{\text{r}} = maxTime_{\widehat{s},i}^{\text{r}} + added\)} \\
\textbf{23} & & \multicolumn{2}{l@{}}{%
\(\mathit{extraTime} = extraTime - added\)} \\
\textbf{24} & & \multicolumn{2}{l@{}}{%
\(\mathit{unsat} = unsat \smallsetminus \left\{ \widehat{s} \right\}\)} \\
\textbf{25} & \multicolumn{3}{l@{}}{%
\textbf{end while}} \\
\textbf{26} & \multicolumn{3}{l@{}}{%
\textbf{return} \(\left( {\mathit{maxTim}e_{s,i}^{\text{r}}} \right)_{s \in S}\)} \\
\end{longtable}
}

The algorithm begins by computing the total time available for repairs \(T^{\text{op}}\) by
deducting travel time between stations from the total time budget \emph{T} (Line 2). It then
conducts a feasibility check and route pruning to ensure that \(T^{\text{op}}\) is sufficient for at
least one repair (Line 3 to Line 6). Then, the algorithm assigns the total time available to each
station based on \(\mathit{BCR}F_{s}^{\text{r}}\), and determines the maximum possible number of
repairs within the allocated time. Stations where the time is insufficient are added to
\(\mathit{unsat}\), while the excess time budget is added to \(\mathit{extraTime}\)(Line 7 to Line
18). If there is extra time left (i.e., \(\mathit{extraTime} > 0\)) and \(\mathit{unsat}\) is not
empty, the algorithm redistributes the extra time to stations in \(\mathit{unsat}\), prioritizing
those with the highest \(\mathit{BCR}F_{s}^{\text{r}}\) until all \(\mathit{extraTime}\) is
allocated or \(\mathit{unsat}\) is empty (Line 19 to Line 25). Finally, the algorithm returns the
station repairing budget for each station.

\protect\phantomsection\label{t0045}
{\def\LTcaptype{none} % do not increment counter
\begin{longtable}[]{@{}llll@{}}
\toprule\noalign{}
\multicolumn{4}{@{}l@{}}{%
\textbf{Algorithm 4:} \(allocateLoadingBudget\left( r_{k}' \right)\)} \\
\midrule\noalign{}
\endhead
\bottomrule\noalign{}
\endlastfoot
\textbf{1} & \multicolumn{3}{l@{}}{%
Initialization:\(\mathit{maxTim}e_{s,k}^{\text{t}} = 0,\forall s \in S\)} \\
\textbf{2} & \multicolumn{3}{l@{}}{%
\(T^{\text{op}} = T - \sum\limits_{l = 0}^{{|r_{k}'|} - 1}t_{s_{k_{l}}',s_{k_{l + 1}}'}^{\text{t}}\)//
The total time budget left for truck \emph{k} to load and unload bikes} \\
\textbf{3} & \multicolumn{3}{l@{}}{%
\textbf{while} \(T^{\text{op}} < 2t^{\text{load}}\) \textbf{do}} \\
& & \multicolumn{2}{l@{}}{%
// remove the station with the lowest benefit-cost ratio function value from the route} \\
\textbf{4} & & \multicolumn{2}{l@{}}{%
Remove station \(s' = \arg\min_{s \in S}BCRF_{s}^{\text{t}}\) from \(r_{k}'\)} \\
\textbf{5} & & \multicolumn{2}{l@{}}{%
\(T^{\text{op}} = T - \sum\limits_{l = 0}^{{|r_{k}'|} - 1}t_{s_{k_{l}}',s_{k_{l + 1}}'}^{\text{t}}\)} \\
\textbf{6} & \multicolumn{3}{l@{}}{%
\textbf{end while}} \\
\textbf{7} & \multicolumn{3}{l@{}}{%
\(\mathit{unsat} = \varnothing\)} \\
\textbf{8} & \multicolumn{3}{l@{}}{%
\(\mathit{extraTime} = 0\)} \\
\textbf{9} & \multicolumn{3}{l@{}}{%
\(\mathit{NS} = {\text{The set of stations in}\mspace{6mu}}r_{k}'\)} \\
\textbf{10} & \multicolumn{3}{l@{}}{%
\textbf{for} \emph{s} \textbf{in} \emph{NS} \textbf{do}} \\
\textbf{11} & & \multicolumn{2}{l@{}}{%
\(\mathit{maxTim}e_{s,k}^{\text{t}} = \frac{T^{\text{op}} \times BCRF_{s}^{\text{t}}\left( {p_{s},b_{s}} \right)}{\sum_{s' \in NS}{\mathit{BCR}F_{s'}^{\text{t}}\left( {p_{s'},b_{s'}} \right)}}\)} \\
\textbf{12} & & \multicolumn{2}{l@{}}{%
\(\mathit{loa}d_{s}^{\max} = \left| {p_{s} - q_{s}} \right| + b_{s}\)} \\
\textbf{13} & & \multicolumn{2}{l@{}}{%
\textbf{if} \(\mathit{maxTim}e_{s,k}^{\text{t}} < 2t^{\text{load}} \cdot load_{s}^{\max}\)
\textbf{then}} \\
\textbf{14} & & & \(\mathit{unsat} = unsat \cup \left\{ s \right\}\) \\
\textbf{15} & & \multicolumn{2}{l@{}}{%
\textbf{else}} \\
\textbf{16} & & &
\(\mathit{extraTime} = extraTime + maxTime_{s,k}^{\text{t}} - 2t^{\text{load}} \cdot load_{s}^{\max}\) \\
\textbf{17} & & & \(\mathit{maxTim}e_{s,k}^{\text{t}} = 2t^{\text{load}} \cdot load_{s}^{\max}\) \\
\textbf{18} & & \multicolumn{2}{l@{}}{%
\textbf{end if}} \\
\textbf{19} & \multicolumn{3}{l@{}}{%
\textbf{end for}} \\
\textbf{20} & \multicolumn{3}{l@{}}{%
\textbf{while} \(\mathit{extraTime} > 2t^{\text{load}}\) \textbf{and}
\(\mathit{unsat} \neq \varnothing\) \textbf{do}} \\
\textbf{21} & & \multicolumn{2}{l@{}}{%
\(\widehat{s} =\) The node in \(\mathit{unsat}\) with the highest
\(\mathit{BCR}F_{s}^{\text{t}}\)} \\
\textbf{22} & & \multicolumn{2}{l@{}}{%
\(\mathit{added} = \min\left\{ {\mathit{extraTime},2t^{\text{load}} \cdot load_{\widehat{s}}^{\max} - maxTime_{\widehat{s},k}^{\text{t}}} \right\}\)} \\
\textbf{23} & & \multicolumn{2}{l@{}}{%
\(\mathit{maxTim}e_{\widehat{s},k}^{\text{t}} = maxTime_{\widehat{s},k}^{\text{t}} + added\)} \\
\textbf{24} & & \multicolumn{2}{l@{}}{%
\(\mathit{extraTime} = extraTime - added\)} \\
\textbf{25} & & \multicolumn{2}{l@{}}{%
\(\mathit{unsat} = unsat \smallsetminus \left\{ \widehat{s} \right\}\)} \\
\textbf{26} & \multicolumn{3}{l@{}}{%
\textbf{end while}} \\
\textbf{27} & \multicolumn{3}{l@{}}{%
\textbf{return} \(\left( {\mathit{maxTim}e_{s,k}^{\text{t}}} \right)_{s \in S}\)} \\
\end{longtable}
}

Note that for a visited surplus station, we must consider that each bike loaded onto the truck also
needs to be unloaded at another node (either the depot or a deficit station). Consequently, we set
the station (un)loading budget as twice the loading time for each bike at this station. For a
visited deficit station, at the current stage, we assume that all bikes being unloaded are
additionally loaded from the depot. Therefore, we also set the station (un)loading budget as twice
the loading time for each bike at this station. It is obvious that this approach leads to an
overestimation, as it double-counts the time for bikes that are simply being transferred from
surplus to deficit stations. To correct this, we make appropriate adjustments to the station
(un)loading budgets during the detailed station-by-station assignment process, as outlined in
\textbf{Algorithms 6, 7}, and \textbf{8} in subsequent sections.

\protect\phantomsection\label{s0160}
\subsection{Appendix D. Repairing, loading, and unloading quantity assignment
algorithms}\label{st190}

\protect\phantomsection\label{p0795}
We start by introducing the repairing quantity assignment. The algorithm takes the route of
repairman \emph{i} and the station repairing budget as the input and returns the number of broken
bikes repaired by repairman \emph{i} at each station. For surplus stations, no bike is repaired at
the current stage. For deficit stations, the repairing quantity is constrained by the deviation of
usable bike inventory from the target level and the station repairing budget.

\protect\phantomsection\label{t0050}
{\def\LTcaptype{none} % do not increment counter
\begin{longtable}[]{@{}lll@{}}
\toprule\noalign{}
\multicolumn{3}{@{}l@{}}{%
\textbf{Algorithm 5:}
\(assignRepair\left( {r_{i},\left( {\mathit{maxTim}e_{s,i}^{\text{r}}} \right)_{s \in S}} \right)\)} \\
\midrule\noalign{}
\endhead
\bottomrule\noalign{}
\endlastfoot
\textbf{1} & \multicolumn{2}{l@{}}{%
Initialization:\(g_{s,i} = 0,\forall s \in S\)} \\
\textbf{2} & \multicolumn{2}{l@{}}{%
\textbf{for} \emph{j}~\textbf{=}~1 to \(\left| r_{i} \right| - 2\) \textbf{do}} \\
\textbf{3} & &
\(g_{s_{i_{j}},i} = \min\left\{ {\max\left\{ {q_{s_{\mathit{ij}}} - p_{s_{\mathit{ij}}},0} \right\},b_{s_{i_{j}}},\left\lfloor \mathit{maxTim}e_{s_{i_{j}},i}^{\text{r}}/t^{\text{repair}} \right\rfloor} \right\}\) \\
\textbf{4} & \multicolumn{2}{l@{}}{%
\textbf{end for}} \\
\textbf{5} & \multicolumn{2}{l@{}}{%
\textbf{return} \(\left( g_{s,i} \right)_{s \in S}\)} \\
\end{longtable}
}

Next, we present the algorithms for (un)loading quantity assignment of trucks, including usable bike
loading (\textbf{Algorithm 6}), usable bike unloading (\textbf{Algorithms 7 and 8}), and broken bike
loading (\textbf{Algorithm 9}). Note that we use \emph{l} to represent the order index for defining
when the concerned node is visited by truck and assigned a loading quantity in these four
algorithms.

\protect\phantomsection\label{t0055}
{\def\LTcaptype{none} % do not increment counter
\begin{longtable}[]{@{}ll@{}}
\toprule\noalign{}
\multicolumn{2}{@{}l@{}}{%
\textbf{Algorithm 6:}
\(assignLoadU\left( {r_{k}',l,\left( {\mathit{maxTimeRemai}n_{s,k}^{\text{t}}} \right)_{s \in S},\left( {\mathit{unsa}t_{s,k}^{\text{load,u}}} \right)_{s \in S}} \right)\)} \\
\midrule\noalign{}
\endhead
\bottomrule\noalign{}
\endlastfoot
\textbf{1} & \(s = s_{k_{l}}'\) \\
\textbf{2} &
\(n = \min\left\{ \left. {Q_{k} - \psi_{l - 1,k}^{\text{u}} - \psi_{l - 1,k}^{\text{b}},p_{s} - q_{s}} \right\} \right.\)//
the quantity without the station loading budget \\
\textbf{3} &
\(y_{l,k}^{\text{u}^{-}} = \min\left\{ \left. {n,\left\lfloor \mathit{maxTimeRemai}n_{s,k}^{\text{t}}/2t^{\text{load}} \right\rfloor} \right\} \right.\)//
the quantity with the station loading budget \\
& // Record the number of bikes not loaded due to the limited station loading budget \\
\textbf{4} &
\(\mathit{unsa}t_{s,k}^{\text{load,u}} = unsat_{s,k}^{\text{load,u}} + \left( {n - y_{l,k}^{\text{u}^{-}}} \right)\) \\
\textbf{5} &
\(\mathit{maxTimeRemai}n_{s,k}^{\text{t}} = maxTimeRemain_{s,k}^{\text{t}} - 2t^{\text{load}} \cdot y_{l,k}^{\text{u}^{-}}\) \\
\textbf{6} & \textbf{return}
\(\left( {\mathit{maxTimeRemai}n_{s,k}^{\text{t}}} \right)_{s \in S}\),\(\left( {\mathit{unsa}t_{s,k}^{\text{load,u}}} \right)_{s \in S}\),\(y_{l,k}^{\text{u}^{-}}\) \\
\end{longtable}
}

\textbf{Algorithm 6} calculates the loading quantity of usable bikes for truck \emph{k} at a given
node. The number of bikes loaded is constrained by the truck\textquotesingle s capacity, the
station\textquotesingle s surplus inventory, and the station loading budget (Line 3). The algorithm
records the unsatisfied loading quantity due to the limited station loading budget (Line 4) and
updates the station loading budget (Line 5). Finally, it returns the updated remaining station
(un)loading budget, unsatisfied loading quantity due to the limited station loading budget, and the
number of bikes loaded onto truck \emph{k} (Line 6).

\protect\phantomsection\label{t0060}
{\def\LTcaptype{none} % do not increment counter
\begin{longtable}[]{@{}lllll@{}}
\toprule\noalign{}
\multicolumn{5}{@{}l@{}}{%
\textbf{Algorithm 7:} \(assignUnloadU\begin{pmatrix}
{r_{k}',l,\left( {\mathit{maxTimeRemai}n_{s,k}^{\text{t}}} \right)_{s \in S},} \\
{\left( {\mathit{unsa}t_{s,k}^{\text{load,u}}} \right)_{s \in S},\left( {\mathit{unsa}t_{s,k}^{\text{unload,u}}} \right)_{s \in S}}
\end{pmatrix}\)} \\
\midrule\noalign{}
\endhead
\bottomrule\noalign{}
\endlastfoot
\textbf{1} & \multicolumn{4}{l@{}}{%
\(s = s_{k_{l}}'\)} \\
\textbf{2} & \multicolumn{4}{l@{}}{%
\(y_{l,k}^{\text{u}^{+}} = \min\left\{ {\psi_{l - 1,k}^{\text{u}},C_{s} - p_{s} - b_{s},q_{s} - p_{s}} \right\}\)} \\
& \multicolumn{4}{l@{}}{%
/* add back the double-counted loading time as these unloaded bikes are not from the depot but from
the previously visited surplus stations */} \\
\textbf{3} & \multicolumn{4}{l@{}}{%
\(\mathit{maxTimeRemai}n_{s,k}^{\text{t}} = maxTimeRemain_{s,k}^{\text{t}} + 2t^{\text{load}} \cdot y_{l,k}^{\text{u}^{+}}\)} \\
& \multicolumn{4}{l@{}}{%
// target inventory not reached due to insufficient bike supply on the truck} \\
\textbf{4} & \multicolumn{4}{l@{}}{%
\textbf{if} \(y_{l,k}^{\text{u}^{+}} < q_{s} - p_{s}\) \emph{and}
\(y_{l,k}^{\text{u}^{+}} \neq C_{s} - p_{s} - b_{s}\) \textbf{then}} \\
& & \multicolumn{3}{l@{}}{%
/* line 5 to 20: load more bikes from the visited but still surplus stations and unload them to the
current station until the target inventory is reached or the station loading budget is used up
*/} \\
\textbf{5} & & \multicolumn{3}{l@{}}{%
\textbf{while} \(\exists\widehat{s} \in S:unsat_{\widehat{s},k}^{\text{load},\text{u}} > 0\)
\textbf{and} \(\mathit{maxTimeRemai}n_{s,k}^{\text{t}} \geq 2t^{\text{load}}\) \textbf{do}} \\
\textbf{6} & & & \multicolumn{2}{l@{}}{%
\(s'\)~=~the station with the highest \(\mathit{BCR}F_{\widehat{s}}^{\text{t}}\) and
\(\mathit{unsat}_{s^{'},k}^{\text{load},\text{u}} > 0\) (break ties by choosing the smallest
\(s'\))} \\
\textbf{7} & & & \multicolumn{2}{l@{}}{%
\(l'\)~=~the index of station \(s'\) in the route \(r_{k}'\) (break ties by choosing the one with
the largest order index)} \\
& & & \multicolumn{2}{l@{}}{%
// minimum truck residual capacity from the chosen node to the current node} \\
\textbf{8} & & & \multicolumn{2}{l@{}}{%
\(rc_{k} = \min_{l^{''} = l^{'},l^{'} + 1,...,l - 1}Q_{k} - \psi_{l^{''},k}^{\text{u}} - \psi_{l^{''},k}^{\text{b}}\)} \\
& & & \multicolumn{2}{l@{}}{%
// the extra loaded quantity with no station loading budget} \\
\textbf{9} & & & \multicolumn{2}{l@{}}{%
\(n = \min\left( {q_{s} - p_{s} - y_{l,k}^{\text{u}^{+}},rc_{k},unsat_{s^{'},k}^{\text{load,u}}} \right)\)} \\
& & & \multicolumn{2}{l@{}}{%
// the extra loaded quantity with the remaining station loading budget} \\
\textbf{10} & & & \multicolumn{2}{l@{}}{%
\(\Lambda^{\text{load,u}} = \min\left\{ {n,\left\lfloor \mathit{maxTimeRemai}n_{s,k}^{\text{t}}/2t^{\text{load}} \right\rfloor} \right\}\)} \\
\textbf{11} & & & \multicolumn{2}{l@{}}{%
\(\mathit{unsa}t_{s,k}^{\text{unload},\text{u}} = unsat_{s,k}^{\text{unload},\text{u}} + \left( {n - \Lambda^{\text{load,u}}} \right)\)//
update the unsatisfied unloading quantity} \\
& & & \multicolumn{2}{l@{}}{%
// update the remaining station loading budget of the current station} \\
\textbf{12} & & & \multicolumn{2}{l@{}}{%
\(\mathit{maxTimeRemai}n_{s,k}^{\text{t}} = maxTimeRemain_{s,k}^{\text{t}} - 2t^{\text{load}} \cdot \Lambda^{\text{load,u}}\)} \\
\textbf{13} & & & \multicolumn{2}{l@{}}{%
\(\mathit{unsa}t_{s^{'},k}^{\text{load,u}} = unsat_{s^{'},k}^{\text{load,u}} - \Lambda^{\text{load},\text{u}}\)//
update the unsatisfied loading quantity} \\
\textbf{14} & & & \multicolumn{2}{l@{}}{%
\(p_{s'} = p_{s'} - \Lambda^{\text{load},\text{u}}\)// update the station inventory after extra
loading of bikes} \\
\textbf{15} & & & \multicolumn{2}{l@{}}{%
\textbf{for} \(l''\) \textbf{in} \(l'\) \emph{to} \(l - 1\) \textbf{do} // update the truck
inventories} \\
\textbf{16} & & & &
\(\psi_{l^{''},k}^{\text{u}} = \psi_{l^{''},k}^{\text{u}} + \Lambda^{\text{load,u}}\) \\
\textbf{17} & & & \multicolumn{2}{l@{}}{%
\textbf{end for}} \\
\textbf{18} & & & \multicolumn{2}{l@{}}{%
\(y_{l^{'},k}^{\text{u}^{-}} = y_{l^{'},k}^{\text{u}^{-}} + \Lambda^{\text{load,u}}\)// update the
loading quantity} \\
\textbf{19} & & & \multicolumn{2}{l@{}}{%
\(y_{l,k}^{\text{u}^{+}} = y_{l,k}^{\text{u}^{+}} + \Lambda^{\text{load,u}}\)// update the unloading
quantity} \\
\textbf{20} & & \multicolumn{3}{l@{}}{%
\textbf{end while}} \\
\textbf{21} & & \multicolumn{3}{l@{}}{%
\textbf{if} \(y_{l,k}^{\text{u}^{+}} < q_{s} - p_{s}\) \textbf{then}// the target inventory is still
not reached} \\
\textbf{22} & & & \multicolumn{2}{l@{}}{%
\emph{dv =} \(\max\left\{ {\mathit{dv} \mid dv < l,s_{k_{\mathit{dv}}}' = 0} \right\}\) // order
index of the last visited depot} \\
\textbf{23} & & & \multicolumn{2}{l@{}}{%
// adjust the loading at the previously visited depot (\textbf{Algorithm 8})} \\
\textbf{24} & & & \multicolumn{2}{l@{}}{%
\emph{h}, \(\mathit{maxTimeRemai}n_{s,k}^{\text{t}}\),
\(\mathit{unsa}t_{s,k}^{\text{unload},\text{u}}\)\emph{=}\(adjustDepotLoading\left( {\mathit{dv},l,maxTimeRemain_{s,k}^{\text{t}},unsat_{s,k}^{\text{unload,u}}} \right)\)} \\
\textbf{25} & & & \multicolumn{2}{l@{}}{%
\(y_{l,k}^{\text{u}^{+}} = y_{l,k}^{\text{u}^{+}} + h\)} \\
\textbf{26} & & \multicolumn{3}{l@{}}{%
\textbf{end if}} \\
\textbf{27} & \multicolumn{4}{l@{}}{%
\textbf{end if}} \\
\textbf{28} & \multicolumn{4}{l@{}}{%
\textbf{return} \(\left( {\mathit{maxTimeRemai}n_{s,k}^{\text{t}}} \right)_{s \in S}\),
\(\left( {\mathit{unsa}t_{s,k}^{\text{load,u}}} \right)_{s \in S}\),
\(\left( {\mathit{unsa}t_{s,k}^{\text{unload,u}}} \right)_{s \in S}\),\(y_{l,k}^{\text{u}^{+}}\)} \\
\end{longtable}
}

\textbf{Algorithm 7} determines the unloading quantity by considering the truck's usable bike
inventory, the station's residual capacity, and the deviation of the inventory from the target
inventory levels (Line 2). If the usable bike shortage cannot be satisfied due to insufficient
usable bikes on the truck, the algorithm utilizes the unsatisfied loading quantities from previously
visited stations that are still in a surplus state and tries to load more bikes at those stations to
satisfy the shortage at the current station (Lines 5 to Line 20). Note that during this process, we
also capture the number of usable bikes unable to be loaded from other nodes by truck k due to the
insufficient station loading budget (Line 9 to Line 11). If the station's demand remains unmet, a
final adjustment is made on the previous depot visit (Line 21 to Line 26) by reducing the number of
usable bikes unloaded or loading more usable bikes there. The detailed procedure for adjustment at
the depot is presented in \textbf{Algorithm 8}. Note that we use dv to denote the order index for
the last visited depot in \(r_{k}'\) that requires (un)loading adjustment, h to denote the increased
loading of usable bikes at the station that requires (un)loading adjustment, and s is the station
that requires (un)loading adjustment from the previous depot.

\protect\phantomsection\label{t0065}
{\def\LTcaptype{none} % do not increment counter
\begin{longtable}[]{@{}lll@{}}
\toprule\noalign{}
\multicolumn{3}{@{}l@{}}{%
\textbf{Algorithm 8:}
\(adjustDepotLoading\left( {\mathit{dv},l,maxTimeRemain_{s,k}^{\text{t}},unsat_{s,k}^{\text{unload,u}}} \right)\)} \\
\midrule\noalign{}
\endhead
\bottomrule\noalign{}
\endlastfoot
\textbf{1} & \multicolumn{2}{l@{}}{%
\(rc_{k} = \min_{l^{'} = dv,dv + 1,...,l - 1}Q_{k} - \psi_{l^{'},k}^{\text{u}} - \psi_{l^{'},k}^{\text{b}}\)} \\
\textbf{2} & \multicolumn{2}{l@{}}{%
\(n = \min\left\{ {rc_{k},q_{s} - p_{s} - y_{l,k}^{\text{u}^{+}}} \right\}\)} \\
\textbf{3} & \multicolumn{2}{l@{}}{%
\(h = \min\left\{ {n,\left\lfloor \mathit{maxTimeRemai}n_{s,k}^{\text{t}}/2t^{\text{load}} \right\rfloor} \right\}\)} \\
\textbf{4} & \multicolumn{2}{l@{}}{%
\(\mathit{unsa}t_{s_{k_{l}}',k}^{\text{unload},\text{u}} = unsat_{s_{k_{l}}',k}^{\text{unload},\text{u}} + \left( {n - h} \right)\)} \\
\textbf{5} & \multicolumn{2}{l@{}}{%
\textbf{if} \(y_{\mathit{dv},k}^{\text{u}^{+}} > 0\) \textbf{then}// the original operation at the
depot is unloading} \\
\textbf{6} & &
\(y_{\mathit{dv},k}^{\text{u}^{+}} = \max\left\{ {y_{\mathit{dv},k}^{\text{u}^{+}} - h,0} \right\}\)
// reduce the unloading of usable bikes at the depot \\
& & // if the reduction exceeds the unloading quantity, the final decision at the depot becomes
loading \\
\textbf{7} & &
\(y_{\mathit{dv},k}^{\text{u}^{-}} = \max\left\{ {h - y_{\mathit{dv},k}^{\text{u}^{+}},0} \right\}\) \\
\textbf{8} & &
\(\mathit{maxTimeRemai}n_{s_{k_{l}}',k}^{\text{t}} = maxTimeRemain_{s_{k_{l}}',k}^{\text{t}} - 2t^{\text{load}} \cdot y_{\mathit{dv},k}^{\text{u}^{-}}\) \\
\textbf{9} & \multicolumn{2}{l@{}}{%
\textbf{end if}} \\
\textbf{10} & \multicolumn{2}{l@{}}{%
\textbf{if} \(y_{\mathit{dv},k}^{\text{u}^{-}} > 0\) \textbf{then}// the original operation at the
depot is loading} \\
\textbf{11} & & \(y_{\mathit{dv},k}^{\text{u}^{-}} = y_{\mathit{dv},k}^{\text{u}^{-}} + h\)//
increase the loading quantity \\
\textbf{12} & &
\(\mathit{maxTimeRemai}n_{s,k}^{\text{t}} = maxTimeRemain_{s,k}^{\text{t}} - 2t^{\text{load}} \cdot h\) \\
\textbf{13} & \multicolumn{2}{l@{}}{%
\textbf{end if}} \\
\textbf{14} & \multicolumn{2}{l@{}}{%
\textbf{for} \(l'\)~=~\emph{dv} to \emph{l} − 1 do} \\
\textbf{15} & & \(\psi_{l',k}^{\text{u}} = \psi_{l',k}^{\text{u}} + h\) \\
\textbf{16} & \multicolumn{2}{l@{}}{%
\textbf{end for}} \\
\textbf{17} & \multicolumn{2}{l@{}}{%
\textbf{return} \emph{h},
\(\mathit{maxTimeRemai}n_{s,k}^{\text{t}}\),\(\mathit{unsa}t_{s,k}^{\text{unload},\text{u}}\)} \\
\end{longtable}
}

\textbf{Algorithm 8} initiates by calculating the minimum residual capacity of the truck from the
last visited depot to the upcoming station \emph{s} =\(s_{k_{l}}'\) (Line 1). It then computes the
additional number of bikes that need to be loaded at the depot, which is the lowest value among the
truck's minimum residual capacity, the shortage of usable bikes at the upcoming station, and the
number of additional bikes that can be loaded using the station loading budget of the upcoming
deficit station (Lines 2--3). If the initial operation at the depot is to unload bikes, the
algorithm decreases the unloading quantity by the newly calculated amount. If the adjustment leads
to a surplus, the operation at the depot is switched from unloading to loading. Conversely, if the
depot was initially set to load bikes, the algorithm simply increases the loading quantity by the
computed adjustment (Line 5 to Line 13). The algorithm terminates by updating the truck inventory
along the route from the depot to the upcoming deficit station and returns the adjustment value.
During the process, the unsatisfied unloading demand is limited by station unloading budgets and the
budgets are updated accordingly (Line 4 and Line 12).

\protect\phantomsection\label{p0800}
After usable bikes are loaded or unloaded, we continue to assign the quantities of broken bikes to
be loaded using \textbf{Algorithm 9}. The assignment follows a similar procedure as
\textbf{Algorithm 6}, with the loading quantities of broken bikes constrained by the residual
capacity of the truck, the broken bike inventory at the station, and the station loading budget.

\protect\phantomsection\label{t0070}
{\def\LTcaptype{none} % do not increment counter
\begin{longtable}[]{@{}ll@{}}
\toprule\noalign{}
\multicolumn{2}{@{}l@{}}{%
\textbf{Algorithm 9:}
\(assignLoadB\left( {r_{k}',l,\left( {\mathit{maxTimeRemai}n_{s,k}^{\text{t}}} \right)_{s \in S},\left( {\mathit{unsa}t_{s,k}^{\text{load},\text{b}}} \right)_{s \in S}} \right)\)} \\
\midrule\noalign{}
\endhead
\bottomrule\noalign{}
\endlastfoot
\textbf{1} & \(s = s_{k_{l}}'\) \\
\textbf{2} &
\(n = \min\left\{ \left. {b_{s},Q_{k} - \psi_{l - 1,k}^{\text{u}} - \psi_{l - 1,k}^{\text{b}}} \right\} \right.\) \\
\textbf{3} &
\(y_{l,k}^{\text{b}^{-}} = \min\left\{ \left. {n,\left\lfloor \mathit{maxTimeRemai}n_{s,k}^{\text{t}}/2t^{\text{load}} \right\rfloor} \right\} \right.\) \\
\textbf{4} & // The number of broken bikes not loaded due to the limited station loading budget \\
&
\(\mathit{unsa}t_{s,k}^{\text{load,b}} = unsat_{s,k}^{\text{load,b}} + \left( {n - y_{l,k}^{\text{b}^{-}}} \right)\) \\
\textbf{5} &
\(\mathit{maxTimeRemai}n_{s,k}^{\text{t}} = maxTimeRemain_{s,k}^{\text{t}} - 2t^{\text{load}} \cdot y_{l,k}^{\text{b}^{-}}\) \\
\textbf{6} & \textbf{return} \(\left( {\mathit{maxTimeRemai}n_{s,k}^{\text{t}}} \right)_{s \in S}\),
\(\left( {\mathit{unsa}t_{s,k}^{\text{load},\text{b}}} \right)_{s \in S}\),\(y_{l,k}^{\text{b}^{-}}\) \\
\end{longtable}
}

\protect\phantomsection\label{s0165}
\subsection{Appendix E. (Un)loading budget adjustment algorithm for trucks}\label{st195}

\protect\phantomsection\label{p0805}
\textbf{Algorithm 10} shows the procedure of the (un)loading budget adjustment.

\protect\phantomsection\label{t0075}
{\def\LTcaptype{none} % do not increment counter
\begin{longtable}[]{@{}lll@{}}
\toprule\noalign{}
\multicolumn{3}{@{}l@{}}{%
\textbf{Algorithm 10:} \(adjustLoadingBudget\begin{pmatrix}
{\left( {\mathit{maxTimeIni}t_{s,k}^{\text{t}}} \right)_{s \in S},\left( {\mathit{maxTimeRemai}n_{s,k}^{\text{t}}} \right)_{s \in S},} \\
{\left( {\mathit{unsa}t_{s,k}^{\text{load},\text{u}}} \right)_{s \in S},\left( {\mathit{unsa}t_{s,k}^{\text{unload},\text{u}}} \right)_{s \in S},\left( {\mathit{unsa}t_{s,k}^{\text{load},\text{b}}} \right)_{s \in S}}
\end{pmatrix}\)} \\
\midrule\noalign{}
\endhead
\bottomrule\noalign{}
\endlastfoot
\textbf{1} & \multicolumn{2}{l@{}}{%
\(\mathit{unusedTim}e_{k} = 0\)// the variable that stores the time budget that is assigned but not
used} \\
& \multicolumn{2}{l@{}}{%
/* lines 2 to 6: add up all the bikes not loaded/unloaded at each station due to insufficient
station (un)loading budgets and store the result in \(\mathit{unsa}t_{\widehat{s},k}\), subtract the
overly assigned station (un)loading budgets, and accumulate them into
\(\mathit{unusedTim}e_{k}\)*/} \\
\textbf{2} & \multicolumn{2}{l@{}}{%
\textbf{for} \(\widehat{s}\) \textbf{in} \emph{S} \textbf{do}} \\
\textbf{3} & &
\(\mathit{unsa}t_{\widehat{s},k} = unsat_{\widehat{s},k}^{\text{load,u}\mspace{6mu}} + unsat_{\widehat{s},k}^{\text{load,b}\mspace{6mu}} + unsat_{\widehat{s},k}^{\text{unload,u}\mspace{6mu}}\) \\
\textbf{4} & &
\(\mathit{maxTimeNe}w_{\widehat{s},k}^{\text{t}} = maxTimeInit_{\widehat{s},k}^{\text{t}} - maxTimeRemain_{\widehat{s},k}^{\text{t}}\) \\
\textbf{5} & &
\(\mathit{unusedTim}e_{k} = unusedTime_{k} + maxTimeRemain_{\widehat{s},k}^{\text{t}}\) \\
\textbf{6} & \multicolumn{2}{l@{}}{%
\textbf{end for}} \\
\textbf{7} & \multicolumn{2}{l@{}}{%
\textbf{while} \(\mathit{unusedTim}e_{k} \geq 2t^{\text{load}}\) \textbf{and}
\(\exists\widehat{s} \in S:unsat_{\widehat{s},k} > 0\) \textbf{do}} \\
\textbf{8} & & \(s' =\) The station that satisfies \(\mathit{unsa}t_{\widehat{s},k} > 0\) with the
highest \(\mathit{BCR}F_{\widehat{s}}^{\text{t}}\) (break ties by choosing the smallest
\(\widehat{s}\)) \\
\textbf{9} & &
\(\Delta = \min\left\{ {\mathit{unusedTim}e_{k},2t^{\text{load}} \cdot unsat_{s',k}} \right\}\) \\
\textbf{10} & & \(\mathit{maxTimeNe}w_{s',k}^{\text{t}} = maxTimeNew_{s',k}^{\text{t}} + \Delta\) \\
\textbf{11} & & \(\mathit{unusedTim}e_{k} = unusedTime_{k} - \Delta\) \\
\textbf{12} & &
\(\mathit{unsa}t_{s',k} = unsat_{s',k} - \left\lfloor \Delta/2t^{\text{load}} \right\rfloor\) \\
\textbf{13} & \multicolumn{2}{l@{}}{%
\textbf{end while}} \\
\textbf{14} & \multicolumn{2}{l@{}}{%
\textbf{return} \(\left( {\mathit{maxTimeNe}w_{s,k}^{\text{t}}} \right)_{s \in S}\)} \\
\end{longtable}
}

\textbf{Algorithm 10} begins with a traversal of each station in the network, counting the total
number of bikes not (un)loaded at each station due to insufficient (un)loading budgets, reducing the
overly assigned budget from the initial station (un)loading budget, and sum up these unused times
(Line 1 to Line 6). The algorithm then enters a loop where it redistributes unused time budget to
stations with unsatisfied loading or unloading needs, i.e., the bikes that could not be loaded or
unloaded due to limited station (un)loading budgets (Line 7 to Line 13). Within each iteration, the
algorithm selects the station with the highest \(\mathit{BCR}F_{\widehat{s}}^{\text{t}}\) and
allocates additional time from the unused budget to make up for this insufficiency (Line 8 to Line
12). This process is repeated until there is either no more unused time or all the demands have been
met. Finally, the updated station (un)loading budgets are returned (Line 14).

\protect\phantomsection\label{da005}
\subsection{Data availability}\label{st020}

\protect\phantomsection\label{p0030}
Data will be made available on request.

\protect\phantomsection\label{bi005}
\subsection{References}\label{references}

\protect\phantomsection\label{bs005}
\begin{enumerate}
\tightlist
\item
  Alvarez-Valdes et al., 2016

  R. Alvarez-Valdes, J.M. Belenguer, E. Benavent, J.D. Bermudez, F. Muñoz, E. Vercher, F. Verdejo

  Optimizing the level of service quality of a bike-sharing system

  Omega, 62 (2016), pp. 163-175
\item
  Brinkmann et al., 2019

  J. Brinkmann, M.W. Ulmer, D.C. Mattfeld

  Dynamic lookahead policies for stochastic-dynamic inventory routing in bike sharing systems

  Comput. Oper. Res., 106 (2019), pp. 260-279
\item
  Caggiani et al., 2018

  L. Caggiani, R. Camporeale, M. Ottomanelli, W.Y. Szeto

  A modeling framework for the dynamic management of free-floating bike-sharing systems

  Transp. Res. Part C: Emerg. Technol., 87 (2018), pp. 159-182
\item
  Caggiani and Ottomanelli, 2012

  L. Caggiani, M. Ottomanelli

  A modular soft computing based method for vehicles repositioning in bike-sharing systems

  Procedia-Soc. Behav. Sci., 54 (2012), pp. 675-684
\item
  Caggiani and Ottomanelli, 2013

  L. Caggiani, M. Ottomanelli

  A dynamic simulation based model for optimal fleet repositioning in bike-sharing systems

  Procedia-Soc. Behav. Sci., 87 (2013), pp. 203-210
\item
  Cai et al., 2022

  Y. Cai, G.P. Ong, Q. Meng

  Dynamic bicycle relocation problem with broken bicycles

  Transp. Res. Part E: Logist. Transp. Rev., 165 (2022), Article 102877
\item
  Castro Fernández, 2011

  A. Castro Fernández

  The contribution of bike-sharing to sustainable mobility in Europe

  Vienna University of Technology (2011)

  Ph.D. Thesis
\item
  Chang et al., 2018

  S. Chang, R. Song, S. He, G. Qiu

  Innovative bike-sharing in China: Solving faulty bike-sharing recycling problem

  J. Adv. Transp. (2018), 10.1155/2018/4941029
\item
  Chang et al., 2021

  X. Chang, J. Wu, H. Sun, G. Correia, J. Chen

  Relocating operational and damaged bikes in free-floating systems: A data-driven modeling
  framework for level of service enhancement

  Transp. Res. Part A: Policy Pract., 153 (2021), pp. 235-260
\item
  Chang et al., 2023

  X. Chang, J. Wu, H. Sun, X. Yan

  A smart predict-then-optimize method for dynamic green bike relocation in the free-floating system

  Transp. Res. Part C: Emerg. Technol., 153 (2023), Article 104220
\item
  Chen and Szeto, 2023

  M. Chen, W.Y. Szeto

  A static green bike repositioning problem with heavy and light carriers

  Transp. Res. Part D: Transp. Environ., 118 (2023), Article 103711
\item
  Chen et al., 2020

  M. Chen, D. Wang, Y. Sun, E.O.D. Waygood, W. Yang

  A comparison of users' characteristics between station-based bikesharing system and free-floating
  bikesharing system: case study in Hangzhou

  China. Transportation, 47 (2) (2020), pp. 689-704
\item
  Chen et al., 2024

  Q. Chen, S. Ma, H. Li, N. Zhu, Q.-C. He

  Optimizing bike rebalancing strategies in free-floating bike-sharing systems: An enhanced
  distributionally robust approach

  Transp. Res. Part E: Logist. Transp. Rev., 184 (2024), Article 103477
\item
  Cheng et al., 2021

  Y. Cheng, J. Wang, Y. Wang

  A user-based bike rebalancing strategy for free-floating bike sharing systems: A bidding model

  Transp. Res. Part E: Logist. Transp. Rev., 154 (2021), Article 102438
\item
  Chiariotti et al., 2020

  F. Chiariotti, C. Pielli, A. Zanella, M. Zorzi

  A bike-sharing optimization framework combining dynamic rebalancing and user incentives

  ACM Trans. Auton. Adapt. Syst., 14 (3) (2020), pp. 1-30
\item
  Citi Bike, 2024

  Citi Bike, 2024. June 2024 monthly report. Available at:
  https://mot-marketing-whitelabel-prod.s3.amazonaws.com/nyc/June-2024-Citi-Bike-Monthly-Report.pdf
  (Accessed: 8 September 2024).
\item
  Contardo et al., 2012

  C. Contardo, C. Morency, L.-M. Rousseau

  Balancing a dynamic public bike-sharing system

  Technical Report CIRRELT-2012-09. Montréal. (2012)
\item
  Dell'Amico et al., 2016

  M. Dell'Amico, M. Iori, S. Novellani, T. Stützle

  A destroy and repair algorithm for the bike sharing rebalancing problem

  Comput. Oper. Res., 71 (2016), pp. 149-162
\item
  Dell'Amico et al., 2018

  M. Dell'Amico, M. Iori, S. Novellani, A. Subramanian

  The bike sharing rebalancing problem with stochastic demands

  Transp. Res. Part b: Methodol., 118 (2018), pp. 362-380
\item
  Di Gaspero et al., 2013a

  Di Gaspero, L., Rendl, A., Urli, T., 2013a. A hybrid ACO+CP for balancing bicycle sharing systems.
  In: M.J. Blesa., C. Blum., P. Festa., A. Roli., M. Sampels. (Eds.), Hybrid Metaheuristics. Berlin,
  Heidelberg: Springer Berlin Heidelberg, pp. 198--212.
\item
  Di Gaspero et al., 2013b

  Di Gaspero, L., Rendl, A., Urli, T., 2013b. Constraint-based approaches for balancing bike sharing
  systems. In: C. Schulte. (Ed.), Principles and Practice of Constraint Programming. Berlin,
  Heidelberg: Springer Berlin Heidelberg, pp. 758--773.
\item
  Du et al., 2020

  M. Du, L. Cheng, X. Li, F. Tang

  Static rebalancing optimization with considering the collection of malfunctioning bikes in
  free-floating bike sharing system

  Transp. Res. Part E: Logist. Transp. Rev., 141 (2020), Article 102012
\item
  Fang et al., 2019

  W. Fang, X. Yang, W. Zhao, B. Liu

  Stability analysis of articulated dump truck based on eigenvalue method

  IOP Conf. Ser. Mater. Sci. Eng., 688 (2) (2019), Article 022024
\item
  Fishman et al., 2014

  E. Fishman, S. Washington, N. Haworth

  Bike share's impact on car use: Evidence from the United States, Great Britain, and Australia

  Transp. Res. Part d: Transp. Environ., 31 (2014), pp. 13-20
\item
  Forma et al., 2015

  I.A. Forma, T. Raviv, M. Tzur

  A 3-step math heuristic for the static repositioning problem in bike-sharing systems

  Transp. Res. Part B: Methodol., 71 (2015), pp. 230-247
\item
  Franceschetti et al., 2013

  A. Franceschetti, D. Honhon, T. Van Woensel, T. Bektaş, G. Laporte

  The time-dependent pollution-routing problem

  Transp. Res. Part B: Methodol., 56 (2013), pp. 265-293
\item
  Freund, 2018

  D. Freund

  Models and algorithms for transportation in the sharing economy

  Cornell University (2018)

  Ph.D. Thesis
\item
  Guo et al., 2024

  Y. Guo, J. Li, L. Xiao, H. Allaoui, A. Choudhary, L. Zhang

  Efficient inventory routing for bike-sharing systems: A combinatorial reinforcement learning
  framework

  Transp. Res. Part E: Logist. Transp. Rev., 182 (2024), Article 103415
\item
  Gurobi Optimization, 2023

  Gurobi Optimization, LLC, 2023. Gurobi optimizer reference manual. Available at:
  https://www.gurobi.com.
\item
  Ho and Szeto, 2014

  S.C. Ho, W.Y. Szeto

  Solving a static repositioning problem in bike-sharing systems using iterated tabu search

  Transp. Res. Part E: Logist. Transp. Rev., 69 (2014), pp. 180-198
\item
  Ho and Szeto, 2017

  S.C. Ho, W.Y. Szeto

  A hybrid large neighborhood search for the static multi-vehicle bike-repositioning problem

  Transp. Res. Part B: Methodol., 95 (2017), pp. 340-363
\item
  Hu et al., 2021

  R. Hu, Z. Zhang, X. Ma, Y. Jin

  Dynamic rebalancing optimization for bike-sharing system using priority-based MOEA/D algorithm

  IEEE Access, 9 (2021), pp. 27067-27084
\item
  Huang et al., 2020

  D. Huang, X. Chen, Z. Liu, C. Lyu, S. Wang, X. Chen

  A static bike repositioning model in a hub-and-spoke network framework

  Transp. Res. Part E: Logist. Transp. Rev., 141 (2020), Article 102031
\item
  Jia et al., 2020a

  H. Jia, H. Miao, G. Tian, M. Zhou, Y. Feng, Z. Li, J. Li

  Multiobjective bike repositioning in bike-sharing systems via a modified artificial bee colony
  algorithm

  IEEE Trans. Autom. Sci. Eng., 17 (2) (2020), pp. 909-920
\item
  Jia et al., 2020b

  Y. Jia, W. Zeng, Y. Xing, D. Yang, J. Li

  The bike-sharing rebalancing problem considering multi-energy mixed fleets and traffic
  restrictions

  Sustainability, 13 (1) (2020), p. 270
\item
  Jin et al., 2022

  Y. Jin, C. Ruiz, H. Liao

  A simulation framework for optimizing bike rebalancing and maintenance in large-scale bike-sharing
  systems

  Simul. Model Pract. Theory, 115 (2022), Article 102422
\item
  Kaspi, 2015

  M. Kaspi

  Bike-sharing systems: parking reservation policies and maintenance operations

  Tel Aviv University (2015)

  Ph.D. Thesis
\item
  Kaspi et al., 2016

  M. Kaspi, T. Raviv, M. Tzur

  Detection of unusable bicycles in bike-sharing systems

  Omega, 65 (2016), pp. 10-16
\item
  Kaspi et al., 2017

  M. Kaspi, T. Raviv, M. Tzur

  Bike-sharing systems: User dissatisfaction in the presence of unusable bicycles

  IISE Trans., 49 (2) (2017), pp. 144-158
\item
  Lee et al., 2020

  E. Lee, B. Son, Y. Han

  Optimal relocation strategy for public bike system with selective pick-up and delivery

  Transp. Res. Rec.: J. Transp. Res. Board, 2674 (4) (2020), pp. 325-336
\item
  Li et al., 2016

  Y. Li, W.Y. Szeto, J. Long, C.S. Shui

  A multiple type bike repositioning problem

  Transp. Res. Part B: Methodol., 90 (2016), pp. 263-278
\item
  Liu et al., 2018

  Y. Liu, W.Y. Szeto, S.C. Ho

  A static free-floating bike repositioning problem with multiple heterogeneous vehicles, multiple
  depots, and multiple visits

  Transp. Res. Part C: Emerg. Technol., 92 (2018), pp. 208-242
\item
  Lu et al., 2019

  H. Lu, M. Zhang, S. Su, X. Gao, C. Luo

  Broken bike recycling planning for sharing bikes system

  IEEE Access, 7 (2019), pp. 177354-177361
\item
  Lu et al., 2022

  L. Lu, S. Zhao, Q.-C. He, N. Zhu

  Task assignment in predictive maintenance for free-float bicycle sharing systems

  Comput. Ind. Eng., 169 (2022), Article 108214
\item
  Lu et al., 2020

  Y. Lu, U. Benlic, Q. Wu

  An effective memetic algorithm for the generalized bike-sharing rebalancing problem

  Eng. Appl. Artif. Intell., 95 (2020), Article 103890
\item
  Luo et al., 2020

  H. Luo, F. Zhao, W.-Q. Chen, H. Cai

  Optimizing bike sharing systems from the life cycle greenhouse gas emissions perspective

  Transp. Res. Part c: Emerg. Technol., 117 (2020), Article 102705
\item
  O'Mahony and Shmoys, 2015

  E. O'Mahony, D.B. Shmoys

  Data analysis and optimization for (Citi)bike sharing

  In: Proceedings of the Twenty-Ninth AAAI Conference on Artificial Intelligence (2015), pp. 687-694
\item
  Pan et al., 2019

  L. Pan, Q. Cai, Z. Fang, P. Tang, L. Huang

  A deep reinforcement learning framework for rebalancing dockless bike sharing systems

  In: Proceedings of the Thirty-Third AAAI Conference on Artificial Intelligence (2019), pp.
  1393-1400
\item
  Rainer-Harbach et al., 2015

  M. Rainer-Harbach, P. Papazek, G.R. Raidl, B. Hu, C. Kloimüllner

  PILOT, GRASP, and VNS approaches for the static balancing of bicycle sharing systems

  J. Glob. Optim., 63 (3) (2015), pp. 597-629
\item
  Raviv and Kolka, 2013

  T. Raviv, O. Kolka

  Optimal inventory management of a bike-sharing station

  IIE Trans., 45 (10) (2013), pp. 1077-1093
\item
  Raviv et al., 2013

  T. Raviv, M. Tzur, I.A. Forma

  Static repositioning in a bike-sharing system: models and solution approaches

  Eur. J. Transp. Logist., 2 (3) (2013), pp. 187-229
\item
  Regue and Recker, 2014

  R. Regue, W. Recker

  Proactive vehicle routing with inferred demand to solve the bikesharing rebalancing problem

  Transp. Res. Part E: Logist. Transp. Rev., 72 (2014), pp. 192-209
\item
  Shi et al., 2024

  Z. Shi, M. Xu, Y. Song, Z. Zhu

  Multi-Platform dynamic game and operation of hybrid Bike-Sharing systems based on reinforcement
  learning

  Transp. Res. Part E: Logist. Transp. Rev., 181 (2024), Article 103374
\item
  Shui and Szeto, 2018

  C.S. Shui, W.Y. Szeto

  Dynamic green bike repositioning problem -- a hybrid rolling horizon artificial bee colony
  algorithm approach

  Transp. Res. Part D: Transp. Environ., 60 (2018), pp. 119-136
\item
  Singla et al., 2015

  A. Singla, M. Santoni, G. Bartók, P. Mukerji, M. Meenen, A. Krause

  Incentivizing users for balancing bike sharing systems

  In: Proceedings of the Twenty-Ninth AAAI Conference on Artificial Intelligence (2015), pp. 723-729
\item
  Szeto et al., 2016

  W.Y. Szeto, Y. Liu, S.C. Ho

  Chemical reaction optimization for solving a static bike repositioning problem

  Transp. Res. Part D: Transp. Environ., 47 (2016), pp. 104-135
\item
  Usama et al., 2019

  M. Usama, Y. Shen, O. Zahoor

  A free-floating bike repositioning problem with faulty bikes

  Procedia Comput. Sci., 151 (2019), pp. 155-162
\item
  Usama et al., 2020

  M. Usama, Y. Shen, O. Zahoor, Q. Bao

  Dockless bike-sharing system: Solving the problem of faulty bikes with simultaneous rebalancing
  operation

  J. Transp. Land Use, 13 (1) (2020), pp. 491-515
\item
  Vidal, 2022

  T. Vidal

  Hybrid genetic search for the CVRP: Open-source implementation and SWAP* neighborhood

  Comput. Oper. Res., 140 (2022), Article 105643
\item
  Vidal et al., 2012

  T. Vidal, T.G. Crainic, M. Gendreau, N. Lahrichi, W. Rei

  A hybrid genetic algorithm for multidepot and periodic vehicle routing problems

  Oper. Res., 60 (3) (2012), pp. 611-624
\item
  Vidal et al., 2013

  T. Vidal, T.G. Crainic, M. Gendreau, C. Prins

  A hybrid genetic algorithm with adaptive diversity management for a large class of vehicle routing
  problems with time-windows

  Comput. Oper. Res., 40 (1) (2013), pp. 475-489
\item
  Vidal et al., 2014

  T. Vidal, T.G. Crainic, M. Gendreau, C. Prins

  A unified solution framework for multi-attribute vehicle routing problems

  Eur. J. Oper. Res., 234 (3) (2014), pp. 658-673
\item
  Wang and Wang, 2021

  J. Wang, Y. Wang

  A two-stage incentive mechanism for rebalancing free-floating bike sharing systems: Considering
  user preference

  Transp. Res. Part F: Traffic Psychol. Behav., 82 (2021), pp. 54-69
\item
  Wang and Szeto, 2018

  Y. Wang, W.Y. Szeto

  Static green repositioning in bike sharing systems with broken bikes

  Transp. Res. Part D: Transp. Environ., 65 (2018), pp. 438-457
\item
  Wang and Szeto, 2021

  Y. Wang, W.Y. Szeto

  An enhanced artificial bee colony algorithm for the green bike repositioning problem with broken
  bikes

  Transp. Res. Part C: Emerg. Technol., 125 (2021), Article 102895
\item
  Wu, 2021

  Wu, C., 2021. Determination of tort liability in traffic accidents based on performance parameter
  standard of electric bicycle. In: 2021 2nd international conference on computers, information
  processing and advanced education. New York, USA: Association for Computing Machinery (CIPAE
  2021), pp. 1485--1489.
\item
  Xu et al., 2019

  G. Xu, Y. Li, T. Xiang, D. Zhao

  A bike repositioning problem with broken bikes

  Syst. Eng., 37 (02) (2019), pp. 91-99
\item
  You, 2019

  P.-S. You

  A two-phase heuristic approach to the bike repositioning problem

  Appl. Math. Model, 73 (2019), pp. 651-667
\item
  Zhang et al., 2020

  D. Zhang, W. Xu, B. Ji, S. Li, Y. Liu

  An adaptive tabu search algorithm embedded with iterated local search and route elimination for
  the bike repositioning and recycling problem

  Comput. Oper. Res., 123 (2020), Article 105035
\item
  Zhang et al., 2017

  D. Zhang, C. Yu, J. Desai, H.Y.K. Lau, S. Srivathsan

  A time-space network flow approach to dynamic repositioning in bicycle sharing systems

  Transp. Res. Part B: Methodol., 103 (2017), pp. 188-207
\item
  Zhang et al., 2018

  S. Zhang, G. Xiang, Z. Huang

  Bike-sharing static rebalancing by considering the collection of bicycles in need of repair

  J. Adv. Transp. (2018), 10.1155/2018/8086378
\end{enumerate}

\begin{description}
\item[\textsuperscript{a}]
\protect\phantomsection\label{np125}
We used four threads and set the parameter \emph{Heuristics} in Gurobi to 0. To avoid no feasible
solution found by Gurobi within the time limit for large networks, for all the experiments run with
Gurobi in Section 5, we used Gurobi's ``\emph{MIP starts}'' feature that enables users to provide it
with initial solutions. We formed the solution by setting trucks to travel to no station and stay at
the depot, and each repairer visits only one station (Station 1 if \emph{\textbar R}\textbar=1,
stations 1 and 2 if \textbar{}\emph{R}\textbar=2).
\end{description}

\begin{description}
\item[\textsuperscript{b}]
\protect\phantomsection\label{np130}
https://www.zhipin.com/job\_detail/81bcad2cd24453bf1HF73NW4FlJU.html.
\end{description}

\end{otherlanguage}

\end{document}
